\chapter{Stability functionals and filtrations}
\index{K-stability!valuative}\index{stability functional}\index{filtration!stability}
\label{chap:bj-stability}

K-stability begins with a degeneration and ends with an inequality.  The
objects between these two ends are varied: weights, filtrations,
non-Archimedean metrics, Monge--Amp\`ere measures, and valuations.  This
chapter makes the passage between them without identifying objects that
only agree after Hausdorff quotient or completion.

The results come from three layers.  The intersection-theoretic functionals
on test configurations are those of \cite{BHJ17}.  The metric geometry of
all norms and filtrations is from \cite{BJKS1}.  Divisorial stability and
stability for arbitrary filtrations are from \cite{BJKS2}.  These later
results are not part of the foundational BJ22 paper, and
\Cref{sec:ch50-dated} records for each of them a date, a source, and the
exact hypotheses under which it is used here.

Unless otherwise stated, \(X\) is an integral projective variety of
dimension \(n\), \(L\) is an ample \(\mathbf Q\)-line bundle, and
\[
             R=\bigoplus_{m\geq0}R_m,\qquad
             R_m=H^0(X,mL),\qquad V=(L^n).
\tag{50.1}\label{eq:stability-section-ring}
\]
Degrees are always taken sufficiently divisible.  Normality and
\(\operatorname{char}k=0\) will be imposed when log discrepancies enter.

\section{Orientation}
\label{sec:ch50-orientation}

\subsection{What this chapter does}
\label{sec:ch50-what}

\Cref{chap:bj-test-configurations} produced, from a one-parameter
degeneration of \((X,L)\), three equivalent packets of data: a test
configuration, a flag ideal, and a linearly bounded multiplicative
filtration of the section ring.  \Cref{chap:bj-fs-pl-psh} attached to the
third of these a function on \(X^{\An}\), the Fubini--Study function
\eqref{eq:norm-fs-function}, and identified the finitely generated case
with the piecewise linear metrics.
\Crefrange{chap:bj-energy}{chap:bj-linear-growth} then built an analysis
around those functions: an energy \(E_\omega\), a finite-energy class
\(\cE^1(\omega)\), a dual energy \(E^\vee_\omega\) on measures, two strong
topologies exchanged by the Monge--Amp\`ere operator, and an exact
description of the valuations that the theory sees.

\emph{This chapter is the payoff of \Cref{chap:bj-test-configurations}.}
It closes the circle by running the construction backwards: it takes the
analytic apparatus and converts it into numerical invariants of
degenerations --- the Donaldson--Futaki invariant, the Mabuchi and Ding
energies, the minimum norm, the entropy --- and then into stability
conditions.  Four things are done.

\emph{First, the correspondence between norms and metrics is proved as a
statement about metric spaces, not only about sets.}  The set-level
dictionary of \Cref{chap:bj-test-configurations} between filtrations and
Fubini--Study functions is not injective: distinct norms can have the same
Fubini--Study function, and the failure is exactly homogenization.  We prove
in \Cref{sec:ch50-norms} that \(\FS\) and the sup-norm operator form an
adjoint pair \eqref{eq:ch50-galois}, that the resulting closure operator
computes the homogenization up to \(d_1\)-distance zero, and that the
spectral pseudometrics \(d_1\) and \(d_\infty\) of
\eqref{eq:stability-finite-dimensional-dp} are genuine pseudometrics with
explicit descriptions --- \(d_\infty\) as a uniform distance
(\Cref{prop:ch50-dinfty-sup}) and \(d_1\) as an integral of dimensions of
intersections of filtration steps (\Cref{thm:ch50-d1-integral}).  The
simultaneous-diagonalization lemma that makes the relative spectrum
well-defined, \Cref{thm:ch50-joint-orthogonal}, is proved from scratch:
over a trivially valued field it is the elementary statement that two flags
in a finite-dimensional vector space admit a common adapted basis.

\emph{Second, the energy functionals are computed as moments of a
Duistermaat--Heckman measure, and the comparisons between them are proved
with sharp constants.}  \Cref{sec:ch50-dh} identifies \(E\), \(\sup\), \(J\),
\(I\) and the minimum norm \(\|\cdot\|_{\min}\) with explicit functionals of
\(\DHm_\chi\), and proves the two-sided comparison
\eqref{eq:stability-J-L1-comparison} between \(J\) and the centred
\(L^1\)-norm with the sharp constant \(c_n=2n^n/(n+1)^{n+1}\).  The proof
given in \Cref{thm:ch50-J-L1} uses nothing but the Brunn--Minkowski
concavity of \Cref{cor:ch33-dh-structure}\,(2) and one calculus estimate;
it is short, it is sharp, and it exhibits the extremal measure.

\emph{Third, entropy is treated as a genuinely extended-real functional.}
The log discrepancy \(A_{X,B}\) is \([0,+\infty]\)-valued and lower
semicontinuous, so \(\Ent_{X,B}\) is weakly lower semicontinuous
(\Cref{thm:ch50-entropy-lsc}); but it is \emph{not} finite on
\(\cM^1(\omega)\), and \Cref{ex:ch50-infinite-entropy} produces an explicit
measure on \((\mathbf P^2)^{\An}\) with dual energy at most \(2/3\) and
infinite entropy.  Every later statement that needs finite entropy says so.

\emph{Fourth, the stability conditions are separated by date.}  Uniform
K-stability in the sense of \eqref{eq:stability-uniform-K} tests only
finitely generated filtrations; divisorial stability in the sense of
\eqref{eq:stability-divisorial-threshold} tests divisorial measures;
filtration stability tests everything.  These are three different
conditions, they became available at three different times, and
\Cref{rem:stability-gap-general-polarization} records precisely what is and
is not known about their coincidence.  In the log Fano case a further
reduction is available, and it is imported from four separate papers, each
dated.

\subsection{The six tasks}
\label{sec:ch50-tasks}

\begin{enumerate}
\item\label{itm:ch50-task:norms}
      \emph{Norms, filtrations, and finite-energy metrics.}  Prove the
      joint-orthogonality lemma, construct the relative spectrum, prove
      that \(d_1\) and \(d_\infty\) are pseudometrics with the stated
      closed forms, prove the adjunction between \(\FS\) and the sup-norm
      operator, and state the Hausdorff-quotient and completion theorems of
      \cite{BJKS1} with their exact scope.
      \Cref{sec:ch50-norms}.
\item\label{itm:ch50-task:energy}
      \emph{Non-Archimedean energy functionals.}  Identify \(E\), \(J\),
      \(I\) and \(\|\cdot\|_{\min}\) as functionals of the
      Duistermaat--Heckman measure and of the Monge--Amp\`ere measure,
      prove \eqref{eq:stability-minimum-IJ} and the sharp comparison
      \eqref{eq:stability-J-L1-comparison}, and compute all of them on
      \(\mathbf P^1\).  \Cref{sec:ch50-dh}.
\item\label{itm:ch50-task:entropy}
      \emph{Log discrepancy and entropy.}  Define \(A_{X,B}\) on
      \(X^{\An}\), prove lower semicontinuity and hence weak lower
      semicontinuity of \(\Ent_{X,B}\), prove the model formula
      \eqref{eq:stability-model-entropy}, and settle the finiteness
      question by an explicit counterexample.
      \Cref{sec:ch50-entropy}.
\item\label{itm:ch50-task:mabuchi}
      \emph{The Mabuchi functional.}  Prove the Chen--Tian decomposition
      \eqref{eq:stability-chen-tian}, the comparison with the
      Donaldson--Futaki invariant \eqref{eq:stability-M-DF-correction}, and
      the identification \eqref{eq:stability-M-beta} of \(M^{\NA}\) with a
      functional of the Monge--Amp\`ere measure alone.
      \Cref{sec:ch50-mabuchi}.
\item\label{itm:ch50-task:ding}
      \emph{The Ding functional.}  Prove the elementary calculus of
      \(L^{\NA}\), prove \(D^{\NA}\le M^{\NA}\) in full, and isolate
      exactly which part of the equality case and of the uniform
      equivalence \eqref{eq:stability-ding-M-uniform} is imported.
      \Cref{sec:ch50-ding}.
\item\label{itm:ch50-task:uniform}
      \emph{Uniform K-stability and valuative criteria.}  State the three
      stability conditions, prove the elementary equivalences among the
      valuative invariants \(A\), \(S\), \(T\), \(\beta\), \(j\) and
      \(\delta\), and import the K-stability equivalences with their dates.
      \Cref{sec:ch50-uniform,sec:ch50-valuative}.
\end{enumerate}

Four worked examples carried to explicit numbers occupy
\Cref{sec:ch50-examples}: the functionals of an explicit test configuration
on \(\mathbf P^1\) (\Cref{ex:ch50-p1-dnc}), a toric pair whose stability is
decided from the polytope (\Cref{ex:ch50-toric-F1}), a filtration that is
not finitely generated (\Cref{ex:ch50-harmonic}), and a polarized variety
that is K-semistable but not uniformly K-stable
(\Cref{ex:ch50-semistable-not-uniform}).

\subsection{Standing hypotheses and notation}
\label{sec:ch50-hypotheses}

The ground field \(k\) is algebraically closed and trivially valued.  The
variety \(X\) is integral, projective over \(k\), of dimension \(n\ge1\),
and \emph{unibranch}, so that the envelope theorem of
\Crefrange{chap:bj-envelope-statement}{chap:bj-envelope-consequences} is
available.  We write \(L\) for an ample \(\mathbf Q\)-line bundle,
\(\omega:=c_1(L)\in\Amp(X)\), and \(V:=(\omega^n)=(L^n)>0\) as in
\eqref{eq:bj-volume}.  The section ring is \eqref{eq:stability-section-ring}
and \(N_m:=\dim_kR_m\).

The classes \(\PSH(\omega)\), \(\cE^1(\omega)\), \(\cE^1_0(\omega)\),
\(\cM(\omega)\) and \(\cM^1(\omega)\) are those of
\Cref{def:bj-psh,def:bj-e1,def:bj-finite-energy-measure} and
\eqref{eq:ch47-normalized-classes}; the functionals \(E_\omega\),
\(E^\vee_\omega\), \(I\), \(J_\psi\), \(J\) are those of
\eqref{eq:bj-energy-definition}, \eqref{eq:bj-dual-energy},
\eqref{eq:bj-I-J-potentials} and \eqref{eq:ch47-J-functional}.  The loci
\[
 X^{\divis}\subset X^{\lin}\subset X^{\val}\subset X^{\An}
\]
are those of \eqref{eq:ch34-loci}, and \(T_\omega=T_L\) is the exhaustion of
\Cref{def:bj-linear-growth} and \Cref{prop:ch49-T-equals-T}.  When
log discrepancies enter --- from \Cref{sec:ch50-entropy} onwards --- we
assume in addition that \(X\) is normal, that
\(\operatorname{char}k=0\), and that \(B\) is a \(\mathbf Q\)-Weil divisor
on \(X\) with \(K_{X,B}:=K_X+B\) \(\mathbf Q\)-Cartier.  \emph{We do not
assume \(B\) effective} except where the log Fano hypothesis is imposed
explicitly.

\subsection{The normalization, inherited}
\label{sec:ch50-normalization}

\begin{warning}{Probability normalization is in force}
\label{warn:ch50-normalization}
\Cref{warn:ch46-normalization} declared, for the whole of
\Crefrange{chap:bj-energy}{chap:bj-stability}, a departure from the
front-matter convention of \Cref{app:normalizations}: in these chapters
\(\MA_\omega(\phi)\) is a \emph{probability} measure, \(\MA^{\mathrm u}_\omega
=V\MA_\omega\) is the unnormalized one, \(E_\omega(0)=0\) and
\(E_\omega(\phi+c)=E_\omega(\phi)+c\).  That convention is in force here
without amendment, and it is the reason every functional of this chapter is
an average rather than a total: \(E\) is the first moment of a
\emph{probability} Duistermaat--Heckman measure
(\Cref{prop:ch46-energy-dh}), the entropy \eqref{eq:stability-entropy}
integrates against a probability measure, and the Donaldson--Futaki
invariant carries the factor \(V^{-1}\).

Three consequences are used constantly.  First, \(\MA_\omega(0)
=\delta_{v_{\triv}}\) and \(\sup_{X^{\An}}\phi=\phi(v_{\triv})\), by
\Cref{lem:ch47-three-normalizations}; so \(J(\phi)=\sup\phi-E(\phi)\) is
simultaneously the normalized energy and the functional \(J_0\).  Second,
dual energy is finite only on probability measures
(\Cref{lem:ch47-mass-forced}); so a Duistermaat--Heckman measure and a
Monge--Amp\`ere measure are objects of the same kind, both of mass one, and
no factor \(V\) intervenes when they are compared.  Third, the
Duistermaat--Heckman conventions of \Cref{warn:ch33-dh-normalizations} must
be consulted before any weight asymptotics is quoted: the unnormalized
measure \eqref{eq:ch33-dh-unnormalized} has mass \(V\), and the toric
normalization of \Cref{chap:model-calculations} differs from both by
\(n!\).
\end{warning}

\subsection{The erratum firewall}
\label{sec:ch50-firewall}

This chapter is about exactly the functionals \(E\), \(I\), \(J\), \(H\),
\(M\), \(R\), \(D\) whose \emph{complex} asymptotic slopes along rays were
the subject of \cite{BHJ19}, part of which was retracted.  The firewall is
therefore not a formality here, and it is stated before anything else is
proved.

\begin{warning}{Binding: no erratum-affected input, direct or transitive}
\label{warn:ch50-erratum-firewall}
\index{erratum!Boucksom--Hisamoto--Jonsson}
No statement of this chapter uses, directly or transitively, any of
\[
 \begin{gathered}
 \text{\cite[Theorem~5.6]{BHJ19}},\qquad
 \text{\cite[Theorem~5.4]{BHJ19}},\\
 \text{\cite[Theorem~C]{BHJ19}},\qquad
 \text{\cite[Corollaries~D and~E]{BHJ19}},
 \end{gathered}
\]
that is, any of the results quarantined in
\eqref{eq:rf-quarantined-results}; nor any reformulation of one of them
under another name.  The prohibition of \Cref{warn:rf-no-surrogate} is in
force, and so is the standing constraint adopted in
\Cref{warn:ch33-erratum-firewall}.  In particular:
\begin{enumerate}
\item\label{itm:ch50-fw:noslope}
      \emph{No asymptotic-slope statement is asserted or used.}  Nowhere in
      this chapter is a limit \(\lim_{t\to\infty}t^{-1}F(\varphi_t)\) of a
      complex energy functional along a ray identified with a
      non-Archimedean functional.  Every functional here is defined
      directly on the non-Archimedean side, by intersection theory
      (\Cref{chap:bj-energy}), by integration against a Monge--Amp\`ere
      measure, or by an infimum over valuations.
\item\label{itm:ch50-fw:nocoercivity}
      \emph{No implication from uniform K-stability to coercivity is
      used.}  In particular \eqref{eq:stability-ding-M-uniform} is an
      equivalence \emph{between two non-Archimedean functionals on the
      space of positive model metrics}, imported from
      \cite[Theorem~8.6]{BHJ17}; it is not the K\"ahler-geometric
      coercivity statement of \cite[Theorem~C]{BHJ19}, and it must not be
      read as one.
\item\label{itm:ch50-fw:differentnames}
      \emph{Two ``Theorem~C''s must not be confused.}  We cite
      \cite[Theorem~C]{BJKS1}, on the Monge--Amp\`ere bijection between
      divisorial norms and divisorial measures.  That is a theorem of a
      different paper, by different authors in a different year, and it is
      unaffected by the erratum.
\item\label{itm:ch50-fw:bhj17ok}
      \emph{\cite{BHJ17} is not affected.}  The 2017 paper on uniform
      K-stability, Duistermaat--Heckman measures and singularities of pairs
      is an algebraic paper over an arbitrary field; the erratum concerns
      the later K\"ahler-geometric asymptotics of \cite{BHJ19}.  All the
      intersection-theoretic identities of
      \Cref{sec:ch50-mabuchi,sec:ch50-ding} come from \cite{BHJ17} and are
      unaffected.
\end{enumerate}
The complete dependency ledger is \Cref{thm:rf-erratum-audit}; asymptotic
slope statements appear in this book only in
\Cref{chap:rays-filtrations}, each with its own post-erratum source.
\end{warning}

\subsection{The dating discipline}
\label{sec:ch50-dating}

\begin{important}{Every imported statement carries a date}
\label{imp:ch50-dating}
K-stability moved quickly after the foundational trivially valued
pluripotential theory of \cite{BJ22} was written, and a chapter on
stability functionals is exposed to the temptation of absorbing later
theorems into a single narrative.  We do not do this.  The rule adopted
here, and applied without exception, is:
\begin{enumerate}
\item\label{itm:ch50-date:proved}
      a statement proved in this book is announced as such, and its proof
      is given;
\item\label{itm:ch50-date:imported}
      a statement imported from the literature is announced as such, with
      the paper, the numbered statement inside that paper, the year of
      first public availability and the year of publication, and with its
      hypotheses displayed in full rather than absorbed into the standing
      hypotheses of \Cref{sec:ch50-hypotheses};
\item\label{itm:ch50-date:foundational}
      no result later than \cite{BJ22} is presented as part of the
      foundational theory, and no such result is used silently in the proof
      of a statement that is presented as foundational.
\end{enumerate}
\Cref{sec:ch50-dated} is the resulting table.  It separates the layers:
\cite{BHJ17} (2015 preprint, 2017 journal), \cite{BJ22} together with its
addendum \cite{BJ22b} (2018 preprint, 2022 and 2024 journal),
\cite{BJKS1} (2021 preprint, 2025 journal), \cite{BJKS2} (2022 preprint,
2023 journal), and the Fano-specific line \cite{Li17,Fuj19,BJ20,LXZ22}.
\end{important}

\subsection{Imported statements}
\label{sec:ch50-imported}

\begin{important}{The seven imported statements}
\label{imp:ch50-imported}
The following are used without proof.  Everything else in this chapter is
proved here or in an earlier chapter of this book.
\begin{enumerate}
\item\label{itm:ch50-import:dp}
      \emph{The \(d_p\) pseudometrics for \(1<p<\infty\).}  That
      \eqref{eq:stability-finite-dimensional-dp} satisfies the triangle
      inequality for \(1<p<\infty\) is proved in \cite{BJKS1}.  We prove the
      cases \(p=1\) (\Cref{thm:ch50-d1-integral}) and \(p=\infty\)
      (\Cref{prop:ch50-dinfty-sup}) and the interpolation
      \(d_1\le d_p\le d_\infty\)
      (\Cref{prop:ch50-dp-interpolation}); nothing later uses the
      intermediate cases.
\item\label{itm:ch50-import:quotient}
      \emph{The Fubini--Study quotient and the completion theorem},
      \Cref{thm:stability-fs-quotient} and
      \Cref{thm:stability-norm-completion}: \cite[Theorems~A and~B]{BJKS1}.
\item\label{itm:ch50-import:divisorial}
      \emph{The divisorial Monge--Amp\`ere bijection}
      \eqref{eq:stability-divnorm-divmeasure}: \cite[Theorem~C]{BJKS1}.
\item\label{itm:ch50-import:S}
      \emph{The identity \(S_L(v)=E^\vee_\omega(\delta_v)\)}:
      \cite[Theorem~7.22(iii)]{BJKS1}.  We verify it independently on
      \(\mathbf P^1\) and on the toric slice.
\item\label{itm:ch50-import:beta}
      \emph{Differentiability of \(\omega\mapsto\|\mu\|_\omega\) on the
      ample cone, and the strong lower semicontinuous extension of
      \(\beta\)}: \cite[Theorem~2.15 and \S\S3.4, 4.1]{BJKS2}.
\item\label{itm:ch50-import:filtration}
      \emph{The identity of divisorial and filtration thresholds},
      \Cref{thm:stability-divisorial-filtrations}:
      \cite[Theorem~5.8]{BJKS2}.
\item\label{itm:ch50-import:fano}
      \emph{The Fano valuative criteria}, \Cref{thm:stability-valuative-fano}
      and the coercivity half of \eqref{eq:stability-ding-M-uniform}:
      \cite{Li17,Fuj19,BJ20} and \cite[Theorem~8.6]{BHJ17}.  These use the
      minimal model program and are not consequences of pluripotential
      theory.
\end{enumerate}
\end{important}

\begin{important}{The envelope theorem is used without further comment}
\label{imp:ch50-envelope}
As in \Cref{imp:ch46-envelope-imported,imp:ch47-envelope,imp:ch48-envelope}
and \Cref{imp:ch49-envelope}, the variety \(X\) is unibranch and
\Cref{thm:bj-final-envelope} is available; the firewall
\Cref{imp:ch45-firewall-discharged} has been discharged and \(P_\omega\) is
an operator (\Cref{imp:ch45-obstacle-operator}), with the calculus of
\Cref{prop:bj-final-obstacle-calculus} and
\Cref{prop:ch45-obstacle-calculus-2}.  Four consequences are used.
\begin{enumerate}
\item\label{itm:ch50-env:cy}
      The Calabi--Yau theorem \Cref{cor:bj-cy-unibranch}: every
      \(\mu\in\cM^1(\omega)\) is \(\MA_\omega(\phi)\) for a
      \(\phi\in\cE^1(\omega)\), unique up to a constant.  This is what
      makes \eqref{eq:stability-M-all-norms} a definition rather than a
      formal expression.
\item\label{itm:ch50-env:homeo}
      \(\MA_\omega:\cE^1(\omega)/\mathbf R\to\cM^1(\omega)\) is a
      homeomorphism for the strong topologies
      (\Cref{thm:ch48-ma-homeomorphism}); this is the analytic half of
      \Cref{thm:stability-norm-completion}.
\item\label{itm:ch50-env:uniqueness}
      Uniqueness of solutions, \Cref{thm:ch47-uniqueness}, which is used
      whenever a metric is identified from its Monge--Amp\`ere measure ---
      in particular in every example of \Cref{sec:ch50-examples}.
\item\label{itm:ch50-env:completeness}
      Completeness of \(\cE^1(\omega)\) for the strong topology
      (\Cref{thm:ch48-completeness}), without which the second clause of
      \Cref{thm:stability-norm-completion} would be vacuous.
\end{enumerate}
\end{important}

\subsection{Relation to \texorpdfstring{\Cref{part:bfj}}{Part III}}
\label{sec:ch50-part3}

There is no counterpart of this chapter over a nontrivially valued base,
and the reason is the same as in \Cref{sec:ch49-part3}: the objects being
metrized here are degenerations of \((X,L)\) over the affine line, and the
trivially valued analytification of \(X\) is the space in which such a
degeneration becomes a \emph{point-like} datum, namely a function on the
cone \(X^{\An}\).  Over a discretely valued field \(K\) the analogue of a
test configuration is a model, and the analogue of the present chapter is
the arithmetic or geometric theory of heights: the functionals \(E\) and
\(E^\vee\) of \Cref{chap:bfj-energy,chap:bfj-finite-energy-measures} exist there,
but there is no scaling action, no exhaustion function, no
Duistermaat--Heckman measure, and no entropy attached to a log discrepancy
normalized against a distinguished point.

What \emph{is} shared is a formal pattern: vertical divisors are evaluated
against relative intersection classes, projection formulae control change
of model, and model Monge--Amp\`ere measures are atomic at components of a
special fibre.  The ambient pairings are nevertheless different.  Over a
discretely valued field, \Cref{chap:bfj-energy} only uses the legal pairing
with an actual vertical divisor in one slot and \(n\) relative closed-model
classes in the other slots; it does not define an absolute
\((n+1)\)-fold power of closed model forms.  In the present trivially
valued setting, the pairing of \Cref{def:bj-pl-energy-pairing} is a genuine
absolute intersection number on a test configuration compactified over
\(\mathbf P^1\).  The identities in \Cref{sec:ch50-mabuchi} belong to this
latter framework.  They cannot be asserted over the discretely valued
base merely by reinterpreting the same symbols; a separate comparison
theorem would be required.

\subsection{Roadmap}
\label{sec:ch50-roadmap}

\Cref{sec:ch50-norms} treats norms, filtrations and the metric geometry of
the space they form.  \Cref{sec:ch50-dh} computes the energy functionals as
moments.  \Cref{sec:ch50-entropy} constructs the entropy and settles its
finiteness and semicontinuity.  \Cref{sec:ch50-mabuchi} and
\Cref{sec:ch50-ding} are the Mabuchi and Ding functionals.
\Cref{sec:ch50-uniform} defines the three stability conditions and
\Cref{sec:ch50-valuative} gives the valuative criteria in the log Fano
case.  \Cref{sec:ch50-examples} contains the four worked examples,
\Cref{sec:ch50-scope} audits the hypotheses,
\Cref{sec:ch50-dictionary} is the chapter dictionary,
\Cref{sec:ch50-dated} the dated scope table, and
\Cref{sec:ch50-source} the provenance and historical note.

\section{Norms, filtrations, and finite-energy metrics}
\label{sec:ch50-norms}

\subsection{Graded norms and filtrations}
\label{sec:ch50-norm-recall}

Recall from \Cref{chap:bj-test-configurations} that an additive norm
\(\chi\) on \(R\) is a collection of maps
\[
 \chi_m:R_m\longrightarrow\mathbf R\cup\{+\infty\},\qquad
 \chi_m(0)=+\infty,
\]
that is non-Archimedean on each \(R_m\), invariant under
\(k^\times\), superadditive under multiplication, and linearly bounded:
\[
 \begin{split}
 \chi_m(s+s')&\geq\min\{\chi_m(s),\chi_m(s')\},\\
 \chi_{m+m'}(ss')&\geq\chi_m(s)+\chi_{m'}(s'),\\
 |\chi_m(s)|&\leq Cm\quad(s\ne0).
 \end{split}
\tag{50.2}\label{eq:stability-norm-axioms}
\]
It is equivalent to the decreasing filtration
\[
             \mathcal F^\lambda R_m
             =\{s\in R_m\mid\chi_m(s)\geq\lambda\}.
\tag{50.3}\label{eq:stability-norm-filtration}
\]
Write \(\mathcal N_{\mathbf R}\) for the space of such norms, after the
usual identification on Veronese subalgebras, and
\(\mathcal T_{\mathbf R}\subset\mathcal N_{\mathbf R}\) for the finite-type
norms.  The integral finite-type locus
\(\mathcal T_{\mathbf Z}\) is identified by the Rees construction with
ample test configurations.

The dictionary between the two descriptions is
\Cref{lem:ch33-filtration-norm-bijection}; the linear-boundedness
conditions correspond under \eqref{eq:linear-bounds-filtration}; and the
Fekete limits
\[
 \lambda_{\max}(\chi):=\lim_{m\to\infty}\frac1m\max_{0\ne s\in R_m}\chi_m(s)
 =\sup_{m\ge1}\frac1m\max_{0\ne s\in R_m}\chi_m(s)
\tag{50.4}\label{eq:ch50-lambda-max}
\]
exist by \eqref{eq:ch33-fekete-limits}.  Throughout this section we work
with \(\Lambda\subset\mathbf R\) a subgroup, and write
\(\mathcal N_\Lambda\), \(\mathcal T_\Lambda\) for the norms whose jumping
values lie in \(\Lambda\); the two cases of interest are
\(\Lambda=\mathbf Z\), where \(\mathcal T_{\mathbf Z}\) consists of the
norms of ample test configurations, and \(\Lambda=\mathbf R\).

\begin{remark}[Why a filtration is more than a degeneration]
\label{rem:ch50-why-filtrations}
Three reasons run through the chapter.  The first is that
\(\mathcal T_{\mathbf R}\) is not complete for any of the metrics below
(\Cref{thm:stability-norm-completion}), so a variational problem posed on
test configurations has no reason to have a solution there.  The second is
that the natural operations --- homogenization, taking a minimum of two
norms, restricting to a subvariety, base change --- leave the finite-type
locus.  The third is that the analytic side is already complete: the class
\(\cE^1(\omega)\) of \Cref{def:bj-e1} and the space \(\cM^1(\omega)\) of
\Cref{def:bj-finite-energy-measure} were constructed in
\Crefrange{chap:bj-energy}{chap:bj-strong-topologies} without any
finiteness hypothesis, and \Cref{ex:ch33-harmonic-not-fg} exhibits a
perfectly good element of \(\PSH(\omega)\) arising from a filtration that
is not of finite type.  Nothing forces us to work in the smaller category,
and \Cref{sec:ch50-uniform} shows that the stability conditions genuinely
differ according to which category is used for testing.
\end{remark}

\subsection{Diagonalization of a single norm}
\label{sec:ch50-diagonal}

Over a trivially valued field the finite-dimensional theory of
non-Archimedean norms is linear algebra: a norm is a filtration by
subspaces, and orthogonality is adaptedness.  We record this once.

Let \(W\) be a finite-dimensional \(k\)-vector space.  By a
\emph{norm} on \(W\) we mean a map \(\chi:W\to\mathbf R\cup\{+\infty\}\)
with \(\chi(0)=+\infty\), \(\chi(w)<+\infty\) for \(w\ne0\),
\(\chi(\lambda w)=\chi(w)\) for \(\lambda\in k^\times\), and
\(\chi(w+w')\ge\min\{\chi(w),\chi(w')\}\).  Its filtration is
\(F^\lambda_\chi W:=\{\chi\ge\lambda\}\).

\begin{lemma}
\label{lem:ch50-norm-is-filtration}
Let \(\chi\) be a norm on \(W\), \(\dim W=N\).  Then:
\begin{enumerate}
\item\label{itm:ch50-nf:subspace}
      each \(F^\lambda_\chi W\) is a \(k\)-subspace, the family is
      decreasing and left continuous in \(\lambda\), equals \(W\) for
      \(\lambda\) small and \(0\) for \(\lambda\) large, and has at most
      \(N\) distinct members;
\item\label{itm:ch50-nf:jumps}
      consequently there are real numbers
      \(\lambda_1\ge\lambda_2\ge\cdots\ge\lambda_N\), the \emph{jumping
      values} of \(\chi\), determined by
      \(\dim F^\lambda_\chi W=\#\{j:\lambda_j\ge\lambda\}\);
\item\label{itm:ch50-nf:adapted}
      a basis \(e_1,\dots,e_N\) of \(W\) is \emph{adapted} to \(\chi\) if
      \(F^\lambda_\chi W=\operatorname{span}\{e_j:\chi(e_j)\ge\lambda\}\)
      for every \(\lambda\); adapted bases exist, and for such a basis
      \(\{\chi(e_j)\}=\{\lambda_j\}\) as multisets and
      \[
       \chi\Bigl(\sum_ja_je_j\Bigr)=\min_{a_j\ne0}\chi(e_j)
       \qquad(a_j\in k).
      \tag{50.5}\label{eq:ch50-orthogonal}
      \]
\end{enumerate}
Conversely, every decreasing, left continuous, exhaustive and separated
filtration of \(W\) by subspaces arises from a unique norm.
\end{lemma}

\begin{proof}
\cref{itm:ch50-nf:subspace}  If \(\chi(w),\chi(w')\ge\lambda\) then
\(\chi(w+w')\ge\lambda\) and \(\chi(aw)=\chi(w)\ge\lambda\) for
\(a\in k^\times\), the absolute value on \(k\) being trivial; and
\(\chi(0)=+\infty\).  So \(F^\lambda_\chi W\) is a subspace.  It decreases
in \(\lambda\) by definition, and
\(F^\lambda=\bigcap_{\mu<\lambda}F^\mu\) because the defining condition is
a non-strict inequality, which is left continuity.  Since \(\chi(w)\) is
finite for \(w\ne0\) we have \(\bigcap_\lambda F^\lambda=0\), and since
\(W\) is finite dimensional, \(\chi\) is bounded below on any basis and
hence, by ultrametricity, on \(W\setminus0\); so \(F^\lambda=W\) for
\(\lambda\) small.  A strictly decreasing chain of subspaces of \(W\) has
at most \(N+1\) members, so there are at most \(N\) jumps.

\cref{itm:ch50-nf:jumps}  List the jumps with multiplicity given by the drop
in dimension.

\cref{itm:ch50-nf:adapted}  Existence: refine a basis of the smallest
nonzero step to a basis of each successive step.  Let \(e_1,\dots,e_N\) be
adapted, ordered so that \(\chi(e_1)\ge\cdots\ge\chi(e_N)\), and let
\(w=\sum_ja_je_j\ne0\).  Ultrametricity gives
\(\chi(w)\ge\mu:=\min_{a_j\ne0}\chi(e_j)\).  Let \(j_0\) be the largest
index with \(a_{j_0}\ne0\), so \(\chi(e_{j_0})=\mu\).  If
\(\chi(w)>\mu\), choose \(\mu'\) with \(\mu<\mu'\le\chi(w)\); then
\(w\in F^{\mu'}_\chi W=\operatorname{span}\{e_j:\chi(e_j)\ge\mu'\}\), and
every such \(e_j\) has \(\chi(e_j)>\mu=\chi(e_{j_0})\), hence \(j<j_0\) by
the ordering.  This contradicts \(a_{j_0}\ne0\).  Therefore
\(\chi(w)=\mu\), which is \eqref{eq:ch50-orthogonal}.  Comparing dimensions
of the steps shows that the multiset \(\{\chi(e_j)\}\) is
\(\{\lambda_j\}\).

The converse is immediate: given the filtration, define
\(\chi(w):=\sup\{\lambda:w\in F^\lambda\}\), a supremum attained by left
continuity.
\end{proof}

\begin{remark}[What is special about the trivial valuation]
\label{rem:ch50-trivial-special}
Over a nontrivially valued field a non-Archimedean norm on a
finite-dimensional space is \emph{not} determined by a filtration by
subspaces, and diagonalizability is a theorem requiring either
discreteness, or maximal completeness, or an \(\varepsilon\)-loss.  Here
there is no loss, no hypothesis, and no \(\varepsilon\): the statement is
that a decreasing family of subspaces has an adapted basis.  Every metric
statement of this section rests on that, and this is the structural reason
the trivially valued theory of norms is so much more rigid than its
nontrivially valued counterpart.
\end{remark}

\subsection{Joint orthogonality and the relative spectrum}
\label{sec:ch50-joint}

The definition \eqref{eq:stability-finite-dimensional-dp} of the spectral
distances presupposes a basis orthogonal for two norms at once.  In this
setting that is a purely linear-algebraic fact about pairs of flags, and we
prove it.

\begin{theorem}
\label{thm:ch50-joint-orthogonal}
\index{orthogonal basis!joint}
Let \(W\) be a \(k\)-vector space of finite dimension \(N\), and let
\(\chi,\chi'\) be norms on \(W\).  For \((\lambda,\mu)\in\mathbf R^2\) put
\[
 W^{\lambda,\mu}:=F^\lambda_\chi W\cap F^\mu_{\chi'}W,
 \qquad
 g(\lambda,\mu):=\dim\frac{W^{\lambda,\mu}}
   {W^{\lambda^+,\mu}+W^{\lambda,\mu^+}},
\tag{50.6}\label{eq:ch50-bigraded}
\]
where \(W^{\lambda^+,\mu}:=\bigcup_{\lambda'>\lambda}W^{\lambda',\mu}\) and
similarly for the other entry.  Then:
\begin{enumerate}
\item\label{itm:ch50-jo:count}
      \(g(\lambda,\mu)\ge0\), it vanishes outside a finite set of pairs,
      and \(\sum_{\lambda,\mu}g(\lambda,\mu)=N\);
\item\label{itm:ch50-jo:basis}
      there is a basis \(e_1,\dots,e_N\) of \(W\) adapted simultaneously to
      \(F^\bullet_\chi\) and \(F^\bullet_{\chi'}\); for such a basis
      \eqref{eq:ch50-orthogonal} holds for both norms, and the multiset
      \[
       \Sigma(\chi,\chi'):=\bigl\{(\chi(e_j),\chi'(e_j))\bigr\}_{j=1}^N
      \]
      is independent of the choice: it is the multiset in which
      \((\lambda,\mu)\) occurs with multiplicity \(g(\lambda,\mu)\).
\end{enumerate}
\end{theorem}

\begin{proof}
\emph{Step 1: nonnegativity and the dimension count.}  Both
\(W^{\lambda^+,\mu}\) and \(W^{\lambda,\mu^+}\) are subspaces of
\(W^{\lambda,\mu}\), so \(g\ge0\).  Since each filtration has finitely many
jumps, \(W^{\lambda,\mu}\) takes finitely many values and \(g\) is
supported on the finite set of pairs of jumps.  Next,
\[
 W^{\lambda^+,\mu}\cap W^{\lambda,\mu^+}
 =F^{\lambda^+}_\chi\cap F^\mu_{\chi'}\cap F^\lambda_\chi\cap
     F^{\mu^+}_{\chi'}
 =W^{\lambda^+,\mu^+},
\]
because \(F^{\lambda^+}_\chi\subset F^\lambda_\chi\) and
\(F^{\mu^+}_{\chi'}\subset F^\mu_{\chi'}\).  Hence, by the dimension
formula for a sum of two subspaces,
\[
 g(\lambda,\mu)
 =\dim W^{\lambda,\mu}-\dim W^{\lambda^+,\mu}
    -\dim W^{\lambda,\mu^+}+\dim W^{\lambda^+,\mu^+}.
\]
Fix \(\lambda\) and sum over \(\mu\).  The family
\(\mu\mapsto\dim W^{\lambda,\mu}\) is decreasing, equal to
\(\dim F^\lambda_\chi\) for \(\mu\) small and to \(0\) for \(\mu\) large,
so \(\sum_\mu(\dim W^{\lambda,\mu}-\dim W^{\lambda,\mu^+})
=\dim F^\lambda_\chi\); likewise with \(\lambda\) replaced by
\(\lambda^+\).  Therefore
\(\sum_\mu g(\lambda,\mu)=\dim F^\lambda_\chi-\dim F^{\lambda^+}_\chi\),
and summing over \(\lambda\) gives \(N\).

\emph{Step 2: construction of the basis.}  For each pair \((\lambda,\mu)\)
with \(g(\lambda,\mu)>0\) choose vectors
\(e_{\lambda,\mu,1},\dots,e_{\lambda,\mu,g(\lambda,\mu)}\in W^{\lambda,\mu}\)
whose classes form a basis of the quotient in \eqref{eq:ch50-bigraded}.  We
claim that for every \((\lambda_0,\mu_0)\) the space
\(W^{\lambda_0,\mu_0}\) is spanned by those \(e_{\lambda,\mu,i}\) with
\(\lambda\ge\lambda_0\) and \(\mu\ge\mu_0\).  Since there are finitely many
jumps, we may argue by descending induction on the pair
\((\lambda_0,\mu_0)\) in the finite set of pairs of jumps, the case of the
largest pair being the definition.  In general
\(W^{\lambda_0,\mu_0}=\bigl(W^{\lambda_0^+,\mu_0}+W^{\lambda_0,\mu_0^+}
\bigr)+\operatorname{span}\{e_{\lambda_0,\mu_0,i}\}\) by the choice of the
\(e_{\lambda_0,\mu_0,i}\), and the two summands in brackets are spanned by
admissible vectors by the inductive hypothesis, all of which have
\(\lambda\ge\lambda_0\) and \(\mu\ge\mu_0\).  This proves the claim.
Taking \((\lambda_0,\mu_0)\) small shows that the chosen vectors span
\(W\); by Step 1 their number is \(N\), so they form a basis.

\emph{Step 3: the basis is adapted to both, and the spectrum is
intrinsic.}  Let \(e=e_{\lambda,\mu,i}\).  Then \(\chi(e)\ge\lambda\) by
construction, and if \(\chi(e)>\lambda\) then \(e\in W^{\lambda^+,\mu}\),
so its class in \eqref{eq:ch50-bigraded} vanishes, a contradiction.  Hence
\(\chi(e)=\lambda\), and likewise \(\chi'(e)=\mu\).  Adaptedness to
\(F^\bullet_\chi\) is the case \(\mu_0\) small of the claim in Step 2, and
adaptedness to \(F^\bullet_{\chi'}\) is the case \(\lambda_0\) small.
Finally, for \emph{any} jointly adapted basis \(f_1,\dots,f_N\) one has,
for all \((\lambda,\mu)\),
\[
 W^{\lambda,\mu}
 =\operatorname{span}\{f_j:\chi(f_j)\ge\lambda,\ \chi'(f_j)\ge\mu\},
\]
because a vector lies in \(F^\lambda_\chi\) if and only if all its nonzero
coordinates are on basis vectors with \(\chi(f_j)\ge\lambda\), by
\eqref{eq:ch50-orthogonal}, and similarly for \(\chi'\).  Hence
\(\dim W^{\lambda,\mu}=\#\{j:\chi(f_j)\ge\lambda,\chi'(f_j)\ge\mu\}\), and
the multiset \(\{(\chi(f_j),\chi'(f_j))\}\) is determined by these numbers,
namely by \(g\).
\end{proof}

\begin{definition}[Relative spectrum]
\label{def:ch50-relative-spectrum}
\index{relative spectrum}
For norms \(\chi,\chi'\) on a finite-dimensional \(W\), the
\emph{relative spectrum} is the multiset
\[
 \{\chi(e_j)-\chi'(e_j)\}_{j=1}^N
\]
computed in any jointly orthogonal basis; it is well defined by
\Cref{thm:ch50-joint-orthogonal}\,\cref{itm:ch50-jo:basis}.
\end{definition}

\begin{remark}[Three norms do not diagonalize simultaneously]
\label{rem:ch50-three-flags}
\Cref{thm:ch50-joint-orthogonal} has no analogue for three norms: three
complete flags in a three-dimensional space need not admit a common adapted
basis, as the classical failure of simultaneous refinement for three
generic flags shows.  This is exactly why the triangle inequality for the
distances below is a theorem and not a computation, and why the proof of
\Cref{thm:ch50-d1-integral} passes through a dimension inequality for three
subspaces rather than through a common basis.
\end{remark}

\subsection{The spectral distances in finite dimension}
\label{sec:ch50-dp-finite}

If \(\chi,\chi'\) are norms on a finite-dimensional vector space, a joint
orthogonal basis \(e_1,\dots,e_N\) gives the relative spectrum
\((\chi(e_j)-\chi'(e_j))_{j=1}^N\).  Put
\[
 d_p(\chi,\chi')
 =\left(\frac1N\sum_{j=1}^N
       |\chi(e_j)-\chi'(e_j)|^p\right)^{1/p},
\quad
 d_\infty(\chi,\chi')=\max_j|\chi(e_j)-\chi'(e_j)|.
\tag{50.7}\label{eq:stability-finite-dimensional-dp}
\]
For graded norms, divide by \(m\) on \(R_m\) and pass to the limit
(\(\limsup\) for \(p=\infty\)).  These are pseudometrics on
\(\mathcal N_{\mathbf R}\).

We prove this for \(p=1\) and \(p=\infty\) and record the general case as
an import.  Two auxiliary norms are needed.  Given norms \(\chi,\chi'\) on
\(W\), define \(\chi\wedge\chi'\) and \(\chi\vee\chi'\) by the filtrations
\[
 F^\lambda_{\chi\wedge\chi'}:=F^\lambda_\chi\cap F^\lambda_{\chi'},
 \qquad
 F^\lambda_{\chi\vee\chi'}:=F^\lambda_\chi+F^\lambda_{\chi'}.
\tag{50.8}\label{eq:ch50-lattice-norms}
\]
Both families are decreasing, left continuous, exhaustive and separated, so
by \Cref{lem:ch50-norm-is-filtration} they are norms.

\begin{lemma}
\label{lem:ch50-lattice}
Let \(\chi,\chi'\) be norms on \(W\), with a joint orthogonal basis
\((e_j)\).  Then \((e_j)\) is orthogonal for \(\chi\wedge\chi'\) and for
\(\chi\vee\chi'\), and
\[
 (\chi\wedge\chi')(e_j)=\min\{\chi(e_j),\chi'(e_j)\},
 \qquad
 (\chi\vee\chi')(e_j)=\max\{\chi(e_j),\chi'(e_j)\}.
\]
Moreover \(\chi\wedge\chi'=\min\{\chi,\chi'\}\) pointwise on \(W\).
\end{lemma}

\begin{proof}
In the joint basis, \(F^\lambda_\chi=\operatorname{span}\{e_j:\chi(e_j)\ge
\lambda\}\) and similarly for \(\chi'\); the intersection of two coordinate
subspaces is the coordinate subspace on the intersection of the index sets,
and the sum is the one on the union.  Thus
\(F^\lambda_{\chi\wedge\chi'}=\operatorname{span}\{e_j:\min\ge\lambda\}\)
and \(F^\lambda_{\chi\vee\chi'}=\operatorname{span}\{e_j:\max\ge\lambda\}\),
which is the assertion.  For the last statement, apply
\eqref{eq:ch50-orthogonal} to all three norms in the basis \((e_j)\):
for \(w=\sum a_je_j\ne0\),
\[
 \begin{split}
 (\chi\wedge\chi')(w)
 &=\min_{a_j\ne0}\min\{\chi(e_j),\chi'(e_j)\}\\
 &=\min\Bigl\{\min_{a_j\ne0}\chi(e_j),\ \min_{a_j\ne0}\chi'(e_j)\Bigr\}
 =\min\{\chi(w),\chi'(w)\}.
 \end{split}
\]
\end{proof}

\begin{proposition}
\label{prop:ch50-dinfty-sup}
For norms \(\chi,\chi'\) on a finite-dimensional \(W\ne0\),
\[
 d_\infty(\chi,\chi')=\sup_{0\ne w\in W}|\chi(w)-\chi'(w)|
 =\min\{C\ge0:\ \chi-C\le\chi'\le\chi+C\ \text{on }W\setminus0\}.
\tag{50.9}\label{eq:ch50-dinfty-sup}
\]
In particular \(d_\infty\) is a metric on the set of norms on \(W\).
\end{proposition}

\begin{proof}
Let \((e_j)\) be a joint orthogonal basis and put
\(C:=\max_j|\chi(e_j)-\chi'(e_j)|=d_\infty(\chi,\chi')\).  For
\(w=\sum a_je_j\ne0\), \eqref{eq:ch50-orthogonal} gives
\[
 \chi(w)=\min_{a_j\ne0}\chi(e_j)
 \ \ge\ \min_{a_j\ne0}\bigl(\chi'(e_j)-C\bigr)=\chi'(w)-C,
\]
and symmetrically; so \(\sup_w|\chi-\chi'|\le C\).  Conversely
\(w=e_j\) realizes \(|\chi(e_j)-\chi'(e_j)|\), so the supremum is at least
\(C\).  The second equality in \eqref{eq:ch50-dinfty-sup} is a
restatement.  Symmetry and the triangle inequality are then immediate from
the supremum description, and \(d_\infty(\chi,\chi')=0\) forces
\(\chi=\chi'\).
\end{proof}

\begin{theorem}
\label{thm:ch50-d1-integral}
Let \(\chi,\chi'\) be norms on a \(k\)-vector space \(W\) of dimension
\(N\ge1\).  Then
\[
 N\,d_1(\chi,\chi')
 =\int_{\mathbf R}\Bigl(\dim F^\lambda_\chi+\dim F^\lambda_{\chi'}
       -2\dim\bigl(F^\lambda_\chi\cap F^\lambda_{\chi'}\bigr)\Bigr)
   \,\rd\lambda ,
\tag{50.10}\label{eq:ch50-d1-integral}
\]
the integrand being a compactly supported nonnegative step function.
Equivalently, writing \(\operatorname{tr}(\chi):=\sum_j\lambda_j(\chi)\) for
the sum of the jumping values,
\[
 N\,d_1(\chi,\chi')
 =\operatorname{tr}(\chi)+\operatorname{tr}(\chi')
   -2\operatorname{tr}(\chi\wedge\chi').
\tag{50.11}\label{eq:ch50-d1-trace}
\]
Consequently \(d_1\) is a metric on the set of norms on \(W\).
\end{theorem}

\begin{proof}
\emph{Step 1: the trace formula.}  Fix \(c>0\) so large that all jumping
values of \(\chi\) and \(\chi'\) lie in \([-c,\infty)\).  For any norm
\(\eta\) with jumping values \(\mu_1\ge\cdots\ge\mu_N\ge-c\),
\[
 \operatorname{tr}(\eta)+Nc
 =\sum_j(\mu_j+c)
 =\int_{-c}^{\infty}\#\{j:\mu_j\ge\lambda\}\,\rd\lambda
 =\int_{-c}^{\infty}\dim F^\lambda_\eta\,\rd\lambda,
\tag{50.12}\label{eq:ch50-trace-layer}
\]
by Fubini for the step function \(\lambda\mapsto\dim F^\lambda_\eta\).
Applying this to \(\chi\), \(\chi'\) and \(\chi\wedge\chi'\), whose jumping
values also lie in \([-c,\infty)\), and forming the combination
\eqref{eq:ch50-d1-trace}, the three occurrences of \(Nc\) cancel because
the coefficients \(1,1,-2\) sum to zero.  This gives
\eqref{eq:ch50-d1-integral} from \eqref{eq:ch50-d1-trace}, since
\(F^\lambda_{\chi\wedge\chi'}=F^\lambda_\chi\cap F^\lambda_{\chi'}\) by
definition.  In turn \eqref{eq:ch50-d1-trace} holds because, in a joint
orthogonal basis,
\[
 \sum_j\bigl|\chi(e_j)-\chi'(e_j)\bigr|
 =\sum_j\Bigl(\chi(e_j)+\chi'(e_j)
      -2\min\{\chi(e_j),\chi'(e_j)\}\Bigr),
\]
and \Cref{lem:ch50-lattice} identifies the last minimum with the jumping
values of \(\chi\wedge\chi'\), while
\Cref{lem:ch50-norm-is-filtration}\,\cref{itm:ch50-nf:adapted} identifies
the multisets \(\{\chi(e_j)\}\), \(\{\chi'(e_j)\}\) with the jumping
values.

\emph{Step 2: the triangle inequality.}  Let \(\chi_1,\chi_2,\chi_3\) be
norms on \(W\).  By \eqref{eq:ch50-d1-integral} the inequality
\(d_1(\chi_1,\chi_3)\le d_1(\chi_1,\chi_2)+d_1(\chi_2,\chi_3)\) is
equivalent, after cancelling \(\dim F_{\chi_1}^\lambda\) and
\(\dim F^\lambda_{\chi_3}\) and dividing by \(2\), to
\[
 \dim\bigl(F^\lambda_{\chi_1}\cap F^\lambda_{\chi_2}\bigr)
 +\dim\bigl(F^\lambda_{\chi_2}\cap F^\lambda_{\chi_3}\bigr)
 \ \le\
 \dim F^\lambda_{\chi_2}
 +\dim\bigl(F^\lambda_{\chi_1}\cap F^\lambda_{\chi_3}\bigr)
\tag{50.13}\label{eq:ch50-three-subspaces}
\]
for every \(\lambda\).  This is a statement about three subspaces
\(A,B,C\) of \(W\), namely
\(\dim(A\cap B)+\dim(B\cap C)\le\dim B+\dim(A\cap C)\), and it holds
because
\[
 \begin{split}
 \dim(A\cap B)+\dim(B\cap C)
 &=\dim\bigl((A\cap B)+(B\cap C)\bigr)+\dim(A\cap B\cap C)\\
 &\le\dim B+\dim(A\cap C),
 \end{split}
\]
the first equality being the dimension formula and the last inequality
holding because \((A\cap B)+(B\cap C)\subset B\) and
\(A\cap B\cap C\subset A\cap C\).

\emph{Step 3: separation.}  If \(d_1(\chi,\chi')=0\) then the relative
spectrum vanishes, so \(\chi(e_j)=\chi'(e_j)\) for a joint orthogonal
basis, whence \(\chi=\chi'\) by \eqref{eq:ch50-orthogonal}.  Symmetry is
clear.
\end{proof}

\begin{remark}[Reading \eqref{eq:ch50-d1-integral}]
\label{rem:ch50-d1-reading}
The integrand of \eqref{eq:ch50-d1-integral} is the codimension of
\(F^\lambda_\chi\cap F^\lambda_{\chi'}\) inside
\(F^\lambda_\chi+F^\lambda_{\chi'}\); it measures, level by level, how far
the two filtrations are from agreeing.  This description makes two things
evident that are not evident from \eqref{eq:stability-finite-dimensional-dp}:
\(d_1\) is monotone under refinement of either filtration, and it is
additive in direct sums of filtered vector spaces.  It is also the form in
which \(d_1\) globalizes: replacing \(\dim\) by an asymptotic normalized
dimension turns \eqref{eq:ch50-d1-integral} into the volume formula
\eqref{eq:stability-d1-volume-energy}.
\end{remark}

\begin{proposition}
\label{prop:ch50-dp-interpolation}
For norms on a finite-dimensional \(W\ne0\) and \(1\le p\le q\le\infty\),
\[
 d_1(\chi,\chi')\ \le\ d_p(\chi,\chi')\ \le\ d_q(\chi,\chi')
 \ \le\ d_\infty(\chi,\chi').
\tag{50.14}\label{eq:ch50-dp-interpolation}
\]
The same inequalities hold for graded norms after passing to the limit.
\end{proposition}

\begin{proof}
This is H\"older's inequality for the uniform probability measure on
\(\{1,\dots,N\}\): the \(L^p\) norm with respect to a probability measure
is nondecreasing in \(p\).  The passage to graded norms preserves
inequalities between limits.
\end{proof}

\begin{remark}[The intermediate \(d_p\) are pseudometrics: an import]
\label{rem:ch50-dp-import}
That \(d_p\) satisfies the triangle inequality for \(1<p<\infty\) is
proved in \cite{BJKS1} (preprint 2021, published 2025).  It is not needed
anywhere below: the stability functionals of this chapter are continuous
for \(d_1\), and the only other distance used is \(d_\infty\).  We
therefore do not reproduce the proof, and we flag the statement as
imported in \Cref{imp:ch50-imported}\,\cref{itm:ch50-import:dp}.  It is
worth noting where the difficulty lies: by
\Cref{rem:ch50-three-flags} there is no common basis for three norms, and
the argument of Step 2 above works for \(p=1\) only because the
\(\ell^1\)-norm of the relative spectrum is an integral of dimensions, a
linearity that no other exponent enjoys.
\end{remark}

\subsection{The spectral distances on graded norms}
\label{sec:ch50-dp-graded}

For \(\chi,\chi'\in\mathcal N_{\mathbf R}\), apply the finite-dimensional
construction to \(W=R_m\), rescale by \(m\), and pass to the limit:
\[
 \begin{gathered}
 d_p(\chi,\chi'):=\lim_{m\to\infty}\frac1m\,d_p(\chi_m,\chi'_m)
 \quad(1\le p<\infty),\\
 d_\infty(\chi,\chi'):=\limsup_{m\to\infty}
      \frac1m\,d_\infty(\chi_m,\chi'_m).
 \end{gathered}
\tag{50.15}\label{eq:ch50-dp-graded}
\]

\begin{proposition}
\label{prop:ch50-dp-graded}
Let \(\chi,\chi'\in\mathcal N_{\mathbf R}\).
\begin{enumerate}
\item\label{itm:ch50-dpg:infty}
      The quantity \(d_\infty(\chi,\chi')\) of
      \eqref{eq:ch50-dp-graded} is finite, and it is a pseudometric on
      \(\mathcal N_{\mathbf R}\).  If \(\chi\) and \(\chi'\) are
      both homogeneous, then \(m\mapsto m^{-1}d_\infty(\chi_m,\chi'_m)\)
      is nondecreasing along divisibility, so the \(\limsup\) is a limit
      and equals \(\sup_m m^{-1}d_\infty(\chi_m,\chi'_m)\).
\item\label{itm:ch50-dpg:one}
      The limit defining \(d_1\) exists, is finite, and \(d_1\) is a
      pseudometric on \(\mathcal N_{\mathbf R}\).
\item\label{itm:ch50-dpg:interp}
      \(d_1\le d_p\le d_\infty\) for all \(p\), and all of them are
      invariant under simultaneous translation
      \((\chi,\chi')\mapsto(\chi+c,\chi'+c)\) and homogeneous of degree
      one under the base change \(\chi\mapsto\chi^{[d]}\) of
      \Cref{prop:ch33-transformations}.
\end{enumerate}
\end{proposition}

\begin{proof}
\cref{itm:ch50-dpg:infty}  By \Cref{prop:ch50-dinfty-sup},
\(d_\infty(\chi_m,\chi'_m)=\sup_{0\ne s\in R_m}|\chi_m(s)-\chi'_m(s)|\).
Linear boundedness \eqref{eq:stability-norm-axioms} bounds each of
\(|\chi_m(s)|\) and \(|\chi'_m(s)|\) by \(Cm\), so
\(m^{-1}d_\infty(\chi_m,\chi'_m)\le2C\) and the \(\limsup\) is finite.
Symmetry is clear, and the triangle inequality passes to the \(\limsup\):
for three norms, \(d_\infty(\chi_m,\chi''_m)\le
d_\infty(\chi_m,\chi'_m)+d_\infty(\chi'_m,\chi''_m)\) in each degree by
\Cref{prop:ch50-dinfty-sup}, and \(\limsup(a_m+b_m)\le\limsup a_m
+\limsup b_m\).

If \(\chi\) and \(\chi'\) are homogeneous, then for \(r\ge1\) and
\(0\ne s\in R_m\),
\[
 \bigl|\chi_{rm}(s^r)-\chi'_{rm}(s^r)\bigr|
 =r\bigl|\chi_m(s)-\chi'_m(s)\bigr| ,
\]
so \(d_\infty(\chi_{rm},\chi'_{rm})\ge r\,d_\infty(\chi_m,\chi'_m)\) and
therefore
\[
 (rm)^{-1}d_\infty(\chi_{rm},\chi'_{rm})
 \ge m^{-1}d_\infty(\chi_m,\chi'_m).
\]
A nondecreasing net along a directed set
converges to its supremum, so the \(\limsup\) is a limit and equals the
supremum.

\cref{itm:ch50-dpg:one}  Existence of the limit is
\Cref{prop:ch50-d1-volume} below.  It expresses
\[
 N_m^{-1}d_1(\chi_m,\chi'_m)
\]
as a difference of normalized traces, each of which has a limit by
\Cref{thm:dh-limit-filtration}.  The triangle
inequality and symmetry pass to the limit from
\Cref{thm:ch50-d1-integral}.

\cref{itm:ch50-dpg:interp}  The inequalities pass to the limit from
\eqref{eq:ch50-dp-interpolation}.  Translation replaces every jumping value
in degree \(m\) by itself plus \(mc\), hence does not change the relative
spectrum; base change of degree \(d\) multiplies every jumping value by
\(d\), by \Cref{prop:ch33-transformations}\,(2).
\end{proof}

\begin{warning}{\(d_1\) is a pseudometric, not a metric}
\label{warn:ch50-pseudometric}
On \(\mathcal N_{\mathbf R}\) the distances \eqref{eq:ch50-dp-graded} are
only pseudometrics, and their kernels are genuinely nontrivial: distinct
norms can be at distance zero.  Two mechanisms produce this, and they
should not be conflated.  The first is homogenization
(\Cref{prop:ch50-homogenization}): \(\chi\) and \(\chi^{\hom}\) have the
same Fubini--Study function and, by
\Cref{rem:stability-three-equivalences}, satisfy
\(d_1(\chi,\chi^{\hom})=0\).  The second is asymptotic equivalence in
degree: changing \(\chi_m\) by \(o(m)\) uniformly changes nothing in the
limit.  Every statement below about ``the'' space of norms therefore refers
to a Hausdorff quotient, and we say which one.
\end{warning}

\subsection{The Fubini--Study operator and its adjoint}
\label{sec:ch50-fs-adjoint}

The bridge to the analytic side is the operator of
\eqref{eq:norm-fs-function}:
\[
 \FS(\chi)=\phi_\chi,\qquad
 \phi_\chi(v)=\sup_{\substack{m\geq1\\0\ne s\in R_m}}
        \frac{\chi_m(s)-v(s)}{m}
 \qquad(v\in X^{\An}).
\tag{50.16}\label{eq:ch50-fs-operator}
\]
By \Cref{prop:ch33-base-ideal-flag}\,(3) the supremum is an increasing
supremum of the functions \(m^{-1}\phi_{\mathfrak a_{(m)}}\) attached to
the degree-\(m\) flag ideals, each of which is a Fubini--Study function for
the class \(\omega\); hence \(\phi_\chi\in\PSH(\omega)\) by
\Cref{def:bj-psh} and \Cref{thm:bj-basic-closure}.  In the opposite
direction, for a function \(\phi:X^{\An}\to\mathbf R\cup\{-\infty\}\) which
is bounded above, put
\[
 \chi^\phi_m(s):=\inf_{v\in X^{\An}}\bigl\{v(s)+m\phi(v)\bigr\}
 \qquad(0\ne s\in R_m),\qquad \chi^\phi_m(0):=+\infty .
\tag{50.17}\label{eq:ch50-supnorm}
\]
In multiplicative notation \(\chi^\phi_m(s)=-\log\sup_{X^{\An}}
\|s\|_{m\phi}\): it is the sup-norm of \(s\) for the metric \(m\phi\).

\begin{proposition}
\label{prop:ch50-supnorm}
Let \(\phi\in\PSH(\omega)\) be bounded.  Then \(\chi^\phi\) is a linearly
bounded multiplicative norm on \(R\) which is moreover
\emph{homogeneous}:
\[
 \chi^\phi_{rm}(s^r)=r\,\chi^\phi_m(s)\qquad(r\ge1).
\tag{50.18}\label{eq:ch50-supnorm-homog}
\]
\end{proposition}

\begin{proof}
Finiteness and \(k^\times\)-invariance are clear, the latter because
\(v(as)=v(s)\) for \(a\in k^\times\).  For the ultrametric inequality, fix
\(v\); then \(v(s+s')\ge\min\{v(s),v(s')\}\) gives
\(v(s+s')+m\phi(v)\ge\min\{v(s)+m\phi(v),\,v(s')+m\phi(v)\}
\ge\min\{\chi^\phi_m(s),\chi^\phi_m(s')\}\), and taking the infimum over
\(v\) preserves the bound.  For multiplicativity,
\[
 v(ss')+(m+m')\phi(v)
 =\bigl(v(s)+m\phi(v)\bigr)+\bigl(v(s')+m'\phi(v)\bigr)
 \ \ge\ \chi^\phi_m(s)+\chi^\phi_{m'}(s'),
\]
again for every \(v\), whence \(\chi^\phi_{m+m'}(ss')\ge\chi^\phi_m(s)
+\chi^\phi_{m'}(s')\).  Homogeneity \eqref{eq:ch50-supnorm-homog} is the
same computation with \(s'=s^{r-1}\), read as an equality: \(v(s^r)
=rv(s)\), so the function being minimized is \(r\) times the one for
\(s\), and the infima agree.  For linear boundedness, note that
\(v(s)\ge0\) for every \(v\in X^{\An}\) and every regular section, while
\(v_{\triv}(s)=0\); taking \(v=v_{\triv}\) in \eqref{eq:ch50-supnorm} and
bounding below by \(\inf\phi\) gives
\[
 m\inf_{X^{\An}}\phi\ \le\ \chi^\phi_m(s)\ \le\ m\,\phi(v_{\triv})
 =m\sup_{X^{\An}}\phi
\]
for every \(0\ne s\in R_m\), which is \eqref{eq:stability-norm-axioms} with
\(C=\sup|\phi|\).
\end{proof}

\begin{theorem}
\label{thm:ch50-galois}
For \(\chi\in\mathcal N_{\mathbf R}\) and \(\phi\in\PSH(\omega)\) bounded,
\[
 \FS(\chi)\le\phi
 \quad\Longleftrightarrow\quad
 \chi\le\chi^\phi .
\tag{50.19}\label{eq:ch50-galois}
\]
Consequently \(\FS\) and \(\phi\mapsto\chi^\phi\) are order reversing on
neither side but order preserving, \(\chi\mapsto\chi^{\FS(\chi)}\) is a
closure operator on \(\mathcal N_{\mathbf R}\), and
\[
 \FS\bigl(\chi^{\FS(\chi)}\bigr)=\FS(\chi).
\tag{50.20}\label{eq:ch50-closure}
\]
\end{theorem}

\begin{proof}
Both sides of \eqref{eq:ch50-galois} say the same thing, namely that
\(\chi_m(s)\le v(s)+m\phi(v)\) for all \(m\), all \(0\ne s\in R_m\) and all
\(v\in X^{\An}\): the left side is this statement read as
\(\phi(v)\ge m^{-1}(\chi_m(s)-v(s))\) with a supremum over \((m,s)\), the
right side reads it as \(\chi_m(s)\le\inf_v(v(s)+m\phi(v))\).  No
cancellation of infinities occurs, because \(v(s)\in[0,+\infty]\) and
\(\phi\) is real valued and bounded.

Given \eqref{eq:ch50-galois}, the composite \(C(\chi):=\chi^{\FS(\chi)}\)
satisfies \(\chi\le C(\chi)\) (take \(\phi=\FS(\chi)\)), is monotone
(because both operators are), and satisfies
\(\FS(C(\chi))\le\FS(\chi)\) by \eqref{eq:ch50-galois} applied to
\(C(\chi)\le\chi^{\FS(\chi)}\); together with monotonicity of \(\FS\) this
gives \eqref{eq:ch50-closure} and \(C\circ C=C\).
\end{proof}

\begin{proposition}
\label{prop:ch50-homogenization}
Let \(\chi\in\mathcal N_{\mathbf R}\).  For \(0\ne s\in R_m\) the limit
\[
                 \chi^{\hom}_m(s)
                  :=\lim_{r\to\infty}\frac1r\chi_{rm}(s^r)
                  =\sup_{r\ge1}\frac1r\chi_{rm}(s^r)
\tag{50.21}\label{eq:ch50-homogenization-def}
\]
exists and is finite, \(\chi\le\chi^{\hom}\le\chi^{\FS(\chi)}\), and
\[
 \FS(\chi^{\hom})=\FS(\chi).
\tag{50.22}\label{eq:ch50-fs-homog}
\]
\end{proposition}

\begin{proof}
Superadditivity \eqref{eq:stability-norm-axioms} gives
\(\chi_{(r+r')m}(s^{r+r'})\ge\chi_{rm}(s^r)+\chi_{r'm}(s^{r'})\), so
Fekete's lemma applies and the limit equals the supremum; it is finite
because \(\chi_{rm}(s^r)\le Crm\).  Taking \(r=1\) gives
\(\chi\le\chi^{\hom}\).  For the upper bound and for
\eqref{eq:ch50-fs-homog}, write \(\phi:=\FS(\chi)\).  For every \(v\) and
every \(r\),
\[
 \frac1r\chi_{rm}(s^r)
 \ \le\ \frac1r\bigl(v(s^r)+rm\phi(v)\bigr)
 = v(s)+m\phi(v),
\]
using \(\chi\le\chi^\phi\) from \eqref{eq:ch50-galois} in degree \(rm\) and
\(v(s^r)=rv(s)\).  Taking the supremum over \(r\) and then the infimum over
\(v\) gives \(\chi^{\hom}_m(s)\le\chi^\phi_m(s)\).  Monotonicity of \(\FS\)
now gives \(\FS(\chi)\le\FS(\chi^{\hom})\le\FS(\chi^\phi)=\FS(\chi)\) by
\eqref{eq:ch50-closure}.
\end{proof}

\begin{remark}[Three equivalence relations]
\label{rem:stability-three-equivalences}
The following operations should not be confused.
\begin{enumerate}
\item Translation is
      \((\chi+c)_m(s)=\chi_m(s)+mc\); it changes the linearization and adds
      \(c\) to the corresponding potential.
\item Homogenization is
      \[
                 \chi^{\hom}(s)
                  :=\lim_{r\to\infty}r^{-1}\chi(s^r).
      \tag{50.23}\label{eq:stability-homogenization}
      \]
      It is the least homogeneous norm dominating \(\chi\).
\item Asymptotic equivalence means distance zero for one, hence every
      finite, spectral \(d_p\)-pseudometric.  Homogenization does not
      identify norms pointwise, but
      \(d_1(\chi,\chi^{\hom})=0\).
\end{enumerate}
\end{remark}

\begin{proof}[Discussion and provenance]
Items (1) and (2) are \Cref{prop:ch33-transformations}\,(1) and
\Cref{prop:ch50-homogenization}; the minimality assertion in (2) holds
because a homogeneous norm \(\eta\ge\chi\) satisfies
\(\eta_m(s)=r^{-1}\eta_{rm}(s^r)\ge r^{-1}\chi_{rm}(s^r)\) for all \(r\).
The vanishing \(d_1(\chi,\chi^{\hom})=0\) in (3) is
proved in \cite{BJKS1} (preprint 2021, published 2025); by
\Cref{prop:ch50-d1-volume} below it is equivalent to the equality of the
two asymptotic normalized traces, that is, to
\(\operatorname{vol}(\chi)=\operatorname{vol}(\chi^{\hom})\) in the
notation of \eqref{eq:stability-volume-norm}, and this is where the
argument of that paper is needed.  We use the statement only through
\Cref{thm:stability-fs-quotient}, which is imported in any case.  That
pointwise equality fails is already visible in
\Cref{ex:ch50-nonhomogeneous} below.
\end{proof}

\begin{example}[A norm that is not homogeneous]
\label{ex:ch50-nonhomogeneous}
Let \(X=\mathbf P^1\), \(L=\cO(1)\) and let \(\chi\) be the norm whose
filtration is \(\mathcal F^\lambda R_m=R_m\) for \(\lambda\le0\),
\(\mathcal F^\lambda R_m=\operatorname{span}\{x^ay^{m-a}:a\ge1\}\) for
\(0<\lambda\le1\), and \(0\) for \(\lambda>1\); that is, \(\chi_m(s)=1\)
if \(x\mid s\) and \(0\) otherwise, independently of \(m\).  This is
\emph{not} a norm: superadditivity \eqref{eq:stability-norm-axioms} fails
already at \(s=s'=x\), \(m=m'=1\), where
\(\chi_2(x^2)=1<2=\chi_1(x)+\chi_1(x)\); equivalently the filtration is
not multiplicative,
\(\mathcal F^1R_1\cdot\mathcal F^1R_1\not\subseteq\mathcal F^2R_2=0\).
A correct example on the same \(X\) and \(L\) is
\[
 \chi_m(x^ay^{m-a}):=\bigl\lfloor a/2\bigr\rfloor,
\]
extended to a general \(0\ne s\in R_m\) by the minimum over the monomials
occurring in \(s\).  All three axioms hold: the non-Archimedean inequality
by construction, superadditivity because
\(\lfloor(a+a')/2\rfloor\ge\lfloor a/2\rfloor+\lfloor a'/2\rfloor\), and
linear boundedness because \(0\le\chi_m\le m/2\).  Homogenization removes
the floor: \(\chi_{rm}\bigl((x^ay^{m-a})^r\bigr)=\lfloor ra/2\rfloor\), so
\(\chi^{\hom}_m(x^ay^{m-a})=a/2\), and \(\chi<\chi^{\hom}\) at every odd
\(a\), the gap being exactly \(\tfrac12\) --- consistently with
\(\chi\le\chi^{\hom}\) of \Cref{prop:ch50-homogenization}.  Yet
\(|\chi_m-\chi^{\hom}_m|\le\tfrac12\) uniformly in \(m\) and in \(s\), so
after the normalization by \(m\) in \eqref{eq:ch50-dp-graded} one gets
\(d_1(\chi,\chi^{\hom})=0\), and the two have the same
Duistermaat--Heckman measure and, by \eqref{eq:ch50-fs-homog}, the same
Fubini--Study function.  This is the
smallest example of the phenomenon recorded in
\Cref{warn:ch50-pseudometric}: the two norms differ in every odd degree,
but by \(O(1)\), and the normalization by \(m\) is all the theory sees.
\end{example}

\begin{lemma}
\label{lem:ch50-model-norm}
Let \((\cX,\cL)\) be a normal test configuration for \((X,L)\) with
\(\cL\) relatively semiample, let \(\phi=\phi_\cL\) be its
Fubini--Study function, and let \(\chi\) be the associated filtration
\eqref{eq:ch33-tc-filtration}.  Then, for every sufficiently
divisible \(m\) and every \(0\ne s\in R_m\),
\[
 \chi_m(s)=\min_{E\subset\cX_0}\bigl(v_E(s)+m\,\phi(v_E)\bigr),
\tag{50.24}\label{eq:ch50-model-norm}
\]
the minimum over the irreducible components of the central fibre.  In
particular \(\chi\) is homogeneous and \(\chi=\chi^{\phi}\), so that
\(\chi\) is a closed point of the closure operator of
\Cref{thm:ch50-galois}.
\end{lemma}

\begin{proof}
By \Cref{def:ch33-tc-filtration}, \(\chi_m(s)\) is the largest
\(\lambda\) for which \(\varpi^{-\lambda}\bar s\) is a regular section
of \(m\cL\) over \(\cX\).  Since \(\cX\) is normal and
\(\varpi^{-\lambda}\bar s\) is regular away from \(\cX_0\), regularity
is equivalent to nonnegativity of the order along each irreducible
component \(E\) of \(\cX_0\); this is
\Cref{lem:ch33-effectivity} applied to the vertical fractional ideal
generated by \(\varpi^{-\lambda}\bar s\) relative to \(m\cL\).  Writing
\(\cL=p^*L+D\) with \(D\) vertical, one has
\(\ord_E(\varpi^{-\lambda})=-\lambda b_E\),
\(b_E^{-1}\ord_E(p^*s)=v_E(s)\) and \(b_E^{-1}\ord_E(D)=\phi(v_E)\), the
first two by \eqref{eq:ch33-bE} and the third by
\Cref{lem:ch33-effectivity}.  Regarding \(\bar s\) as a rational section
of \(m\cL=mp^*L+mD\) adds \(mD\) to its divisor, so
\[
 \ord_E\bigl(\varpi^{-\lambda}\bar s\bigr)
 =-\lambda b_E+b_E\,v_E(s)+m\,b_E\,\phi(v_E),
\]
and the condition \(\ord_E(\varpi^{-\lambda}\bar s)\ge0\) for every \(E\)
reads \(\lambda\le v_E(s)+m\phi(v_E)\), which is
\eqref{eq:ch50-model-norm}.  As a check,
\Cref{ex:ch33-dnc-p1} has two components, with
\((v_E,\phi(v_E))=(\ord_p,-c)\) and
\((v_{\widetilde E},\phi(v_{\widetilde E}))=(v_{\triv},0)\), so
\eqref{eq:ch50-model-norm} gives
\(\chi_m(x^ay^{m-a})=\min\{a-mc,\,0\}\), which is
\eqref{eq:ch33-dnc-weights}.

Homogeneity is immediate from \eqref{eq:ch50-model-norm}, since
\(v_E(s^r)=rv_E(s)\).  Finally \(\chi\le\chi^\phi\) by
\eqref{eq:ch50-galois}, while \eqref{eq:ch50-model-norm} exhibits
\(\chi_m(s)\) as a minimum over a subset of \(X^{\An}\) of the very
quantity whose infimum over all of \(X^{\An}\) is \(\chi^\phi_m(s)\);
hence \(\chi_m(s)\ge\chi^\phi_m(s)\), and the two agree.
\end{proof}

\begin{proposition}
\label{prop:ch50-dinfty-fs}
Let \(\chi,\chi'\) be the filtrations of two normal test configurations
for \((X,L)\) with relatively semiample polarizations, and let
\(\phi=\FS(\chi)\), \(\phi'=\FS(\chi')\) be their Fubini--Study
functions.  Then \(m^{-1}d_\infty(\chi_m,\chi'_m)\) is nondecreasing along
divisibility, its limit exists, and
\[
 d_\infty(\chi,\chi')=\sup_{X^{\An}}|\phi-\phi'| .
\tag{50.25}\label{eq:ch50-dinfty-fs}
\]
\end{proposition}

\begin{proof}
Write \(c_m:=m^{-1}d_\infty(\chi_m,\chi'_m)
=m^{-1}\sup_{0\ne s\in R_m}|\chi_m(s)-\chi'_m(s)|\), by
\Cref{prop:ch50-dinfty-sup}, and \(C:=\sup_{X^{\An}}|\phi-\phi'|\), a
finite number because both functions are bounded.

\emph{\(c_m\le C\) for every \(m\).}  From \(\phi\le\phi'+C\) and
\eqref{eq:ch50-supnorm} we get \(\chi^{\phi}_m\le\chi^{\phi'}_m+mC\); by
\Cref{lem:ch50-model-norm}, \(\chi=\chi^\phi\) and \(\chi'=\chi^{\phi'}\),
so \(\chi_m\le\chi'_m+mC\).  Exchanging the two gives
\(|\chi_m-\chi'_m|\le mC\) on \(R_m\setminus0\), that is \(c_m\le C\).

\emph{\(\sup_mc_m\ge C\).}  If \(c_m\le C'\) for every \(m\), then
\(\chi_m(s)\le\chi'_m(s)+mC'\) for all \(m\) and \(s\), whence by
\eqref{eq:ch50-fs-operator} \(\phi\le\phi'+C'\); symmetrically
\(\phi'\le\phi+C'\), so \(C\le C'\).

\emph{Monotonicity.}  Both \(\chi\) and \(\chi'\) are homogeneous by
\Cref{lem:ch50-model-norm}, so
\Cref{prop:ch50-dp-graded}\,\cref{itm:ch50-dpg:infty} applies: \(c_m\) is
nondecreasing along divisibility and \(d_\infty(\chi,\chi')
=\lim_mc_m=\sup_mc_m\).  Combining the three displays gives
\eqref{eq:ch50-dinfty-fs}.
\end{proof}

\begin{remark}[From model norms to all finite-type norms]
\label{rem:ch50-dinfty-finite-type}
\Cref{prop:ch50-dinfty-fs} covers exactly the norms of normal semiample
test configurations, that is, the integral finite-type locus
\(\mathcal T_{\mathbf Z}\) after passing to a normal representative.  For
\(\mathcal T_{\mathbf R}\) the same statement holds, and the reduction is
by translation, base change and homogenization, none of which changes
either side of \eqref{eq:ch50-dinfty-fs}: the right-hand side because of
\Cref{prop:ch33-transformations}\,(4) and \eqref{eq:ch50-fs-homog}, the
left-hand side because of
\Cref{prop:ch50-dp-graded}\,\cref{itm:ch50-dpg:interp}.  The only point
requiring care is that a general finite-type norm need not satisfy
\(\chi=\chi^{\FS(\chi)}\) --- redundant sections in a presentation change
\(\chi\) without changing \(\FS(\chi)\) --- so one must first replace
\(\chi\) by its closure; the assertion that this does not change
\(d_\infty\) is part of \Cref{thm:stability-fs-quotient}, hence imported.
\end{remark}

\subsection{Volumes, and \texorpdfstring{\(d_1\)}{d1} in the limit}
\label{sec:ch50-volume}

\begin{proposition}
\label{prop:ch50-d1-volume}
For \(\chi\in\mathcal N_{\mathbf R}\) put
\[
 \operatorname{vol}(\chi)
 :=\lim_m\frac1m\frac1{N_m}
              \sum_{j=1}^{N_m}\lambda_{m,j}.
\tag{50.26}\label{eq:stability-volume-norm}
\]
The limit exists and equals \(\int_{\mathbf R}\lambda\,\rd\DHm_\chi\), by
\Cref{thm:dh-limit-filtration} and \eqref{eq:ch33-dh-moment-formula}, the
supports of the \(\nu_m\) lying in a fixed compact interval.  Then
\(\chi\wedge\chi'\), defined degreewise by \eqref{eq:ch50-lattice-norms},
is again a linearly bounded multiplicative norm, and
\[
 d_1(\chi,\chi')
 =\operatorname{vol}(\chi)+\operatorname{vol}(\chi')
       -2\operatorname{vol}(\chi\wedge\chi').
\tag{50.27}\label{eq:ch50-d1-volume}
\]
In particular the limit in \eqref{eq:ch50-dp-graded} exists for \(p=1\).
\end{proposition}

\begin{proof}
Multiplicativity of \(\chi\wedge\chi'=\min\{\chi,\chi'\}\)
(\Cref{lem:ch50-lattice}) is immediate:
\(\min\{\chi,\chi'\}_{m+m'}(ss')\ge\min\{\chi_m(s)+\chi_{m'}(s'),
\chi'_m(s)+\chi'_{m'}(s')\}\ge\min\{\chi,\chi'\}_m(s)
+\min\{\chi,\chi'\}_{m'}(s')\); linear boundedness and the ultrametric
inequality are clear.  Now divide \eqref{eq:ch50-d1-trace} in degree \(m\)
by \(mN_m\) and let \(m\to\infty\); each of the three normalized traces
converges to the corresponding \(\operatorname{vol}\) by
\Cref{thm:dh-limit-filtration} and \eqref{eq:ch33-dh-moment-formula},
the integrand \(\lambda\) being bounded on the supports.
\end{proof}

\begin{remark}[Volume, energy, and the first moment]
\label{rem:ch50-volume-energy}
Comparing \eqref{eq:stability-volume-norm} with
\Cref{prop:ch46-energy-dh} shows that for a norm coming from a test
configuration, \(\operatorname{vol}(\chi)=E_\omega(\FS(\chi))\); this is the
second line of \eqref{eq:stability-d1-volume-energy}, and it is proved in
full generality in \Cref{thm:ch50-energy-dh}.  Thus
\eqref{eq:ch50-d1-volume} says that the \(d_1\)-distance between two norms
is the energy of the first plus the energy of the second minus twice the
energy of their minimum: a formula which, on the analytic side, is the
statement that \(d_1\) is the Monge--Amp\`ere metric associated with the
concave functional \(E\).  The identification with the functional \(I\) of
\eqref{eq:bj-I-J-potentials} is not immediate from this and is part of
\Cref{thm:stability-fs-quotient}.
\end{remark}

\subsection{The Hausdorff quotient and the completion}
\label{sec:ch50-quotient}

\begin{theorem}\label{thm:stability-fs-quotient}
For every subgroup \(\Lambda\subset\mathbf R\) and
\(p\in[1,\infty]\), the Fubini--Study operator identifies the Hausdorff
quotient of \((\mathcal T_\Lambda,d_p)\) with the space
\(\mathcal H_\Lambda(L)\) of \(\Lambda\)-Fubini--Study functions.
For \(p=\infty\), the induced distance is the uniform distance.
\end{theorem}

\begin{proof}[Proof architecture]
For \(p=\infty\), the infimum-norm construction is inverse to
Fubini--Study after homogenization; the spectral-radius formula
\eqref{eq:stability-homogenization} proves the isometry.  For \(p=1\), set
\(\operatorname{vol}\) as in \eqref{eq:stability-volume-norm}.
The identities
\[
 \begin{split}
 d_1(\chi,\chi')
 &=\operatorname{vol}(\chi)+\operatorname{vol}(\chi')
       -2\operatorname{vol}(\chi\wedge\chi'),\\
 \operatorname{vol}(\chi)&=E(\FS(\chi))
 \qquad(\chi\text{ of finite type})
 \end{split}
\tag{50.28}\label{eq:stability-d1-volume-energy}
\]
reduce separation to strict monotonicity of energy.  The intermediate
\(d_p\)'s follow from \(d_1\leq d_p\leq d_\infty\).
This is \cite[Theorem~A]{BJKS1}; the formula for arbitrary norms uses the
upper energy extension in \cite[Theorem~5.1]{BJKS1}.
\end{proof}

\begin{remark}[What is proved here and what is quoted]
\label{rem:ch50-quotient-status}
The two displayed identities of \eqref{eq:stability-d1-volume-energy} are
proved in this book: the first is \Cref{prop:ch50-d1-volume} and the
second is \Cref{thm:ch50-energy-dh}.  The case \(p=\infty\) of
\Cref{thm:stability-fs-quotient} is \Cref{prop:ch50-dinfty-fs} together
with \Cref{prop:ch50-homogenization}.  What is imported is the
\emph{separation} statement for \(p=1\) --- that \(d_1(\chi,\chi')=0\)
forces \(\FS(\chi)=\FS(\chi')\) --- and the surjectivity of \(\FS\) onto
\(\mathcal H_\Lambda(L)\) for a general subgroup \(\Lambda\).  Separation
is where strict monotonicity of \(E\) is used, in the form
\Cref{cor:ch47-strict-monotone}; the printed proof in
\cite{BJKS1} runs through the upper energy extension of arbitrary
norms, which we have not constructed.  This is
\Cref{imp:ch50-imported}\,\cref{itm:ch50-import:quotient}.
\end{remark}

The completion has two faces.  On the potential side it is
\(\mathcal E^1(L)\), while on the measure side it is
\(\mathcal M^1(L)\).  The second description does not require the envelope
property.

Translations act isometrically on \((\mathcal N_{\mathbf R},d_1)\).  We
write
\begin{equation}
\label{eq:ch50-quotient-d1}
 \overline d_1([\chi],[\chi'])
 :=\inf_{c\in\mathbf R}d_1(\chi+c,\chi')
\tag{50.28a}
\end{equation}
for the resulting quotient pseudometric on
\(\mathcal N_{\mathbf R}/\mathbf R\), and similarly on
\(\mathcal T_{\mathbf R}/\mathbf R\) and
\(\mathcal E^1(L)/\mathbf R\).  Whenever one of these spaces is called a
metric space below, points at \(\overline d_1\)-distance zero have also been
identified.  This Hausdorffization is essential on the norm side: two
norms can be asymptotically equivalent modulo translation without being
equal modulo translation.

\begin{theorem}\label{thm:stability-norm-completion}
There are metrics \(d_1\) inducing the strong topologies on
\(\mathcal E^1(L)\) and \(\mathcal M^1(L)\) such that:
\begin{enumerate}
\item \(\mathcal H_{\mathbf R}(L)\) is \(d_1\)-dense in
      \(\mathcal E^1(L)\);
\item the Monge--Amp\`ere operator extends uniquely to an isometry with
      dense image
      \[
       \MA:(\mathcal N_{\mathbf R}/\mathbf R,\overline d_1)
               \longrightarrow(\mathcal M^1(L),d_1),
      \tag{50.29}\label{eq:stability-norm-ma-isometry}
      \]
      where the convention following \eqref{eq:ch50-quotient-d1} is in
      force;
\item \(\mathcal M^1(L)\) is complete, while \(\mathcal E^1(L)\) is
      complete if and only if the envelope property holds.
\end{enumerate}
Hence \(\mathcal M^1(L)\) is always the Hausdorff completion of
\(\mathcal T_{\mathbf R}/\mathbf R\).  Under the envelope property,
\(\mathcal E^1(L)\) is the completion before quotienting by constants.
\end{theorem}

\begin{proof}
The analytic metric assertions are deep inputs from
\cite[Theorem~B and Theorems~5.5, 5.11]{BJKS1}, built on the energy and
strong-topology estimates of \cite{BJ22}.  Precisely, we import the
following three facts.

First, \(d_1\) extends uniquely from \(\cH_{\mathbf R}(L)\) to a metric on
\(\cE^1(L)\) inducing the strong topology, and
\(\cH_{\mathbf R}(L)\) is dense.  Second, \(\overline d_1\) on
\(\cE^1(L)/\mathbf R\) extends uniquely across its Monge--Amp\`ere image to
a metric on \(\cM^1(L)\) inducing the strong topology; the embedding
\[
\label{eq:ch50-e1-ma-dense}
 \MA:(\cE^1(L)/\mathbf R,\overline d_1)
 \lhook\joinrel\longrightarrow(\cM^1(L),d_1)
\tag{50.29a}
\]
is isometric with dense image, and \((\cM^1(L),d_1)\) is complete.  This
uses, in particular, the comparison of the Darvas metric with the BJ22
quasi-metrics and the completeness of the measure-side strong quasi-metric
\cite[Theorems~10.12 and~10.14]{BJ22}; these estimates are not reproved
here.  Third,
\[
\label{eq:ch50-canonical-approximants-dense}
 \chi_d\in\mathcal T_{\mathbf R},
 \qquad d_1(\chi_d,\chi)\longrightarrow0
 \quad(\chi\in\mathcal N_{\mathbf R})
\tag{50.29b}
\]
for the canonical finite-type approximants; this is
\cite[Theorem~3.18]{BJKS1}.  The isometry with dense image
\[
\label{eq:ch50-test-ma-dense}
 \MA:(\mathcal T_{\mathbf R}/\mathbf R,\overline d_1)
 \longrightarrow(\cM^1(L),d_1)
\tag{50.29c}
\]
is \cite[Lemma~7.1]{BJKS1}: it is obtained by composing the quotient
isometry of \Cref{thm:stability-fs-quotient} with
\eqref{eq:ch50-e1-ma-dense}.

We now give the formal extension argument in full.  Let
\(q:\mathcal N_{\mathbf R}\to\mathcal N_{\mathbf R}/\mathbf R\) be the quotient map,
with the Hausdorff convention of \eqref{eq:ch50-quotient-d1}.  For
\(q(\chi)\), choose a net \((\tau_a)\) in \(\mathcal T_{\mathbf R}\) such that
\[
\label{eq:ch50-norm-approximating-net}
 \overline d_1(q(\tau_a),q(\chi))\longrightarrow0;
\tag{50.29d}
\]
the canonical approximants in
\eqref{eq:ch50-canonical-approximants-dense} provide such a net.  Since
\eqref{eq:ch50-test-ma-dense} is an isometry, \((\MA(\tau_a))\) is Cauchy
in the complete metric space \(\cM^1(L)\).  Define
\[
\label{eq:ch50-ma-completion-definition}
 \MA(q(\chi)):=\lim_a\MA(\tau_a).
\tag{50.29e}
\]
If \((\tau'_b)\) is another approximating net, then on the product directed
set
\begin{align*}
 d_1\bigl(\MA(\tau_a),\MA(\tau'_b)\bigr)
 &=\overline d_1\bigl(q(\tau_a),q(\tau'_b)\bigr)\\
 &\leq\overline d_1(q(\tau_a),q(\chi))
     +\overline d_1(q(\chi),q(\tau'_b))\longrightarrow0.
\end{align*}
Thus \eqref{eq:ch50-ma-completion-definition} is independent of every
choice, and it is constant on the zero-distance classes already removed by
Hausdorffization.

For \(\chi,\chi'\in\mathcal N_{\mathbf R}\), choose approximating nets
\((\tau_a)\), \((\tau'_b)\).  Continuity of the distance and the triangle
inequality give
\begin{align*}
 d_1\bigl(\MA(q(\chi)),\MA(q(\chi'))\bigr)
 &=\lim_{(a,b)}d_1\bigl(\MA(\tau_a),\MA(\tau'_b)\bigr)\\
 &=\lim_{(a,b)}\overline d_1\bigl(q(\tau_a),q(\tau'_b)\bigr)\\
 &=\overline d_1\bigl(q(\chi),q(\chi')\bigr).
\end{align*}
Hence the extension is an isometry.  It is unique because
\(q(\mathcal T_{\mathbf R})\) is dense by
\eqref{eq:ch50-canonical-approximants-dense}, and its image is dense because
it contains the dense image in \eqref{eq:ch50-test-ma-dense}.  Completeness
of \(\cM^1(L)\) now says exactly that it is the Hausdorff completion of both
\((\mathcal T_{\mathbf R}/\mathbf R,\overline d_1)\) and
\((\mathcal N_{\mathbf R}/\mathbf R,\overline d_1)\).

On the unquotiented potential side,
\Cref{thm:stability-fs-quotient} identifies the Hausdorff quotient of
\((\mathcal T_{\mathbf R},d_1)\) isometrically with the dense subspace
\(\cH_{\mathbf R}(L)\subset\cE^1(L)\).  Therefore \(\cE^1(L)\) is that
completion exactly when it is complete.  The imported completeness
criterion \cite[Theorem~5.5(ii)]{BJKS1} says that this occurs if and only if
the envelope property holds.  This proves all three clauses and both final
completion assertions.
\end{proof}

In our unibranch setting
\Crefrange{chap:bj-envelope-statement}{chap:bj-envelope-consequences}
establishes the envelope property,
so both descriptions are available.  The distinction remains important:
a general filtration is not finitely generated, and its finite-energy
representative need not be PL or a model function.

\begin{remark}[Which half of \Cref{thm:stability-norm-completion} is
book-internal]
\label{rem:ch50-completion-status}
The underlying strong-topology statements in clause~(3) are proved earlier
in this book: completeness of \(\cM^1(L)\) is
\Cref{thm:bj-strong-measures}, completeness of \(\cE^1(L)\) under the
envelope property is \Cref{thm:ch48-completeness}, and the failure without
it is \Cref{cor:ch45-failure-nonunibranch}.  Clause (1) is the density
statement \Cref{thm:ch34-pl-density} combined with the strong
approximation of \Cref{prop:bj-strong-monotone}.  What is imported from
\cite{BJKS1} is the exact compatible metric package listed at the start of
the proof, including \eqref{eq:ch50-e1-ma-dense}, measure-side metric
completeness, and the canonical approximation
\eqref{eq:ch50-canonical-approximants-dense}.  Once these are granted, the
passage from \eqref{eq:ch50-test-ma-dense} to clause~(2) is the explicit
completion argument \eqref{eq:ch50-norm-approximating-net}--
\eqref{eq:ch50-ma-completion-definition}.  It shows that the assignment
\(\chi\mapsto\MA(\FS(\chi))\), defined on finite-type norms by
\Cref{prop:ch46-pl-measure}, extends to an
\emph{isometry} for \(d_1\) on the Hausdorff quotient of all of
\(\mathcal N_{\mathbf R}\).  The essential point is that a general norm
has a well-defined energy, obtained in \cite[Theorem~5.1]{BJKS1} as an
upper extension; the corresponding potential \(\FS(\chi)\) may fail to lie
in \(\cE^1(\omega)\), and its measure is defined by the strong limit
\eqref{eq:ch50-ma-completion-definition}, not by applying the
Monge--Amp\`ere operator to that raw pointwise function.  When the envelope
property holds, \Cref{cor:bj-cy-unibranch} subsequently represents the
limiting measure by a unique finite-energy potential modulo constants.
This asymmetry is exactly what
\Cref{thm:stability-norm-completion} is warning about, and it is the
reason \eqref{eq:stability-M-all-norms} is stated for measures.
\end{remark}

\section{Duistermaat--Heckman measures and energy functionals}
\label{sec:ch50-dh}

\subsection{The spectral measure of a norm}
\label{sec:ch50-spectral-measure}

Let \(\lambda_{m,1}\geq\cdots\geq\lambda_{m,N_m}\) be the jumping values of
\(\chi|_{R_m}\).  The empirical measures
\[
             \nu_m=\frac1{N_m}\sum_{j=1}^{N_m}
                    \delta_{\lambda_{m,j}/m}
\tag{50.30}\label{eq:stability-empirical-dh}
\]
converge weakly to the Duistermaat--Heckman probability measure
\(\DHm_\chi\).  Its support is compact by linear boundedness.

The convergence, the compactness of the support, the description of
\(\DHm_\chi\) as the push-forward of Lebesgue measure on an Okounkov body
by a concave transform, and the identification of the top of the support
with the Fekete limit \eqref{eq:ch50-lambda-max}, are
\Cref{thm:dh-limit-filtration}, \Cref{thm:ch33-dh-concave-transform} and
\Cref{cor:ch33-dh-structure}.  We shall use all four, and in particular the
Brunn--Minkowski concavity
\Cref{cor:ch33-dh-structure}\,(2), which is the engine of
\Cref{thm:ch50-J-L1}.

\begin{proposition}
\label{prop:ch50-sup-lambdamax}
Let \(\chi\in\mathcal N_{\mathbf R}\) and \(\phi=\FS(\chi)\).  Then
\[
 \sup_{X^{\An}}\phi=\phi(v_{\triv})
 =\lambda_{\max}(\chi)
 =\max\operatorname{supp}\DHm_\chi .
\tag{50.31}\label{eq:ch50-sup-lambdamax}
\]
\end{proposition}

\begin{proof}
Every \(v\in X^{\An}\) is a semivaluation with a center, so
\(v(s)\ge0\) for every regular section \(s\); and
\(v_{\triv}(s)=0\) for every \(s\ne0\).  Hence, from
\eqref{eq:ch50-fs-operator},
\[
 \phi(v)=\sup_{m,s}\frac{\chi_m(s)-v(s)}{m}
 \ \le\ \sup_{m,s}\frac{\chi_m(s)}{m}
 =\phi(v_{\triv}),
\]
and the right-hand side is \(\lambda_{\max}(\chi)\) by
\eqref{eq:ch50-lambda-max}.  The identification
\(\lambda_{\max}=\max\operatorname{supp}\DHm_\chi\) is
\Cref{cor:ch33-dh-structure}\,(3).  (For \(\phi\in\PSH(\omega)\) the first
equality is also \Cref{lem:ch47-three-normalizations}\,\cref{itm:ch47-norm:max};
the computation above is independent of it and does not use pshness.)
\end{proof}

\begin{theorem}
\label{thm:ch50-energy-dh}
Let \(\chi\in\mathcal T_{\mathbf R}\) be of finite type, with
\(\phi=\FS(\chi)\in\mathcal H_{\mathbf R}(L)\).  Then
\[
 E(\phi)=\int_{\mathbf R}\lambda\,\rd \DHm_\chi(\lambda),
\qquad
 J(\phi)=\sup\operatorname{supp}\DHm_\chi
                  -\int_{\mathbf R}\lambda\,\rd \DHm_\chi(\lambda).
\tag{50.32}\label{eq:stability-dh-energy-J}
\]
Equivalently \(E_\omega(\phi)=\operatorname{vol}(\chi)\) and
\(J(\phi)=\lambda_{\max}(\chi)-\operatorname{vol}(\chi)\), in the notation
of \eqref{eq:ch50-lambda-max} and \eqref{eq:stability-volume-norm}.
\end{theorem}

\begin{proof}
Assume first that \(\chi\in\mathcal T_{\mathbf Z}\) is integral and that
\(\phi\) is the PL function of an integrally closed test configuration
\((\cX,\cL)\) with \(\cL\) relatively ample.  Then the first identity is
\Cref{prop:ch46-energy-dh}.  A general \(\chi\in\mathcal T_{\mathbf R}\) is
obtained from an integral one by translation \(\chi\mapsto\chi+c\) with
\(c\in\mathbf R\), by base change \(\chi\mapsto\chi^{[d]}\), and by
replacing \(L\) with a positive multiple; under these operations
\(E_\omega(\phi)\) transforms by \(E\mapsto E+c\), \(E\mapsto dE\) and
\(E\mapsto rE\) respectively, by
\eqref{eq:bj-energy-basic-identities} and
\Cref{thm:ch46-energy-homogeneous}, and the first moment of \(\DHm\)
transforms in the same way by
\Cref{prop:ch33-dh-transformations}.  Both sides are therefore determined
by the integral case, and they agree.  Finally, if
\(\cL\) is only semiample, dominate \(\phi\) by relatively ample data and
use continuity of both sides along the resulting decreasing net, by
\Cref{prop:ch46-energy-usc}\,\cref{itm:ch46-Eusc:decreasing} and
\Cref{prop:ch33-dh-continuity}.

The second identity in \eqref{eq:stability-dh-energy-J} is now
\eqref{eq:ch47-J-functional} combined with
\Cref{prop:ch50-sup-lambdamax}.
\end{proof}

\begin{remark}[Beyond finite type]
\label{rem:ch50-energy-beyond-finite-type}
For a general \(\chi\in\mathcal N_{\mathbf R}\) the identity
\(E_\omega(\FS(\chi))=\operatorname{vol}(\chi)\) can fail, and the failure
is instructive.  The right-hand side is always defined, by
\Cref{thm:dh-limit-filtration}.  The left-hand side is defined whenever
\(\FS(\chi)\in\cE^1(\omega)\), which is automatic here because
\(\FS(\chi)\) is bounded --- indeed
\(\lambda_{\min}\le\FS(\chi)\le\lambda_{\max}\) by
\eqref{eq:ch50-fs-operator} and linear boundedness --- so both sides are
finite.  What can fail is the equality, because \(\FS(\chi)\) only sees
\(\chi\) through the closure \(\chi^{\FS(\chi)}\) of
\Cref{thm:ch50-galois}, and the Duistermaat--Heckman measure of a norm is
not in general that of its closure.  The correct general statement is
\(E_\omega(\FS(\chi))\ge\operatorname{vol}(\chi)\), with equality if and
only if \(d_1(\chi,\chi^{\FS(\chi)})=0\); this is
\cite[Theorem~5.1]{BJKS1} and the surrounding discussion there.  Every use below is either
for finite-type norms or for the measure-theoretic functional
\eqref{eq:stability-M-all-norms}, which is defined through
\(\MA(\chi)\) and not through \(\DHm_\chi\).  \Cref{ex:ch50-harmonic}
computes both sides for a non-finitely generated filtration and finds them
equal, as it must, that filtration being homogeneous.
\end{remark}

\subsection{The centred spectral norms}
\label{sec:ch50-spectral-norms}

The centered spectral norms are
\[
 \|\chi\|_p
 :=\left(\int_{\mathbf R}|\lambda-\bar\lambda|^p
                    \,\rd \DHm_\chi(\lambda)\right)^{1/p},
\qquad
 \bar\lambda=\int_{\mathbf R}\lambda\,\rd \DHm_\chi(\lambda).
\tag{50.33}\label{eq:stability-spectral-lp}
\]
They are invariant under translation and homogeneous under base change.
For positive model metrics,
\[
        c_nJ(\phi)\leq\|\chi\|_1\leq2J(\phi)
\tag{50.34}\label{eq:stability-J-L1-comparison}
\]
with \(c_n=2n^n/(n+1)^{n+1}\)
\cite[Theorem~7.9]{BHJ17}.  In particular, any of these quantities
vanishes exactly on a translate of the trivial metric.

We give a complete proof of \eqref{eq:stability-J-L1-comparison}, with
both constants, and identify the extremal measure.  Only two inputs are
used: the Brunn--Minkowski concavity of
\Cref{cor:ch33-dh-structure}\,(2) and
\Cref{prop:ch50-sup-lambdamax}.  Notice that
\eqref{eq:stability-J-L1-comparison} is really a statement about a class of
probability measures on \(\mathbf R\), and we isolate it as such.

\begin{theorem}
\label{thm:ch50-J-L1}
\index{Duistermaat--Heckman measure!centred \(L^1\)-norm}
Let \(\mu\) be a compactly supported probability measure on \(\mathbf R\)
such that
\[
 F(\tau):=\mu\bigl([\tau,+\infty)\bigr)
\]
satisfies: \(F^{1/n}\) is concave on the interval
\(\{F>0\}=[\lambda_{\min},a]\), where
\(a:=\max\operatorname{supp}\mu\).  Write
\(\bar\lambda:=\int\lambda\,\rd\mu\), \(A:=a-\bar\lambda\) and
\(\|\mu\|_1:=\int|\lambda-\bar\lambda|\,\rd\mu\).  Then
\[
 \frac{2n^n}{(n+1)^{n+1}}\,A\ \le\ \|\mu\|_1\ \le\ 2A .
\tag{50.35}\label{eq:ch50-J-L1-abstract}
\]
Both constants are attained, the right one only in the limit and the left
one exactly for
\(\rd\mu=n(a-\lambda)^{n-1}(a-\lambda_{\min})^{-n}\,\rd\lambda\) on
\([\lambda_{\min},a]\).

Consequently, for every \(\chi\in\mathcal N_{\mathbf R}\) with associated
\(\phi=\FS(\chi)\) of finite type,
\eqref{eq:stability-J-L1-comparison} holds, with \(A=J(\phi)\) and
\(\|\mu\|_1=\|\chi\|_1\).
\end{theorem}

\begin{proof}
We may assume \(A>0\); otherwise \(\mu=\delta_a\) and both sides vanish.
Translating, assume \(a=0\), so that \(\bar\lambda=-A<0\) and
\(\lambda_{\min}=-\ell\) with \(\ell\ge A\).

\emph{Upper bound.}  Since \(\int(\lambda-\bar\lambda)\rd\mu=0\), we have
\(\int(\lambda-\bar\lambda)_+\rd\mu=\int(\lambda-\bar\lambda)_-\rd\mu\) and
hence \(\|\mu\|_1=2\int(\lambda-\bar\lambda)_+\rd\mu\).  On the support,
\(\lambda\le a\), so \((\lambda-\bar\lambda)_+\le a-\bar\lambda=A\) and
\(\|\mu\|_1\le2A\).  Equality would force \(\mu\) to be carried by
\(\{a\}\cup\{\lambda:\lambda\le\bar\lambda\}\) with all the upper mass at
\(a\); a two-atom measure \((1-\varepsilon)\delta_{-\ell}
+\varepsilon\delta_0\) has \(\|\mu\|_1/A=2\varepsilon\to2\) as
\(\varepsilon\to1^-\),
but no measure with \(F^{1/n}\) concave attains it, concavity forcing
\(F\) to be continuous on the interior of its support.

\emph{Lower bound, step 1: the mass above the mean.}  Put
\(v(s):=F(-s)^{1/n}\) for \(s\in[0,\ell]\).  Since \(F\) is decreasing and
\(F^{1/n}\) is concave on \([-\ell,0]\), the function \(v\) is increasing
and concave on \([0,\ell]\), with \(v(\ell)=F(-\ell)^{1/n}=1\) and
\(v(0)=F(0)^{1/n}=\mu(\{a\})^{1/n}\ge0\).  We claim
\[
 v(A)\ \ge\ \frac{n}{n+1}.
\tag{50.36}\label{eq:ch50-mass-above-mean}
\]
Write \(\gamma:=v(A)\); we may assume \(\gamma>0\), for otherwise
\(F\equiv0\) on \([-A,0]\) and then \(\mu\) is carried by
\([-\ell,-A]\), contradicting \(\int\lambda\,\rd\mu=-A\) unless
\(\mu=\delta_{-A}\), a case in which \(\ell=A\) and
\(v\equiv1\).  We may also assume \(\gamma<1\): if \(v(A_1)=v(A)\) for some \(A_1<A\) then
\(v\), being concave and nondecreasing, is constant on \([A_1,\ell]\), so
\(\gamma=v(\ell)=1\) and \eqref{eq:ch50-mass-above-mean} holds trivially.
Hence \(v\) is strictly increasing on \([0,A]\) and \(\gamma\in(v_0,1)\).
Put \(v_0:=v(0)\ge0\) and let \(\psi\) be the inverse of \(v\),
\[
 \psi(w):=\bigl|\{s\in[0,\ell]:v(s)<w\}\bigr|
 \qquad(v_0\le w\le1),
\]
which is convex and nondecreasing, being the inverse of an increasing
concave function, with \(\psi(v_0)=0\) and \(\psi(\gamma)=A\).  The
layer-cake formula rewrites the mass identity as
\[
 A=\int(a-\lambda)\,\rd\mu
  =\int_{-\ell}^{0}\bigl(1-F(\tau)\bigr)\rd\tau
  =\int_0^{\ell}\bigl(1-v(s)^n\bigr)\rd s
  =\int_0^{\ell}\!\!\int_{v(s)}^{1}nw^{n-1}\rd w\,\rd s
  =n\int_{v_0}^{1}w^{n-1}\psi(w)\,\rd w .
\]
Write \(m:=n\int_{v_0}^1w^{n-1}\rd w=1-v_0^n\) and
\(q:=n\int_{v_0}^1w^{n}\rd w=\tfrac n{n+1}\bigl(1-v_0^{n+1}\bigr)\), and let
\(p\) be a subgradient of \(\psi\) at \(\gamma\); \(p>0\) because
\(\psi(v_0)=0<A=\psi(\gamma)\).  Convexity supplies two inequalities,
\[
 \psi(w)\ \ge\ \psi(\gamma)+p\,(w-\gamma)\ \ (v_0\le w\le1),
 \qquad
 0=\psi(v_0)\ \ge\ \psi(\gamma)+p\,(v_0-\gamma).
\tag{50.37}\label{eq:ch50-concave-upper}
\]
Integrating the first against \(nw^{n-1}\rd w\) over \([v_0,1]\) and using
\(\psi(\gamma)=A\) gives \(A\ge Am+p(q-\gamma m)\), that is
\(Av_0^{n}\ge p\,(q-\gamma m)\); the second gives \(A\le p\,(\gamma-v_0)\).
If \(q-\gamma m\le0\), then \(\gamma\ge q/m\ge\frac n{n+1}\), because
\(1-v_0^{n+1}\ge1-v_0^{n}\).  Otherwise
\(p(q-\gamma m)\le Av_0^{n}\le p(\gamma-v_0)v_0^{n}\), whence
\(q\le\gamma m+\gamma v_0^{n}-v_0^{n+1}=\gamma-v_0^{n+1}\) and
\[
 \gamma\ \ge\ q+v_0^{n+1}
 =\frac n{n+1}+\frac{v_0^{n+1}}{n+1}\ \ge\ \frac n{n+1},
\]
which is \eqref{eq:ch50-mass-above-mean}.  The argument has to be run this
way round.  The tempting bound \(v(s)\le\gamma s/A\) is available only for
\(s\ge A\); on \([0,A]\) concavity together with \(v(0)\ge0\) forces the
\emph{opposite} inequality \(v(s)\ge\gamma s/A\), so it cannot be
integrated over \([0,A/\gamma]\).  Both estimates above are attained:
\(v(s)=s\) gives \(v_0=0\) and \(\gamma=\frac n{n+1}\), while
\(v(s)=\frac12(1+s)\) on \([0,1]\) with \(n=1\) gives \(v_0=\frac12\),
\(A=\frac14\), \(q=\frac38\) and \(\gamma=\frac58=q+v_0^{2}\).

\emph{Lower bound, step 2: concavity above the mean.}  For
\(\tau\in[\bar\lambda,a]\), concavity of \(F^{1/n}\) and
\(F(a)\ge0\) give
\[
 F(\tau)^{1/n}\ \ge\ F(\bar\lambda)^{1/n}\,\frac{a-\tau}{a-\bar\lambda},
\]
whence
\[
 \int_{\bar\lambda}^{a}F(\tau)\,\rd\tau
 \ \ge\ F(\bar\lambda)\int_{\bar\lambda}^{a}
      \Bigl(\frac{a-\tau}{A}\Bigr)^{n}\rd\tau
 =\frac{F(\bar\lambda)\,A}{n+1}.
\]
Since \(\int(\lambda-\bar\lambda)_+\rd\mu
=\int_{\bar\lambda}^{a}\mu((\tau,\infty))\,\rd\tau
=\int_{\bar\lambda}^{a}F(\tau)\,\rd\tau\), the two sets of \(\tau\) where
\(F\) and \(\mu((\tau,\infty))\) differ being countable, we obtain, using
\(F(\bar\lambda)=v(A)^n\ge(n/(n+1))^n\) from
\eqref{eq:ch50-mass-above-mean},
\[
 \|\mu\|_1=2\int_{\bar\lambda}^{a}F(\tau)\,\rd\tau
 \ \ge\ \frac{2}{n+1}\Bigl(\frac{n}{n+1}\Bigr)^{n}A
 =\frac{2n^n}{(n+1)^{n+1}}\,A .
\]

\emph{Sharpness.}  Take \(a=0\), \(\ell=1\) and
\(F(\tau)=(-\tau)^n\) on \([-1,0]\), that is
\(\rd\mu=n(-\lambda)^{n-1}\rd\lambda\).  Then
\(F^{1/n}\) is linear, hence concave; \(\bar\lambda=-\int_0^1s\cdot
ns^{n-1}\rd s=-\frac n{n+1}\), so \(A=\frac n{n+1}\); and
\[
 \|\mu\|_1=2\int_{-A}^0(-\tau)^n\rd\tau=\frac{2A^{n+1}}{n+1}
 =\frac{2}{n+1}\Bigl(\frac n{n+1}\Bigr)^{n+1},
\]
so that \(\|\mu\|_1/A=\frac2{n+1}\bigl(\frac n{n+1}\bigr)^n=c_n\).  Every
inequality used above is an equality for this \(\mu\).

\emph{Application.}  For \(\mu=\DHm_\chi\), the hypothesis on \(F^{1/n}\)
is \Cref{cor:ch33-dh-structure}\,(2), \(a=\lambda_{\max}(\chi)\) is
\Cref{cor:ch33-dh-structure}\,(3), \(\bar\lambda=E_\omega(\phi)\) is
\Cref{thm:ch50-energy-dh}, and therefore \(A=\lambda_{\max}-\bar\lambda
=J(\phi)\) by \eqref{eq:stability-dh-energy-J}.
\end{proof}

\begin{remark}[The extremal measure is a cone]
\label{rem:ch50-extremal-cone}
The extremal measure of \Cref{thm:ch50-J-L1} is the Duistermaat--Heckman
measure of a \emph{cone}: by \eqref{eq:ch33-dh-is-pushforward} it is the
push-forward of Lebesgue measure on the Okounkov body \(\Delta\) by an
affine concave transform whose superlevel sets are homothets of \(\Delta\)
shrinking to a point.  The deformation to the normal cone of a point on
\(\mathbf P^n\), polarized at the boundary of the ample cone, realizes it;
for \(n=1\) this is \Cref{ex:ch33-dnc-p1} with \(c=1\), whose measure
\(\mathbf 1_{[-1,0]}\rd\lambda\) is exactly the case \(n=1\) of
\eqref{eq:ch50-J-L1-abstract}, with
\(\|\mu\|_1=\tfrac14\) and \(A=\tfrac12\), ratio \(c_1=\tfrac12\).  This is
the same extremal configuration that saturates the two constants
\(\frac1{n+1}\) and \(\frac n{n+1}\) of
\eqref{eq:bj-dirac-energy-T}, as \Cref{cor:ch49-minkowski} explains: all
three sharp constants of the theory are the constants in Minkowski's
centroid inequality, and all are attained on cones.
\end{remark}

\begin{corollary}
\label{cor:ch50-vanishing}
For \(\chi\in\mathcal T_{\mathbf R}\) with \(\phi=\FS(\chi)\), the
following are equivalent:
\(\|\chi\|_1=0\); \(\|\chi\|_p=0\) for some, equivalently every,
\(p\in[1,\infty]\); \(J(\phi)=0\); \(\phi\) is constant; \(\DHm_\chi\) is a
Dirac mass; \(\chi\) is a translate of the trivial norm up to
\(d_1\)-distance zero.
\end{corollary}

\begin{proof}
By \eqref{eq:stability-J-L1-comparison} and
\eqref{eq:ch50-dp-interpolation}, \(\|\chi\|_1=0\) is equivalent to
\(J(\phi)=0\), by the two-sided comparison
\eqref{eq:stability-J-L1-comparison}, and to \(\|\chi\|_p=0\) for every
\(p\in[1,\infty]\): the centred \(L^1\) norm of a probability measure
vanishes exactly when the measure is the Dirac mass at its own mean, and
then every centred \(L^p\) norm vanishes.  (One may \emph{not} argue by
\(\|\chi\|_\infty\le2J\); that inequality is false.  On
\Cref{ex:ch50-p1-dnc} with \(c=\tfrac15\) and \(L=\cO(2)\) one has
\(\bar\lambda=-c^2/4=-0.01\), \(\lambda_{\min}=-0.2\), \(J=0.01\) and
\(\|\chi\|_\infty=0.19>0.02=2J\).  Only \eqref{eq:stability-J-L1-comparison},
for \(p=1\), is available.)  That \(J(\phi)=0\) forces
\(\phi\) constant is \Cref{cor:ch47-J-vanishes}.  A constant \(\phi\) has
\(E=\sup\phi\), so \(\DHm_\chi\) has zero variance and is a Dirac mass by
\eqref{eq:stability-spectral-lp}; and then
\(\operatorname{vol}(\chi)=\operatorname{vol}(\chi\wedge\chi_{\triv}+c)\)
gives \(d_1(\chi,\chi_{\triv}+c)=0\) by \eqref{eq:ch50-d1-volume}.
\end{proof}

\subsection{The minimum norm}
\label{sec:ch50-minimum-norm}

For any \(\chi\in\mathcal N_{\mathbf R}\), define its minimum norm by
\[
                   \|\chi\|_{\min}
                   :=E^\vee(\MA(\chi)).
\tag{50.38}\label{eq:stability-minimum-norm}
\]
It is \(d_1\)-continuous, translation invariant, homogeneous, and
comparable to \(d_1(\chi,\chi_{\triv})\) after optimizing by translations.
For an ample test configuration it agrees, up to the fixed factor \(V\),
with Dervan's minimum norm; for a finite-type norm,
\[
 \|\chi\|_{\min}
   =E(\phi)-\int\phi\,\MA(\phi)=I(\phi)-J(\phi).
\tag{50.39}\label{eq:stability-minimum-IJ}
\]
These are \cite[Definition~7.19 and Proposition~7.20]{BJKS1}.

The identity \eqref{eq:stability-minimum-IJ} is book-internal, and we prove
it, together with the elementary properties.

\begin{proposition}
\label{prop:ch50-minimum-norm}
Let \(\phi\in\cE^1(\omega)\) and \(\mu:=\MA_\omega(\phi)\).  Write
\(I(\phi):=I(\phi,0)\) with \(I\) as in \eqref{eq:bj-I-J-potentials}.
Then
\[
 E^\vee_\omega(\mu)=E_\omega(\phi)-\int\phi\,\rd\mu
 =J_\phi(0)=I(\phi)-J(\phi)\ \in[0,+\infty),
\tag{50.40}\label{eq:ch50-min-norm-identities}
\]
and:
\begin{enumerate}
\item\label{itm:ch50-mn:transl}
      \(E^\vee_\omega(\MA_\omega(\phi+c))=E^\vee_\omega(\MA_\omega(\phi))\)
      for \(c\in\mathbf R\);
\item\label{itm:ch50-mn:zero}
      \(E^\vee_\omega(\mu)=0\) if and only if \(\phi\) is constant, if and
      only if \(\mu=\delta_{v_{\triv}}\);
\item\label{itm:ch50-mn:comparison}
      \(\frac1{n+1}I(\phi)\le E^\vee_\omega(\MA_\omega(\phi))\le
      \frac n{n+1}I(\phi)\) and
      \(\frac1{n+1}I(\phi)\le J(\phi)\le\frac n{n+1}I(\phi)\); in
      particular
      \[
       \frac1n\,J(\phi)\ \le\ \|\chi\|_{\min}\ \le\ n\,J(\phi)
      \tag{50.41}\label{eq:ch50-min-vs-J}
      \]
      whenever \(\phi=\FS(\chi)\);
\item\label{itm:ch50-mn:scaling}
      \(E^\vee_\omega(\MA_\omega(t\cdot\phi))=t\,E^\vee_\omega
      (\MA_\omega(\phi))\) for the scaling action of \(t>0\).
\end{enumerate}
\end{proposition}

\begin{proof}
The first two equalities in \eqref{eq:ch50-min-norm-identities} are the
Fenchel equality \eqref{eq:bj-dual-energy-ma} of
\Cref{prop:bj-fenchel-equalities}, together with the definition
\eqref{eq:bj-I-J-potentials} of \(J_\phi\) and the normalizations
\(E_\omega(0)=0\), \(\MA_\omega(0)=\delta_{v_{\triv}}\).  The third follows
from \(I(\phi)=I(\phi,0)=J_0(\phi)+J_\phi(0)\), again
\eqref{eq:bj-I-J-potentials}, and
\(J_0(\phi)=J(\phi)\), which is
\Cref{lem:ch47-three-normalizations}\,\cref{itm:ch47-norm:J}.
Nonnegativity and finiteness are
\Cref{lem:ch47-IJ-elementary}\,\cref{itm:ch47-IJ:nonneg}.

\cref{itm:ch50-mn:transl}  \(\MA_\omega(\phi+c)=\MA_\omega(\phi)\), so the
left side does not move; alternatively, \(E(\phi+c)=E(\phi)+c\) and
\(\int(\phi+c)\rd\mu=\int\phi\,\rd\mu+c\) since \(\mu\) is a probability
measure.

\cref{itm:ch50-mn:zero}  If \(\phi\) is constant then
\(\mu=\delta_{v_{\triv}}\) and both sides of
\eqref{eq:ch50-min-norm-identities} vanish.  Conversely
\(E^\vee_\omega(\mu)=0\) gives \(I(\phi)=J(\phi)\); by
\cref{itm:ch50-mn:comparison} this forces \(I(\phi)=0\), whence \(\phi\) is
constant by \Cref{thm:ch47-uniqueness}\,\cref{itm:ch47-uniq:I}.

\cref{itm:ch50-mn:comparison}  Both chains are
\eqref{eq:bj-I-J-comparison} of \Cref{prop:bj-energy-I-identities}, applied
to the pair \((\phi,0)\) in the two possible orders.  Combining them,
\(\|\chi\|_{\min}=I-J\ge I-\frac n{n+1}I=\frac{I}{n+1}\ge\frac1n\cdot
\frac{n}{n+1}I\ge\frac1nJ\), and symmetrically.

\cref{itm:ch50-mn:scaling}  This is
\Cref{prop:ch47-scaling}\,\cref{itm:ch47-scal:JI} together with
\Cref{prop:ch47-dual-scaling}.
\end{proof}

\begin{remark}[Which functional measures what]
\label{rem:ch50-which-functional}
Four translation-invariant nonnegative functionals have now appeared:
\(J\), \(I\), \(\|\cdot\|_{\min}=I-J\) and \(\|\cdot\|_1\).  By
\eqref{eq:ch50-min-vs-J} and \eqref{eq:stability-J-L1-comparison} they are
all comparable, with constants depending only on \(n\):
\[
 \frac{c_n}{n}\,\|\chi\|_{\min}\ \le\ c_nJ\ \le\ \|\chi\|_1
 \ \le\ 2J\ \le\ 2n\,\|\chi\|_{\min},
 \qquad
 \frac{1}{n+1}I\le J\le\frac{n}{n+1}I .
\]
Hence any one of them may be used to define uniform K-stability, and the
resulting notions agree up to the value of \(\varepsilon\); this is why
\eqref{eq:stability-uniform-K} and \eqref{eq:stability-M-coercive-model}
are equivalent.  They are \emph{not} interchangeable when a sharp constant
is at stake, and they behave differently under limits: \(J\) and \(I\) are
defined on all of \(\cE^1(\omega)\), \(\|\cdot\|_{\min}\) is defined on all
of \(\cM^1(\omega)\) --- and therefore, by
\eqref{eq:stability-norm-ma-isometry}, on all of
\(\mathcal N_{\mathbf R}\) --- while \(\|\cdot\|_1\) requires a
Duistermaat--Heckman measure and hence a norm.  In this chapter
\(\|\cdot\|_{\min}\) is the functional of record, because it is the one
that extends to the completion.
\end{remark}

\subsection{The product configuration on the projective line}
\label{sec:ch50-p1-product}

\begin{example}[The product action on \(\mathbf P^1\)]
\label{ex:stability-p1-product}
Let \(L=\cO_{\mathbf P^1}(1)\) and give \(x^jy^{m-j}\) weight \(j\).
Then \(\DHm_\chi\) is Lebesgue probability measure on \([0,1]\).  Thus
\[
 E(\phi)=\frac12,\qquad J(\phi)=\frac12,\qquad
 \|\chi\|_1=\int_0^1\left|\lambda-\frac12\right|\,\rd\lambda=\frac14.
\tag{50.42}\label{eq:stability-p1-numbers}
\]
The total space is a product, but the induced metric is nontrivial because
the action on the general fibre is nontrivial.
\end{example}

We carry the example to the end, because every functional of this chapter
can be evaluated on it by hand and because it is the basic obstruction to
uniform K-stability.

\begin{proposition}
\label{prop:ch50-p1-product}
In the situation of \Cref{ex:stability-p1-product}, write
\(p_0=[0:1]\) and \(p_\infty=[1:0]\), and let
\(\phi=\FS(\chi)\).  Then:
\begin{enumerate}
\item\label{itm:ch50-p1p:fs}
      \(\phi(t\,\ord_{p_0})=\max\{1-t,0\}\), \(\phi\equiv1\) on the ray
      through \(p_\infty\) and at every other point of
      \((\mathbf P^1)^{\An}\), and \(\sup\phi=\phi(v_{\triv})=1\);
\item\label{itm:ch50-p1p:ma}
      \(\MA_\omega(\phi)=\delta_{\ord_{p_0}}\), and \(\phi\) is, up to an
      additive constant, the Green function of \(\ord_{p_0}\);
\item\label{itm:ch50-p1p:numbers}
      \(E(\phi)=\tfrac12\), \(J(\phi)=\tfrac12\),
      \(\|\chi\|_1=\tfrac14\),
      \(\|\chi\|_{\min}=\tfrac12\), \(I(\phi)=1\);
\item\label{itm:ch50-p1p:mabuchi}
      \(H^{\NA}(\phi)=1\), \(R^{\NA}(\phi)=-2\), \(\overline S=2\) and
      \(M^{\NA}(\phi)=\DF(\cX,\cL)=0\).
\end{enumerate}
\end{proposition}

\begin{proof}
\cref{itm:ch50-p1p:fs}  The filtration is monomial:
\(\chi_m(\sum_ja_jx^jy^{m-j})=\min\{j:a_j\ne0\}\).  In the chart
\(y=1\) with coordinate \(z=x/y\), a section \(s=\sum_ja_jz^j\) has
\(\ord_{p_0}(s)=\min\{j:a_j\ne0\}=\chi_m(s)\).  Hence for \(v=t\ord_{p_0}\),
\eqref{eq:ch50-fs-operator} gives
\(\phi(v)=\sup_{m,s}m^{-1}(1-t)\chi_m(s)=\max\{1-t,0\}\), the supremum of
\(\chi_m(s)/m\) over \(s\ne0\) being \(1\), attained at \(s=x^m\), and its
infimum being \(0\), attained at \(s=y^m\).  In the chart \(x=1\) with
coordinate \(w=y/x\), the section \(s\) becomes \(\sum_ja_jw^{m-j}\), so
\(\ord_{p_\infty}(s)=m-\max\{j:a_j\ne0\}\); taking \(s=x^m\) gives
\(\chi_m(s)=m\), \(\ord_{p_\infty}(s)=0\) and hence \(\phi\ge1\) on that
ray, while \(\phi\le\sup\phi=1\) by \Cref{prop:ch50-sup-lambdamax}.  At a
point of \((\mathbf P^1)^{\An}\) lying over a closed point \(q\ne
p_0,p_\infty\) the same section \(x^m\) has value \(0\), so again
\(\phi=1\).

\cref{itm:ch50-p1p:ma}  On a curve, and with the probability normalization,
the Monge--Amp\`ere measure of a psh function is computed ray by ray:
\(\MA^{\mathrm u}_\omega(\phi)\) restricted to the open ray through a
closed point \(q\) is the distributional second derivative of
\(t\mapsto\phi(t\ord_q)\), and the residual mass sits at \(v_{\triv}\)
(\Cref{ex:ch46-curve-laplacian} and \eqref{eq:bj-curve-laplacian}).  Here
\(g(t)=\max\{1-t,0\}\) on the \(p_0\)-ray has \(g''=\delta_1\) and
\(g\equiv1\) elsewhere has \(g''=0\); the total mass \(1\) is accounted for
on the \(p_0\)-ray, so no mass remains at \(v_{\triv}\).  Since \(V=1\),
\(\MA_\omega(\phi)=\delta_{\ord_{p_0}}\).  Uniqueness of solutions
(\Cref{thm:ch47-uniqueness}) identifies \(\phi\) with the Green function up
to a constant.

\cref{itm:ch50-p1p:numbers}  The first three numbers are
\eqref{eq:stability-p1-numbers}, obtained from
\eqref{eq:stability-dh-energy-J} and
\eqref{eq:stability-spectral-lp} with \(\DHm_\chi=\mathbf
1_{[0,1]}\rd\lambda\), the latter being \eqref{eq:ch33-p1-product-dh}.  For
the minimum norm, \(\int\phi\,\MA_\omega(\phi)=\phi(\ord_{p_0})=0\) by
\cref{itm:ch50-p1p:fs}, so \eqref{eq:ch50-min-norm-identities} gives
\(\|\chi\|_{\min}=E(\phi)-0=\tfrac12\), and
\(I=J+\|\chi\|_{\min}=1\).  As a check,
\Cref{prop:ch49-toric-green}\,\cref{itm:ch49-tg:energy} on the toric slice
of \(\mathbf P^1\) with \(P=[0,1]\) gives
\(E^\vee_\omega(\delta_{v_1})=f_P(1)+\langle1,b_P\rangle=0+\tfrac12
=\tfrac12\), and \eqref{eq:bj-dirac-energy-T} gives
\(E^\vee_\omega(\delta_{\ord_{p_0}})=\tfrac12T_\omega(\ord_{p_0})
=\tfrac12\), all three agreeing.

\cref{itm:ch50-p1p:mabuchi}  Entropy: \(A_{\mathbf P^1}(\ord_{p_0})=1\), so
\(H^{\NA}(\phi)=\int A\,\rd\MA_\omega(\phi)=1\) by
\eqref{eq:stability-entropy}.  For \(\overline S\), formula
\eqref{eq:stability-average-scalar} with \(B=0\), \(n=1\), \(V=1\) and
\((K_X\cdot L^{0})=\deg K_{\mathbf P^1}=-2\) gives
\(\overline S=2\).  For the Ricci energy, \(K_X=-2L\) numerically, so by
\eqref{eq:stability-ricci-energy} and
\Cref{lem:ch50-ricci-energy} below,
\(R^{\NA}(\phi)=-2\bigl((n+1)E(\phi)-\int\phi\,\MA_\omega(\phi)\bigr)
=-2\bigl(2\cdot\tfrac12-0\bigr)=-2\).  The Chen--Tian formula
\eqref{eq:stability-chen-tian} then gives
\(M^{\NA}(\phi)=1-2+2\cdot\tfrac12=0\).  The central fibre of the product
configuration is \(X\) itself, reduced, so
\eqref{eq:stability-M-DF-correction} gives \(\DF=M^{\NA}=0\).
\end{proof}

\begin{remark}[The first obstruction]
\label{rem:ch50-p1-obstruction}
\Cref{prop:ch50-p1-product} already shows that
\((\mathbf P^1,\cO(1))\) is not uniformly K-stable in the sense of
\eqref{eq:stability-uniform-K}: the product configuration has
\(\DF=0\) but \(\|\chi\|_{\min}=\tfrac12>0\).  The mechanism is the
positive-dimensional automorphism group, and it is the reason
\Cref{ex:stability-p1-delta} finds \(\delta(\mathbf P^1)=1\) exactly.  The
same computation is the model for every product configuration: by
\Cref{prop:ch49-toric-green}\,\cref{itm:ch49-tg:ma} the Fubini--Study
function of a product configuration attached to a one-parameter subgroup
\(u\) of a torus acting on \((X,L)\) is the Green function of \(v_u\), so
its Monge--Amp\`ere measure is \(\delta_{v_u}\) and all its invariants are
computable from the polytope.  \Cref{ex:ch50-toric-F1} exploits this.
\end{remark}

\section{Log discrepancy, entropy, and Mabuchi energy}
\label{sec:ch50-entropy}

\subsection{Log discrepancy on the analytification}
\label{sec:ch50-log-discrepancy}

Assume now that \(X\) is normal, \(k\) is algebraically closed of
characteristic zero, and \(B\) is a \(\mathbf Q\)-Weil divisor such that
\(K_{X,B}:=K_X+B\) is \(\mathbf Q\)-Cartier.  It need not be effective.
For \(v=c\,\ord_F\in X^{\divis}\), put
\[
             A_{X,B}(v)
             :=c\bigl(1+\ord_F(K_Y-\pi^*(K_X+B))\bigr).
\tag{50.43}\label{eq:stability-log-discrepancy}
\]
This is independent of the model \(Y\).  Its maximal lower
semicontinuous extension to \(X^{\An}\) takes values in
\([0,+\infty]\) when \((X,B)\) is subklt; it is \(+\infty\) outside the
finite-log-discrepancy valuation locus.

We make the extension explicit, because the entropy of a measure is an
integral of it and the finiteness of that integral is a genuine issue.

\begin{definition}[The log discrepancy function on \(X^{\An}\)]
\label{def:ch50-A-extension}
\index{log discrepancy!extension to \(X^{\An}\)}
For an SNC pair \((Y,E)\) over \((X,B)\), let \(\sigma_{Y,E}\subset
X^{\An}\) be its quasimonomial cone and \(r_{Y,E}:X^{\An}\to\sigma_{Y,E}\)
the retraction of \Cref{prop:ch34-retraction}.  On \(\sigma_{Y,E}\) define
\(A_{X,B}\) by the affine formula
\eqref{eq:ch32-log-discrepancy-affine}, with \(K_Y-\pi^*K_X\) replaced by
\(K_Y-\pi^*(K_X+B)\).  Then set
\[
 A_{X,B}(v):=\sup_{(Y,E)}A_{X,B}\bigl(r_{Y,E}(v)\bigr)
 \ \in[-\infty,+\infty],\qquad v\in X^{\An},
\tag{50.44}\label{eq:ch50-A-sup-retractions}
\]
the supremum being over all SNC pairs over \((X,B)\).
\end{definition}

\begin{proposition}
\label{prop:ch50-A-properties}
With \(X\), \(B\) as above:
\begin{enumerate}
\item\label{itm:ch50-A:consistent}
      \eqref{eq:ch50-A-sup-retractions} agrees with
      \eqref{eq:stability-log-discrepancy} on \(X^{\divis}\) and with
      \eqref{eq:ch32-log-discrepancy-affine} on every quasimonomial cone;
      the family in \eqref{eq:ch50-A-sup-retractions} is increasing along
      refinements of \((Y,E)\);
\item\label{itm:ch50-A:lsc}
      \(A_{X,B}\) is lower semicontinuous on \(X^{\An}\);
\item\label{itm:ch50-A:homog}
      \(A_{X,B}(tv)=t\,A_{X,B}(v)\) for \(t>0\), and
      \(A_{X,B}(v_{\triv})=0\);
\item\label{itm:ch50-A:sign}
      if \((X,B)\) is sublc then \(A_{X,B}\ge0\) on \(X^{\An}\), and if
      \((X,B)\) is subklt then \(A_{X,B}>0\) on
      \(X^{\An}\setminus\{v_{\triv}\}\);
\item\label{itm:ch50-A:infinite}
      \(A_{X,B}=+\infty\) at every point of \(X^{\An}\) with proper
      support, and \(\{A_{X,B}<+\infty\}\subset X^{\val}\) is a
      scaling-stable subset containing every quasimonomial point.
\end{enumerate}
\end{proposition}

\begin{proof}
\cref{itm:ch50-A:consistent}  Consistency on a fixed cone is
\Cref{prop:ch32-log-discrepancy-affine}\,\cref{itm:ch32-lda:consistent},
whose proof applies verbatim with \(K_X\) replaced by \(K_X+B\), the only
input being the toroidal computation of the relative log canonical divisor
of a weighted blowup.  If \(v\) is quasimonomial for \((Y,E)\) then
\(r_{Y,E}(v)=v\), so the supremum is at least \(A_{X,B}(v)\); the reverse
inequality, and the monotonicity of the family, is the statement that
\(A_{X,B}(r_{Y',E'}(v))\ge A_{X,B}(r_{Y,E}(v))\) whenever \((Y',E')\)
refines \((Y,E)\), which is proved in \cite{JM12} (published 2012), transported
to the trivially valued analytification through
\Cref{prop:ch34-retraction}.  This is the one imported ingredient of the
present proposition, and it is the classical fact that log discrepancy
increases when the retraction is refined.

\cref{itm:ch50-A:lsc}  Each function \(A_{X,B}\circ r_{Y,E}\) is
continuous: \(r_{Y,E}\) is continuous by \Cref{prop:ch34-retraction}, and
\(A_{X,B}\) is continuous on \(\sigma_{Y,E}\) by
\Cref{prop:ch32-log-discrepancy-affine}\,\cref{itm:ch32-lda:positive}, being
affine with finitely many terms on each face and compatible with the face
maps.  A supremum of continuous functions is lower semicontinuous.

\cref{itm:ch50-A:homog}  Each \(r_{Y,E}\) commutes with the scaling action
(\Cref{prop:ch34-retraction}\,\cref{itm:ch34-retr:pl}) and each
\(A_{X,B}|_{\sigma_{Y,E}}\) is homogeneous of degree one by
\eqref{eq:ch32-log-discrepancy-affine}; the supremum of homogeneous
functions is homogeneous.  At \(v_{\triv}\) every retraction returns the
cone point, where the affine formula gives \(0\).

\cref{itm:ch50-A:sign}  On a cone, \(A_{X,B}(v_\alpha)
=\sum_j\alpha_jA_{X,B}(\ord_{E_j})\) with \(\alpha_j\ge0\); sublc says
exactly that every \(A_{X,B}(\ord_{E_j})\ge0\) for every SNC pair, subklt
that every one is \(>0\).  For the strict statement, if
\(v\ne v_{\triv}\) then \(r_{Y,E}(v)\ne v_{\triv}\) for some \((Y,E)\)
--- otherwise \(v\) would be dominated by every retraction and hence
trivial --- and then the affine formula gives a strictly positive value.

\cref{itm:ch50-A:infinite}  Let \(v\) have proper support \(Z\subsetneq X\).
Choose an SNC pair \((Y,E)\) such that some component \(E_j\) dominates
\(Z\); the retractions along the tower of blowups of \(Z\) then have
\(A_{X,B}(r(v))\to+\infty\), because \(v(\mathcal I_Z)=+\infty\) forces the
retraction coefficient along the exceptional divisor of the \(m\)-th
blowup to grow without bound.  Scaling stability and the containment of
quasimonomial points are \cref{itm:ch50-A:homog} and
\cref{itm:ch50-A:consistent}.
\end{proof}

\begin{warning}{Log discrepancy and linear growth are different finiteness
conditions}
\label{warn:ch50-A-vs-T}
Both \(A_{X,B}\) and \(T_\omega\) are homogeneous, lower semicontinuous
gauges on \(X^{\An}\) that are finite on \(X^{\divis}\) and infinite at
proper-support points, and it is tempting to treat them as
interchangeable.  They are not.  On \(\mathbf P^2\) with \(L=\cO(1)\), the
divisorial valuation \(\ord_{C}\) of a smooth curve \(C\) of degree \(d\)
has
\[
 A_{\mathbf P^2}(\ord_C)=1,
 \qquad
 T_L(\ord_C)=\frac1d,
\]
the second by \Cref{prop:ch49-psef-threshold}, so the ratio
\(A/T_\omega\) is unbounded on \(X^{\divis}\), already in dimension two.
Consequently a bound on dual energy gives \emph{no} bound on entropy, which
is the content of \Cref{ex:ch50-infinite-entropy}; and there is no constant
\(C\) with \(A_{X,B}\le C\,T_\omega\).  The inequality that \emph{does}
hold in one direction --- \(T_\omega\le C_\omega A_{X,B}\) on
\(X^{\divis}\) for a constant depending on \(\omega\), by the Izumi-type
bounds --- is not used anywhere in this chapter, and we do not assert it.
\end{warning}

\subsection{Entropy}
\label{sec:ch50-entropy-def}

\begin{definition}[Entropy]
\label{def:stability-entropy}
For a probability measure \(\mu\) for which the extended integral is
defined, set
\[
                  \Ent_{X,B}(\mu)
                   :=\int_{X^{\An}}A_{X,B}\,\rd \mu.
\tag{50.45}\label{eq:stability-entropy}
\]
For a semipositive model metric \(\phi\), write
\(H^{\NA}_B(\phi)=\Ent_{X,B}(\MA(\phi))\).
\end{definition}

When \((X,B)\) is sublc the integrand is nonnegative by
\Cref{prop:ch50-A-properties}\,\cref{itm:ch50-A:sign} and the integral is
defined in \([0,+\infty]\) for every probability measure, with no
convention needed.  We assume sublc from here on whenever entropy appears;
the general case requires the truncation conventions of
\Cref{app:normalizations} and is not used.

\begin{theorem}
\label{thm:ch50-entropy-lsc}
\index{entropy!lower semicontinuity}
Assume \((X,B)\) sublc.  Then
\[
 \Ent_{X,B}:\ \{\text{probability measures on }X^{\An}\}
 \longrightarrow[0,+\infty]
\]
is convex and lower semicontinuous for the weak topology.  Moreover
\(\Ent_{X,B}(\delta_v)=A_{X,B}(v)\) for every \(v\in X^{\An}\), and
\(\Ent_{X,B}\) is affine along convex combinations, in the sense that
\(\Ent_{X,B}(\sum_ic_i\mu_i)=\sum_ic_i\Ent_{X,B}(\mu_i)\) for every
countable convex combination.
\end{theorem}

\begin{proof}
Affineness along convex combinations, hence convexity, is the linearity of
the integral together with the monotone convergence theorem, all terms
being nonnegative; and \(\Ent(\delta_v)=A_{X,B}(v)\) is the definition.

For lower semicontinuity, let \(\mathcal C\) denote the set of continuous
functions \(g:X^{\An}\to\mathbf R\) with \(g\le A_{X,B}\).  We claim
\[
 \Ent_{X,B}(\mu)=\sup_{g\in\mathcal C}\int g\,\rd\mu .
\tag{50.46}\label{eq:ch50-entropy-sup-continuous}
\]
Granting the claim, \(\Ent_{X,B}\) is a supremum of weakly continuous
affine functionals, hence weakly lower semicontinuous, which is the
assertion.

To prove the claim, note first that \(\le\) is clear.  For \(\ge\), fix
\(c<\Ent_{X,B}(\mu)\).  The function \(A_{X,B}\) is lower semicontinuous
and \([0,+\infty]\)-valued on the compact Hausdorff space \(X^{\An}\), and
\(\mu\) is a Radon probability measure.  By
\eqref{eq:ch50-A-sup-retractions}, \(A_{X,B}=\sup_{(Y,E)}g_{Y,E}\) with
\(g_{Y,E}:=A_{X,B}\circ r_{Y,E}\) continuous and real valued, and the
family \((g_{Y,E})\) is directed increasing by
\Cref{prop:ch50-A-properties}\,\cref{itm:ch50-A:consistent}.  Monotone
convergence for an increasing net of continuous functions on a compact
space, against a Radon measure --- that is, the fact that
\(\mu\bigl(\sup_ig_i\bigr)=\sup_i\mu(g_i)\) for a directed increasing net
of lower semicontinuous functions, which is inner regularity of \(\mu\)
--- gives \(\int A_{X,B}\rd\mu=\sup_{(Y,E)}\int g_{Y,E}\rd\mu\).  Hence
\(\int g_{Y,E}\rd\mu>c\) for some \((Y,E)\), and \(g_{Y,E}\in\mathcal C\).
\end{proof}

\begin{remark}[Why lower and not upper]
\label{rem:ch50-entropy-lsc-direction}
The direction matters.  Lower semicontinuity is what makes
\(\Ent_{X,B}\) --- and hence \(M^{\NA}\) --- suitable for a variational
problem in which one \emph{minimizes}: a minimizing net has a weak limit
point by compactness of the space of probability measures on the compact
space \(X^{\An}\), and lower semicontinuity ensures that no entropy is
gained in the limit.  Upper semicontinuity is false:
\Cref{ex:ch50-entropy-jumps} gives finite-energy probability measures
converging even in total variation to \(\delta_{v_{\triv}}\), while their
entropies stay equal to \(1\) and the entropy of the limit is \(0\).
\end{remark}

\begin{example}[Entropy is not upper semicontinuous]
\label{ex:ch50-entropy-jumps}
Let \(X=\mathbf P^2\), \(B=0\), and \(L=\cO(1)\).  For each \(m\ge2\),
choose a smooth plane curve \(C_m\) of degree \(m\) and set
\(v_m:=m\,\ord_{C_m}\), as in \Cref{ex:ch50-infinite-entropy}.  Then
\(T_\omega(v_m)=1\) and \(A_{\mathbf P^2}(v_m)=m\).  The probability
measures
\[
 \nu_m:=\left(1-\frac1m\right)\delta_{v_{\triv}}
          +\frac1m\delta_{v_m}
\]
belong to \(\cM^1(\omega)\), since convexity and
\eqref{eq:bj-dirac-energy-T} give
\(E^\vee_\omega(\nu_m)\le2/(3m)\).  They converge to
\(\delta_{v_{\triv}}\) in total variation, whereas
\[
 \Ent_{\mathbf P^2}(\nu_m)=\frac1mA_{\mathbf P^2}(v_m)=1,
 \qquad
 \Ent_{\mathbf P^2}(\delta_{v_{\triv}})=0.
\]
Thus entropy is not upper semicontinuous, even on \(\cM^1(\omega)\).
\end{example}

\subsection{The model formula}
\label{sec:ch50-model-entropy}

If \(\cX_0=\sum_Eb_EE\) and \(\overline\cL\) is the compactified
polarization, then
\[
 H^{\NA}_B(\phi)
 =\frac1V\sum_E
     A_{X,B}(v_E)b_E(E\cdot\overline\cL^{\,n}),
\qquad
 v_E=b_E^{-1}r(\ord_E).
\tag{50.47}\label{eq:stability-model-entropy}
\]
Thus the entropy is visibly nonnegative when \((X,B)\) is sublc.
On \(\mathcal M^1\) it may equal \(+\infty\).  This is intentional: finite
Monge--Amp\`ere energy does not imply finite log discrepancy energy.
Finite-energy measures ignore proper-support pluripolar points, but
\eqref{eq:stability-entropy} still uses the lower semicontinuous extension,
not an arbitrary value assigned off \(X^{\val}\).

\begin{proposition}
\label{prop:ch50-model-entropy}
Let \((\cX,\cL)\) be a normal test configuration for \((X,L)\) with
\(\cL\) relatively semiample, and let \(\phi=\phi_\cL\) be its PL
function.  Then \eqref{eq:stability-model-entropy} holds, both sides being
finite; in particular \(H_B^{\NA}(\phi)\) is a finite convex combination of
log discrepancies of divisorial valuations.
\end{proposition}

\begin{proof}
By \Cref{prop:ch46-pl-measure} and the atomic formula
\eqref{eq:bj-model-ma-atomic}, the unnormalized Monge--Amp\`ere measure of
\(\phi\) is
\[
 \MA^{\mathrm u}_\omega(\phi)=\sum_Eb_E\,(E\cdot\overline\cL^{\,n})\,
 \delta_{v_E},\qquad v_E=b_E^{-1}r(\ord_E),
\]
a finite sum over the irreducible components \(E\) of \(\cX_0\), with
nonnegative coefficients summing to \((\cX_0\cdot\overline\cL^n)=V\).
Dividing by \(V\) gives the probability measure \(\MA_\omega(\phi)\).
Each \(v_E\) is a divisorial valuation of \(k(X)\) by
\Cref{def:ch33-component-valuation} and \Cref{lem:ch33-gauss-component},
so \(A_{X,B}(v_E)\) is finite and given by
\eqref{eq:stability-log-discrepancy}.  Integrating \(A_{X,B}\) against an
atomic measure evaluates it at the atoms, which is
\eqref{eq:stability-model-entropy}.  Nonnegativity under the sublc
hypothesis is \Cref{prop:ch50-A-properties}\,\cref{itm:ch50-A:sign}.
\end{proof}

\subsection{Entropy is not finite on finite-energy measures}
\label{sec:ch50-entropy-infinite}

\begin{example}[A finite-energy measure of infinite entropy]
\label{ex:ch50-infinite-entropy}
\index{entropy!infinite on \(\cM^1\)}
Let \(X=\mathbf P^2\) over an algebraically closed field of
characteristic zero, \(B=0\), \(L=\cO(1)\), \(\omega=c_1(L)\), \(V=1\).
For \(d\ge2\) choose a smooth plane curve \(C_d\) of degree \(d\) and put
\[
 v_d:=d\cdot\ord_{C_d}\in X^{\divis},
 \qquad
 c_d:=\frac{\kappa}{d^2},\quad
 \kappa:=\Bigl(\sum_{d\ge2}d^{-2}\Bigr)^{-1},
\]
so that \(\sum_{d\ge2}c_d=1\).  Set
\(\mu:=\sum_{d\ge2}c_d\,\delta_{v_d}\), a Radon probability measure on
\(X^{\An}\).  Then
\[
 E^\vee_\omega(\mu)\le\frac23<\infty,
 \qquad
 \Ent_{\mathbf P^2}(\mu)=+\infty .
\tag{50.48}\label{eq:ch50-infinite-entropy}
\]
Thus \(\mu\in\cM^1(\omega)\) and \(\Ent(\mu)=+\infty\).
\end{example}

\begin{proof}
\emph{The invariants of \(v_d\).}  By
\Cref{prop:ch49-psef-threshold}, \(T_L(\ord_{C_d})=\tau_L(C_d)\) is the
pseudoeffective threshold of \(C_d\) in \(\cO(1)\); since
\(\cO(1)-c\,C_d\) has degree \(1-cd\) on \(\mathbf P^2\), it is
pseudoeffective exactly for \(c\le1/d\), so \(T_L(\ord_{C_d})=1/d\).  By
homogeneity, \(T_\omega(v_d)=d\cdot(1/d)=1\).  On the other hand
\(C_d\) is a prime divisor on \(X\) itself, so
\(A_{\mathbf P^2}(\ord_{C_d})=1\) by
\eqref{eq:stability-log-discrepancy}, and
\(A_{\mathbf P^2}(v_d)=d\) by homogeneity
(\Cref{prop:ch50-A-properties}\,\cref{itm:ch50-A:homog}).

\emph{Finite dual energy.}  By \eqref{eq:bj-dirac-energy-T},
\(E^\vee_\omega(\delta_{v_d})\le\frac{n}{n+1}T_\omega(v_d)
=\frac23\) for every \(d\).  Let
\(\mu_N:=\sum_{2\le d\le N}c_d\delta_{v_d}
+\bigl(1-\sum_{2\le d\le N}c_d\bigr)\delta_{v_{\triv}}\), a probability
measure.  Since \(E^\vee_\omega\) is convex
(\Cref{prop:ch47-dual-elementary}\,\cref{itm:ch47-dual:convex}) and
\(E^\vee_\omega(\delta_{v_{\triv}})=0\) by
\Cref{cor:ch47-dual-zero}, we get \(E^\vee_\omega(\mu_N)\le\frac23\).  The
net \(\mu_N\) converges weakly to \(\mu\), and \(E^\vee_\omega\) is weakly
lower semicontinuous
(\Cref{prop:ch47-dual-elementary}\,\cref{itm:ch47-dual:lsc}), so
\(E^\vee_\omega(\mu)\le\frac23\).

\emph{Infinite entropy.}  By \Cref{thm:ch50-entropy-lsc},
\(\Ent(\mu)=\sum_{d\ge2}c_dA_{\mathbf P^2}(v_d)
=\kappa\sum_{d\ge2}d^{-2}\cdot d=\kappa\sum_{d\ge2}d^{-1}=+\infty\).
\end{proof}

\begin{remark}[What the example does and does not show]
\label{rem:ch50-infinite-entropy-scope}
It shows that \(\Ent_{X,B}\) is not finite on \(\cM^1(\omega)\), already
for \(X=\mathbf P^2\) with its most classical polarization and with
\(B=0\), a klt --- indeed smooth --- pair.  It does \emph{not} show that
the sublevel sets of \(\Ent\) fail to be weakly compact: they do not,
\(\Ent\) being lower semicontinuous on a compact space.  Nor does it
contradict the finiteness of the entropy of a \emph{divisorial} measure
\eqref{eq:stability-divisorial-measure}, which is a finite convex
combination of Dirac masses at divisorial points and therefore has finite
entropy.  The mechanism is exactly the unbounded ratio
\(A/T_\omega\) recorded in \Cref{warn:ch50-A-vs-T}: an infinite convex
combination can keep \(T_\omega\), hence \(E^\vee\), bounded while letting
\(A\) escape.  In particular there is no ``entropy \(\le\) constant
\(\times\) energy'' estimate to be had, and the Mabuchi functional of
\eqref{eq:stability-M-all-norms} genuinely takes the value \(+\infty\) on a
dense set of directions.
\end{remark}

\begin{proposition}
\label{prop:ch50-entropy-finite-cases}
Assume \((X,B)\) sublc.
\begin{enumerate}
\item\label{itm:ch50-ef:divisorial}
      If \(\mu=\sum_{i=1}^ra_i\delta_{v_i}\) is a divisorial measure as in
      \eqref{eq:stability-divisorial-measure}, then
      \(\Ent_{X,B}(\mu)=\sum_ia_iA_{X,B}(v_i)<\infty\).
\item\label{itm:ch50-ef:model}
      If \(\phi\) is a semipositive model metric then
      \(H^{\NA}_B(\phi)<\infty\), by
      \Cref{prop:ch50-model-entropy}.
\item\label{itm:ch50-ef:dense}
      The set of \(\mu\in\cM^1(\omega)\) with \(\Ent_{X,B}(\mu)<\infty\)
      is a convex, strongly dense subset of \(\cM^1(\omega)\) containing
      all divisorial measures; it is not closed, and it is not all of
      \(\cM^1(\omega)\) by \Cref{ex:ch50-infinite-entropy}.
\end{enumerate}
\end{proposition}

\begin{proof}
\cref{itm:ch50-ef:divisorial} and \cref{itm:ch50-ef:model} are immediate.
For \cref{itm:ch50-ef:dense}, convexity is
\Cref{thm:ch50-entropy-lsc}.  For density, finite-energy potentials admit
decreasing PL approximants by \Cref{lem:ch46-pl-domination}, and these
converge strongly by \Cref{prop:ch46-energy-usc}.  The bi-Lipschitz estimate
of \Cref{prop:bj-ma-strong-embedding} makes their Monge--Amp\`ere measures
converge strongly, while \Cref{prop:ch46-pl-measure} makes each approximating
measure divisorial and hence of finite entropy by
\cref{itm:ch50-ef:divisorial}.  Surjectivity of the Monge--Amp\`ere map is
\Cref{thm:ch48-ma-homeomorphism}.  Thus finite-entropy measures are strongly
dense.  They are not closed because \(\mu_N\to\mu\) strongly in
\Cref{ex:ch50-infinite-entropy}.  Indeed, if \(\sup\phi=0\), then the
definition of \(T_\omega\) and \(T_\omega(v_d)=1\) give
\[
 \left|\int\phi\,\rd(\mu-\mu_N)\right|
 \le\sum_{d>N}c_d\,|\phi(v_d)-\phi(v_{\triv})|
 \le\sum_{d>N}c_d\longrightarrow0.
\]
The estimate is uniform on every \(J\)-sublevel (in fact on all normalized
finite-energy potentials), while \(\mu_N\to\mu\) weakly.  The uniform
pairing criterion of \Cref{thm:bj-strong-measures} therefore gives strong
convergence.  Each \(\mu_N\) has finite entropy and \(\mu\) has infinite
entropy, proving nonclosedness.
\end{proof}

\subsection{The strongly lower semicontinuous extension}
\label{sec:ch50-strong-lsc}

\begin{important}{Imported: the strong extension of \(\beta\)}
\label{imp:ch50-strong-lsc}
The Mabuchi functional will be defined on all of \(\mathcal N_{\mathbf R}\)
by \eqref{eq:stability-M-all-norms}, through the measure-theoretic
expression \eqref{eq:stability-beta-measure}.  Two facts about that
expression are imported from \cite{BJKS2} (preprint June 2022, published
2023):
\begin{enumerate}
\item\label{itm:ch50-slsc:diff}
      for \(\mu\in\cM^1\), the function \(\omega\mapsto\|\mu\|_\omega
      :=E^\vee_\omega(\mu)\) is of class \(C^1\) on the ample cone
      \(\Amp(X)\subset N^1(X)_{\mathbf R}\)
      \cite[Theorem~2.15]{BJKS2}; this is what gives meaning to the
      directional derivative \(\nabla_{K_{X,B}}\|\mu\|_\omega\) in
      \eqref{eq:stability-beta-measure};
\item\label{itm:ch50-slsc:extension}
      when \((X,B)\) is subklt, the functional
      \(\beta_{X,B;\omega}\) of \eqref{eq:stability-beta-measure} is the
      maximal strongly lower semicontinuous extension, from divisorial
      measures to \(\cM^1(\omega)\), of its own restriction to divisorial
      measures \cite[\S\S3.4 and 4.1]{BJKS2}.
\end{enumerate}
Neither statement is proved here.  \Cref{thm:ch50-entropy-lsc} proves
\emph{weak} lower semicontinuity of the first summand of
\eqref{eq:stability-beta-measure}; the second summand is continuous for the
strong topology by \cref{itm:ch50-slsc:diff} and
\Cref{thm:bj-strong-measures}, but is only upper semicontinuous for the
weak one, so the sum is strongly, not weakly, lower semicontinuous.  This
asymmetry is the reason \Cref{thm:stability-divisorial-filtrations} is a
statement about the \emph{strong} topology.
\end{important}

\section{The Mabuchi functional and the Donaldson--Futaki invariant}
\label{sec:ch50-mabuchi}

\subsection{The average scalar curvature and the Ricci energy}
\label{sec:ch50-ricci}

Put
\[
           \overline S_B
           :=-\frac{n(K_{X,B}\cdot L^{n-1})}{V}.
\tag{50.49}\label{eq:stability-average-scalar}
\]
If a normal test configuration dominates the product and
\(p:\overline\cX\to X\) is the induced morphism, its non-Archimedean Ricci
energy is
\[
        R_B^{\NA}(\phi)
        :=\frac1V(p^*K_{X,B}\cdot\overline\cL^{\,n}).
\tag{50.50}\label{eq:stability-ricci-energy}
\]

\begin{lemma}
\label{lem:ch50-ricci-energy}
Let \(\phi\) be a semipositive model metric on \(L\), determined on a
normal test configuration \(\cX\) dominating \(X\times\mathbf A^1\).  Then
\[
 R_B^{\NA}(\phi)=\frac1V\,(K_{X,B},0)\,(\omega,\phi)^n
\tag{50.51}\label{eq:ch50-ricci-pairing}
\]
in the notation of \Cref{def:bj-pl-energy-pairing}; in particular
\(R^{\NA}_B(\phi)\) does not depend on the model chosen, and it is
determined by the numerical class \(K_{X,B}\).  If moreover
\(K_{X,B}\equiv\gamma\,\omega\) numerically for some
\(\gamma\in\mathbf R\), then
\[
 R^{\NA}_B(\phi)
 =\gamma\Bigl((n+1)E_\omega(\phi)-\int_{X^{\An}}\phi\,\MA_\omega(\phi)\Bigr),
 \qquad
 \overline S_B=-n\gamma .
\tag{50.52}\label{eq:ch50-ricci-proportional}
\]
\end{lemma}

\begin{proof}
The pair \((K_{X,B},0)\) is realized on \(\overline\cX\) by the line bundle
\(p^*K_{X,B}\) with vanishing vertical part, so
\eqref{eq:ch50-ricci-pairing} is
\eqref{eq:bj-pl-energy-intersection}.  Model independence and
multilinearity are \Cref{prop:ch46-pairing-wellposed}.

Assume \(K_{X,B}\equiv\gamma\omega\).  Then
\((K_{X,B},0)=\gamma(\omega,0)\) as pairs, and
\((\omega,0)=(\omega,\phi)-(0,\phi)\), so
\[
 (K_{X,B},0)(\omega,\phi)^n
 =\gamma\Bigl((\omega,\phi)^{n+1}-(0,\phi)(\omega,\phi)^n\Bigr).
\]
By \eqref{eq:bj-energy-definition}, \((\omega,\phi)^{n+1}=(n+1)VE_\omega(\phi)\);
by \Cref{prop:ch46-pl-measure}\,\cref{itm:ch46-pl-ma:duality} applied with
\(f=\phi\) and all pairs equal to \((\omega,\phi)\),
\((0,\phi)(\omega,\phi)^n=\int\phi\,\MA^{\mathrm u}_\omega(\phi)
=V\int\phi\,\MA_\omega(\phi)\).  Dividing by \(V\) gives
\eqref{eq:ch50-ricci-proportional}.  Finally
\((K_{X,B}\cdot L^{n-1})=\gamma(\omega^n)=\gamma V\), so
\(\overline S_B=-n\gamma\).
\end{proof}

\subsection{Two intersection-theoretic invariants}
\label{sec:ch50-df-def}

\begin{definition}[Donaldson--Futaki and Mabuchi invariants]
\label{def:ch50-df}
\index{Donaldson--Futaki invariant}\index{Mabuchi functional!non-Archimedean}
Let \((\cX,\cL)\) be a normal test configuration for \((X,L)\) with
\(\cL\) relatively semiample, and let
\((\overline\cX,\overline\cL)\) be its \(\mathbf P^1\)-compactification.
Write \(B_{\overline\cX}\) for the strict transform of \(B\).  Put
\[
 \begin{split}
 \DF_B(\cX,\cL)
 &:=\frac1V\bigl(K_{(\overline\cX,B_{\overline\cX})/\mathbf P^1}
       \cdot\overline\cL^{\,n}\bigr)
   +\frac{\overline S_B}{(n+1)V}\bigl(\overline\cL^{\,n+1}\bigr),\\
 M^{\NA}_B(\cX,\cL)
 &:=\frac1V\bigl(K^{\log}_{(\overline\cX,B_{\overline\cX})/\mathbf P^1}
       \cdot\overline\cL^{\,n}\bigr)
   +\frac{\overline S_B}{(n+1)V}\bigl(\overline\cL^{\,n+1}\bigr),
 \end{split}
\tag{50.53}\label{eq:ch50-df-M-definitions}
\]
where
\(K_{(\overline\cX,B_{\overline\cX})/\mathbf P^1}
=K_{\overline\cX}+B_{\overline\cX}-\pi^*K_{\mathbf P^1}\) and
\(K^{\log}_{(\overline\cX,B_{\overline\cX})/\mathbf P^1}
=K_{(\overline\cX,B_{\overline\cX})/\mathbf P^1}
+\cX_{0,\mathrm{red}}-\cX_0\).
\end{definition}

The first line of \eqref{eq:ch50-df-M-definitions} agrees with the
normalization \eqref{eq:ytd-df} of \Cref{def:ytd-df}, where the same
invariant is used on the complex side; \Cref{app:normalizations} records
the sign conventions once and for all.

\subsection{The Chen--Tian formula}
\label{sec:ch50-chen-tian}

\begin{theorem}\label{thm:stability-chen-tian}
For a semipositive model metric,
\[
             M_B^{\NA}(\phi)
              =H_B^{\NA}(\phi)+R_B^{\NA}(\phi)
                              +\overline S_BE(\phi).
\tag{50.54}\label{eq:stability-chen-tian}
\]
If \((\cX,\cL)\) is its normal representative, then
\[
 M_B^{\NA}(\phi)
  =\DF_B(\cX,\cL)
   +\frac1V\bigl((\cX_{0,\mathrm{red}}-\cX_0)
                        \cdot\overline\cL^{\,n}\bigr)
  \leq\DF_B(\cX,\cL).
\tag{50.55}\label{eq:stability-M-DF-correction}
\]
Equality holds when the central fibre is reduced, and it holds after every
sufficiently divisible normalized base change.
\end{theorem}

\begin{proof}
\emph{Step 1: the log relative canonical divisor of a test configuration.}
Let \(\cX\) be a normal test configuration dominating
\(X\times\mathbf A^1\), with \(\rho:\overline\cX\to
S:=X\times\mathbf P^1\), \(p=\operatorname{pr}_X\circ\rho\) and
\(\pi:\overline\cX\to\mathbf P^1\).  Since \(K_S=\operatorname{pr}_X^*K_X
+\operatorname{pr}_{\mathbf P^1}^*K_{\mathbf P^1}\) and
\(B_S=\operatorname{pr}_X^*B\), the discrepancy divisor
\[
 D:=K_{\overline\cX}+B_{\overline\cX}-\rho^*(K_S+B_S)
\]
is supported on \(\cX_0\), because \(\rho\) is an isomorphism over
\(S\setminus(X\times\{0\})\).  Substituting and using
\(\pi^*[0]=\cX_0\),
\[
 K^{\log}_{(\overline\cX,B_{\overline\cX})/\mathbf P^1}
 =p^*K_{X,B}+D+\cX_{0,\mathrm{red}}-\cX_0 .
\tag{50.56}\label{eq:ch50-log-canonical-decomposition}
\]

\emph{Step 2: the coefficients are log discrepancies.}  Let \(E\) be an
irreducible component of \(\cX_0\), with multiplicity
\(b_E=\ord_E(\varpi)\) and associated divisorial valuation
\(v_E=b_E^{-1}r(\ord_E)\in X^{\divis}\) as in
\Cref{def:ch33-component-valuation}.  By definition of the log discrepancy
of the pair \((S,B_S)\),
\(\ord_E(D)=A_{S,B_S}(\ord_E)-1\).  By
\Cref{lem:ch33-gauss-component}, \(\ord_E=b_E\,\sigma(v_E)\) with
\(\sigma\) the Gauss extension, and by
\Cref{lem:ch40-gauss-log-discrepancy} together with homogeneity of log
discrepancy \eqref{eq:appJ-discrepancy-homogeneous},
\[
 A_{S,B_S}(\ord_E)=b_E\,A_{S,B_S}(\sigma(v_E))
 =b_E\bigl(A_{X,B}(v_E)+1\bigr).
\]
Since \(\ord_E(\cX_{0,\mathrm{red}})=1\) and
\(\ord_E(\cX_0)=b_E\), the coefficient of \(E\) in the vertical part of
\eqref{eq:ch50-log-canonical-decomposition} is
\[
 \bigl(A_{S,B_S}(\ord_E)-1\bigr)+1-b_E
 =b_EA_{X,B}(v_E)+b_E-b_E
 =b_E\,A_{X,B}(v_E).
\]
Hence
\[
 K^{\log}_{(\overline\cX,B_{\overline\cX})/\mathbf P^1}
 =p^*K_{X,B}+\sum_{E\subset\cX_0}b_E\,A_{X,B}(v_E)\,E .
\tag{50.57}\label{eq:ch50-log-canonical-atoms}
\]

\emph{Step 3: intersecting.}  Intersect
\eqref{eq:ch50-log-canonical-atoms} with \(\overline\cL^{\,n}\) and divide
by \(V\).  The first term gives \(R^{\NA}_B(\phi)\) by
\eqref{eq:stability-ricci-energy}; the second gives
\(\frac1V\sum_Eb_EA_{X,B}(v_E)(E\cdot\overline\cL^n)=H^{\NA}_B(\phi)\) by
\eqref{eq:stability-model-entropy}.  The remaining term of
\eqref{eq:ch50-df-M-definitions} is
\(\frac{\overline S_B}{(n+1)V}(\overline\cL^{n+1})
=\overline S_BE_\omega(\phi)\) by
\eqref{eq:bj-energy-definition}.  This is
\eqref{eq:stability-chen-tian}.

\emph{Step 4: comparison with \(\DF\).}  Subtracting the two lines of
\eqref{eq:ch50-df-M-definitions} gives exactly the first equality of
\eqref{eq:stability-M-DF-correction}.  The divisor
\(\cX_0-\cX_{0,\mathrm{red}}=\sum_E(b_E-1)E\) is effective and supported in
the central fibre; each \(E\) is contained in a fibre of \(\pi\), so
\(\overline\cL|_E\) is nef because \(\overline\cL\) is relatively
semiample, whence \((E\cdot\overline\cL^{\,n})\ge0\).  Therefore the
correction term is \(\le0\) and
\(M^{\NA}_B\le\DF_B\), with equality when every \(b_E=1\), that is when
\(\cX_0\) is reduced.

\emph{Step 5: base change.}  Let \(d\ge1\) and let \((\cX_d,\cL_d)\) be the
normalization of the base change of \((\cX,\cL)\) along
\(\varpi\mapsto\varpi^d\).  Both \(\DF_B\) and \(M^{\NA}_B\) are
homogeneous of degree one under this operation, by
\eqref{eq:ch33-intersection-base-change}; and by semistable reduction in
codimension one --- applied at the finitely many generic points of
\(\cX_0\), where \(\cX\) is a discrete valuation ring by
\Cref{thm:ch33-integrally-closed-criterion} --- the central fibre of
\(\cX_d\) is reduced as soon as every \(b_E\) divides \(d\).  For such
\(d\), Step 4 gives \(M^{\NA}_B(d\cdot\phi)=\DF_B(\cX_d,\cL_d)\).  See also
\cite[Propositions~7.15--7.16 and 7.22]{BHJ17}.
\end{proof}

\begin{remark}[Provenance and the two normalizations]
\label{rem:ch50-chen-tian-provenance}
\Cref{thm:stability-chen-tian} is \cite[\S7]{BHJ17} (preprint 2015,
published 2017); \Cref{def:ch50-df} matches
\cite[Definitions~7.13, 7.17 and 7.20]{BHJ17}.  The proof above is written
out because two of its steps are places where a sign or a multiplicity is
easily lost.  The first is Step 2: the coefficient of \(E\) is
\(b_EA_{X,B}(v_E)\) and not \(A_{X,B}(v_E)\) --- the multiplicity enters
twice, once through \(\ord_E=b_E\sigma(v_E)\) and once through
\(\cX_{0,\mathrm{red}}-\cX_0\), and the two occurrences cancel exactly one
of the two \(b_E\)'s.  The second is Step 4: the correction is
\emph{nonpositive}, so \(M^{\NA}\le\DF\) and not the other way round; a
sign error here reverses every stability implication.  The unit test is
\Cref{ex:ch33-dnc-p1}, whose central fibre is reduced, so that
\(M^{\NA}=\DF\) there, together with a nonreduced example such as the
degree-\(d\) base change of it, where the two differ by a computable
amount.
\end{remark}

\subsection{The measure formulation}
\label{sec:ch50-beta}

There is also an intrinsic measure formula.  Write
\(\|\mu\|_\omega=E_\omega^\vee(\mu)\) for \(\omega=c_1(L)\).  For
\(\mu\in\mathcal M^1\), the function
\(\omega\mapsto\|\mu\|_\omega\) is \(C^1\) on the ample cone
\cite[Theorem~2.15]{BJKS2}.  Hence
\[
 \beta_{X,B;\omega}(\mu)
 :=\Ent_{X,B}(\mu)
       +\nabla_{K_{X,B}}\|\mu\|_\omega
 \in\mathbf R\cup\{+\infty\}.
\tag{50.58}\label{eq:stability-beta-measure}
\]
For model data,
\[
       M_B^{\NA}(\phi)
       =\beta_{X,B;\omega}(\MA(\phi)).
\tag{50.59}\label{eq:stability-M-beta}
\]
When \((X,B)\) is subklt, the right side of
\eqref{eq:stability-beta-measure} is the maximal strongly lower
semicontinuous extension of its restriction to divisorial measures
\cite[\S\S3.4 and 4.1]{BJKS2}.  This supplies the later definition
\[
       M_{X,B;L}(\chi)
       :=\beta_{X,B;c_1(L)}(\MA(\chi)),
       \qquad\chi\in\mathcal N_{\mathbf R}.
\tag{50.60}\label{eq:stability-M-all-norms}
\]
It is allowed to take the value \(+\infty\).

The proportional case of \eqref{eq:stability-M-beta} is elementary, and it
is the case used in every example of this chapter.  We prove it.

\begin{proposition}
\label{prop:ch50-M-proportional}
Assume \(K_{X,B}\equiv\gamma\,\omega\) numerically, \(\gamma\in\mathbf R\).
Then for every semipositive model metric \(\phi\),
\[
 M^{\NA}_B(\phi)
 =H^{\NA}_B(\phi)+\gamma\,E^\vee_\omega\bigl(\MA_\omega(\phi)\bigr)
 =\Ent_{X,B}\bigl(\MA_\omega(\phi)\bigr)
   +\gamma\,\bigl\|\MA_\omega(\phi)\bigr\|_\omega ,
\tag{50.61}\label{eq:ch50-M-proportional}
\]
and \eqref{eq:stability-M-beta} holds, the derivative in
\eqref{eq:stability-beta-measure} being
\(\nabla_{\gamma\omega}\|\mu\|_\omega=\gamma\|\mu\|_\omega\).  More
generally, for any \(\mu\in\cM^1(\omega)\) the right-hand side of
\eqref{eq:ch50-M-proportional} computes
\(\beta_{X,B;\omega}(\mu)\) without any appeal to
\cite[Theorem~2.15]{BJKS2}.
\end{proposition}

\begin{proof}
By \Cref{thm:stability-chen-tian} and
\Cref{lem:ch50-ricci-energy},
\[
 \begin{split}
 M^{\NA}_B(\phi)
 &=H^{\NA}_B(\phi)
  +\gamma\Bigl((n+1)E_\omega(\phi)-\int\phi\,\MA_\omega(\phi)\Bigr)
  -n\gamma E_\omega(\phi)\\
 &=H^{\NA}_B(\phi)
  +\gamma\Bigl(E_\omega(\phi)-\int\phi\,\MA_\omega(\phi)\Bigr),
 \end{split}
\]
and the bracket is \(E^\vee_\omega(\MA_\omega(\phi))\) by
\eqref{eq:ch50-min-norm-identities}.  For the last assertion, if
\(\omega'=\omega+s\theta\) with \(\theta=\gamma\omega\) then
\(\omega'=(1+s\gamma)\omega\) and \(\|\mu\|_{(1+s\gamma)\omega}
=(1+s\gamma)\|\mu\|_\omega\); indeed for \(t>0\) the substitution
\(\phi\mapsto t\phi\) in \eqref{eq:bj-dual-energy} gives
\(E_{t\omega}(t\phi)=tE_\omega(\phi)\) by
\Cref{thm:ch46-energy-homogeneous}, whence
\(\|\mu\|_{t\omega}=t\|\mu\|_\omega\).  Differentiating at \(s=0\) gives
\(\nabla_{\gamma\omega}\|\mu\|_\omega=\gamma\|\mu\|_\omega\), a genuine
derivative of an affine function of \(s\), so no differentiability theorem
is needed.
\end{proof}

\begin{corollary}
\label{cor:ch50-M-dirac}
Assume \(K_{X,B}\equiv\gamma\omega\) and \(v\in X^{\lin}\).  Then
\[
 \beta_{X,B;\omega}(\delta_v)=A_{X,B}(v)+\gamma\,E^\vee_\omega(\delta_v).
\tag{50.62}\label{eq:ch50-beta-dirac}
\]
In the log Fano case \(L=-K_{X,B}\), so \(\gamma=-1\), this reads
\(\beta_{X,B;\omega}(\delta_v)=A_{X,B}(v)-E^\vee_\omega(\delta_v)\), which
by \Cref{imp:ch50-imported}\,\cref{itm:ch50-import:S} is the Fujita--Li
invariant \(\beta_L(v)=A_{X,B}(v)-S_L(v)\) of
\eqref{eq:stability-valuative-invariants}.
\end{corollary}

\begin{proof}
Immediate from \Cref{prop:ch50-M-proportional} and
\(\Ent(\delta_v)=A_{X,B}(v)\).
\end{proof}

\begin{remark}[What \eqref{eq:stability-M-all-norms} does and does not
say]
\label{rem:ch50-M-all-norms-scope}
Definition \eqref{eq:stability-M-all-norms} assigns a value in
\(\mathbf R\cup\{+\infty\}\) to \emph{every} linearly bounded
multiplicative filtration, by first passing to the associated
finite-energy measure through
\eqref{eq:stability-norm-ma-isometry} and then evaluating
\eqref{eq:stability-beta-measure}.  Three cautions.

First, the passage \(\chi\mapsto\MA(\chi)\) is not the naive one: for a
general \(\chi\) the function \(\FS(\chi)\) need not be an
\(\cE^1\)-potential of the measure, and the measure is produced by the
isometric extension of \Cref{thm:stability-norm-completion}\,(2).  For
finite-type \(\chi\) it is the naive one, by
\Cref{prop:ch46-pl-measure}.

Second, the value \(+\infty\) genuinely occurs, by
\Cref{ex:ch50-infinite-entropy} together with
\Cref{prop:ch50-M-proportional}: the measure constructed there has finite
\(\|\cdot\|_\omega\) and infinite entropy, so
\(\beta_{X,B;\omega}=+\infty\) on it.  Any inequality of the form
\(M\ge\varepsilon\|\cdot\|_{\min}\) is therefore vacuous at such a point,
and stability is a statement about the directions where \(M\) is finite.

Third, \eqref{eq:stability-M-all-norms} is not a limit of the
model-theoretic \(M^{\NA}\) along an approximating net: the entropy is only
lower semicontinuous, so approximation gives \(\liminf\), and equality
requires the strong lower semicontinuous extension statement of
\Cref{imp:ch50-strong-lsc}\,\cref{itm:ch50-slsc:extension}.  That statement
is imported, and it is exactly what
\Cref{thm:stability-divisorial-filtrations} needs.
\end{remark}

\section{The Ding functional}
\label{sec:ch50-ding}

\subsection{The log Fano hypothesis and the functional \texorpdfstring{\(L^{\NA}\)}{L NA}}
\label{sec:ch50-L}

Assume that \((X,B)\) is log Fano: it is klt, \(B\geq0\), and
\(L=-K_{X,B}\) is ample.  For a model metric \(\phi\), define
\[
 \begin{split}
 L_B^{\NA}(\phi)
   &:=\inf_{v\in X^{\divis}\cup\{v_{\triv}\}}
        \{A_{X,B}(v)+(\phi-\phi_{\triv})(v)\},\\
 D_B^{\NA}(\phi)&:=L_B^{\NA}(\phi)-E(\phi).
 \end{split}
\tag{50.63}\label{eq:stability-Ding}
\]
The trivial valuation is included, so the infimum is never over an empty
set.  Both summands change by the same constant under translation of
\(\phi\), and \(D_B^{\NA}\) is therefore translation invariant.

Throughout this section \(\phi_{\triv}=0\) is the trivial metric, so that
\((\phi-\phi_{\triv})(v)=\phi(v)\); we write \(\gamma=-1\) for the
proportionality constant of \Cref{lem:ch50-ricci-energy}, since
\(K_{X,B}=-L\), and consequently \(\overline S_B=n\).

\begin{proposition}
\label{prop:ch50-L-elementary}
Let \((X,B)\) be log Fano and let \(\phi,\psi\) be model metrics on
\(L=-K_{X,B}\).  Then:
\begin{enumerate}
\item\label{itm:ch50-L:transl}
      \(L^{\NA}_B(\phi+c)=L^{\NA}_B(\phi)+c\) for \(c\in\mathbf R\), and
      \(L^{\NA}_B(0)=0\);
\item\label{itm:ch50-L:mono}
      \(\phi\le\psi\) implies \(L^{\NA}_B(\phi)\le L^{\NA}_B(\psi)\);
\item\label{itm:ch50-L:lip}
      \(|L^{\NA}_B(\phi)-L^{\NA}_B(\psi)|\le\sup_{X^{\divis}}|\phi-\psi|\);
\item\label{itm:ch50-L:concave}
      \(L^{\NA}_B\) is concave: \(L^{\NA}_B((1-t)\phi+t\psi)\ge
      (1-t)L^{\NA}_B(\phi)+tL^{\NA}_B(\psi)\) for \(t\in[0,1]\);
\item\label{itm:ch50-L:bounds}
      \(L^{\NA}_B(\phi)\le\sup_{X^{\An}}\phi\), hence
      \(D^{\NA}_B(\phi)\le J(\phi)\);
\item\label{itm:ch50-L:scaling}
      \(L^{\NA}_B(t\cdot\phi)=t\,L^{\NA}_B(\phi)\) for the scaling action
      of \(t>0\).
\end{enumerate}
\end{proposition}

\begin{proof}
\cref{itm:ch50-L:transl}  Adding \(c\) to \(\phi\) adds \(c\) to every term
of the infimum.  For \(L^{\NA}_B(0)=0\): the term at \(v_{\triv}\) is
\(A_{X,B}(v_{\triv})+0=0\) by
\Cref{prop:ch50-A-properties}\,\cref{itm:ch50-A:homog}, and every other term
is \(A_{X,B}(v)\ge0\) by klt-ness and
\Cref{prop:ch50-A-properties}\,\cref{itm:ch50-A:sign}.

\cref{itm:ch50-L:mono} and \cref{itm:ch50-L:lip}  An infimum of a family of
functions is monotone and \(1\)-Lipschitz in the family for the uniform
distance; the supremum may be taken over \(X^{\divis}\) alone because that
is where the infimum is taken, and the value at \(v_{\triv}\) is the
supremum of \(\phi\), which is a limit of values at divisorial points.

\cref{itm:ch50-L:concave}  Each map
\(\phi\mapsto A_{X,B}(v)+\phi(v)\) is affine in \(\phi\), and an infimum of
affine functions is concave.

\cref{itm:ch50-L:bounds}  Take \(v=v_{\triv}\) in the infimum:
\(L^{\NA}_B(\phi)\le0+\phi(v_{\triv})=\sup\phi\) by
\Cref{lem:ch47-three-normalizations}\,\cref{itm:ch47-norm:max}.  Subtracting
\(E(\phi)\) gives \(D^{\NA}_B(\phi)\le\sup\phi-E(\phi)=J(\phi)\).

\cref{itm:ch50-L:scaling}  Both \(A_{X,B}\) and \(\phi\) are homogeneous for
the scaling action, the first by
\Cref{prop:ch50-A-properties}\,\cref{itm:ch50-A:homog} and the second by
\eqref{eq:function-scaling-action}; the infimum over the scaling-stable set
\(X^{\divis}\cup\{v_{\triv}\}\) inherits the homogeneity.
\end{proof}

\begin{remark}[\(L^{\NA}\) as a log canonical threshold]
\label{rem:ch50-L-lct}
If \(\phi=m^{-1}\phi_{\mathfrak a}\) is the model function of a flag ideal
\(\mathfrak a\) on \(X_{\triv}\), then the valuative criterion for
multiplier ideals \eqref{eq:appC-valuative-multiplier} identifies
\(L^{\NA}_B(\phi)\) with a log canonical threshold of the pair
\((X\times\mathbf A^1,B_{\mathbf A^1})\) along \(\mathfrak a\), suitably
normalized.  Precisely, \(\inf_v\{A_{X,B}(v)+\phi(v)\}\ge0\) is equivalent
to the sublc-ness of \((X\times\mathbf A^1,\,B_{\mathbf A^1}
+m^{-1}\mathfrak a)\) --- a statement of exactly the shape of
\Cref{thm:appC-valuative-multiplier}.  This is the reason \(D^{\NA}\) is
the functional that appears in the Fano theory: its positivity is a
singularity condition on the total space of the degeneration, and no
scalar curvature is involved.  We do not use this reformulation, and we
record it only because it explains the shape of the definition.
\end{remark}

\subsection{Anticanonical metrics}
\label{sec:ch50-anticanonical}

Following \cite[Definition~7.31]{BHJ17}, call \(\phi\)
\emph{anticanonical} if it has a normal representative
\((\cX,\cL)\) such that \((\cX,B_{\cX}+\cX_{0,\mathrm{red}})\) is sublc and
\[
        \cL=-K^{\log}_{(\cX,B_{\cX})/\mathbf A^1}+c\cX_0
        \quad\text{for some }c\in\mathbf Q.
\]

\begin{lemma}
\label{lem:ch50-anticanonical-constant}
Let \(\phi\) be anticanonical with representative \((\cX,\cL)\) and
constant \(c\).  Then
\[
 A_{X,B}(v_E)+\phi(v_E)=c
\tag{50.64}\label{eq:ch50-anticanonical-components}
\]
for every irreducible component \(E\) of \(\cX_0\), and
\[
 A_{X,B}(v)+\phi(v)\ \ge\ c
 \qquad\text{for every }v\in X^{\divis}\cup\{v_{\triv}\}.
\tag{50.65}\label{eq:ch50-anticanonical-inequality}
\]
Consequently \(L^{\NA}_B(\phi)=c\).
\end{lemma}

\begin{proof}
Write \(\cL=p^*L+D\) with \(D\) a vertical \(\mathbf Q\)-Cartier divisor on
\(\cX\), so that \(\phi(v_E)=b_E^{-1}\ord_E(D)\) by
\Cref{lem:ch33-effectivity} and \eqref{eq:ch33-bE}.  Since
\(L=-K_{X,B}\), formula \eqref{eq:ch50-log-canonical-atoms} gives
\[
 -K^{\log}_{(\cX,B_\cX)/\mathbf A^1}+c\cX_0
 =p^*L-\sum_{E}b_E\,A_{X,B}(v_E)\,E+c\sum_Eb_EE .
\]
Comparing with \(\cL=p^*L+D\) yields
\(\ord_E(D)=b_E\bigl(c-A_{X,B}(v_E)\bigr)\), which is
\eqref{eq:ch50-anticanonical-components}.

For \eqref{eq:ch50-anticanonical-inequality}, let
\(v\in X^{\divis}\cup\{v_{\triv}\}\).  By
\Cref{lem:ch33-gauss-component} there is a normal test configuration
\(\cX'\), which we may take to dominate \(\cX\) by a morphism \(\mu\), and
an irreducible component \(E'\) of \(\cX'_0\) with \(v=v_{E'}\).  Applying
\eqref{eq:ch50-log-canonical-atoms} on \(\cX'\) and pulling back the same
identity from \(\cX\), we obtain
\[
 K^{\log}_{(\cX',B_{\cX'})/\mathbf A^1}
 =\mu^*K^{\log}_{(\cX,B_{\cX})/\mathbf A^1}+\Delta,
 \qquad
 \ord_{E'}(\Delta)=A_{(\cX,\,B_\cX+\cX_{0,\mathrm{red}})}(\ord_{E'}),
\]
the second equality because
\(K^{\log}=K+B_\cX+\cX_{0,\mathrm{red}}\) relative to \(\mathbf A^1\), so
that \(\Delta=K_{\cX'}+B_{\cX'}+\cX'_{0,\mathrm{red}}
-\mu^*(K_{\cX}+B_{\cX}+\cX_{0,\mathrm{red}})\), whose coefficient along
\(E'\) is \(\bigl(A_{(\cX,B_\cX+\cX_{0,\mathrm{red}})}(\ord_{E'})-1\bigr)
+1\); the extra \(1\) is the coefficient of \(E'\) in
\(\cX'_{0,\mathrm{red}}\).  Now
\[
 b_{E'}\phi(v_{E'})=\ord_{E'}\bigl(\mu^*\cL-p'^*L\bigr)
 =\ord_{E'}\bigl(-\mu^*K^{\log}_{(\cX,B_\cX)/\mathbf A^1}\bigr)
   +c\,b_{E'}+\ord_{E'}(p'^*K_{X,B})
\]
and, substituting \(\mu^*K^{\log}_\cX=K^{\log}_{\cX'}-\Delta\) and then
\eqref{eq:ch50-log-canonical-atoms} on \(\cX'\),
\[
 b_{E'}\phi(v_{E'})
 =-b_{E'}A_{X,B}(v_{E'})+c\,b_{E'}+\ord_{E'}(\Delta).
\]
Dividing by \(b_{E'}\) gives
\[
 A_{X,B}(v)+\phi(v)
 =c+b_{E'}^{-1}A_{(\cX,\,B_\cX+\cX_{0,\mathrm{red}})}(\ord_{E'})
 \ \ge\ c,
\tag{50.66}\label{eq:ch50-anticanonical-defect}
\]
the inequality by the sublc hypothesis.  Together with
\eqref{eq:ch50-anticanonical-components} this gives
\(L^{\NA}_B(\phi)=c\).
\end{proof}

\subsection{Ding is dominated by Mabuchi}
\label{sec:ch50-ding-mabuchi}

\begin{theorem}\label{thm:stability-ding-mabuchi}
For every semipositive model metric on \(-K_{X,B}\),
\[
                        D_B^{\NA}(\phi)\leq M_B^{\NA}(\phi).
\tag{50.67}\label{eq:stability-ding-less-M}
\]
Equality holds when \(\phi\) is anticanonical in this precise sense.
Moreover, for every
\(\delta\in[0,1]\),
\[
 M_B^{\NA}\geq\delta J
 \quad\Longleftrightarrow\quad
 D_B^{\NA}\geq\delta J
\tag{50.68}\label{eq:stability-ding-M-uniform}
\]
on the space of positive model metrics.
\end{theorem}

\begin{proof}[Proof of \eqref{eq:stability-ding-less-M} and of the equality
case]
Let \(\phi\) be a semipositive model metric with normal representative
\((\cX,\cL)\), which we may assume dominates \(X\times\mathbf A^1\).  By
\Cref{prop:ch50-model-entropy} the probability measure
\(\MA_\omega(\phi)\) is \(\sum_Ec_E\delta_{v_E}\) with
\(c_E:=V^{-1}b_E(E\cdot\overline\cL^{\,n})\ge0\) and \(\sum_Ec_E=1\).
Hence
\[
 \begin{split}
 H^{\NA}_B(\phi)+\int_{X^{\An}}\phi\,\MA_\omega(\phi)
 &=\sum_Ec_E\Bigl(A_{X,B}(v_E)+\phi(v_E)\Bigr)\\
 &\ge\ \min_E\Bigl(A_{X,B}(v_E)+\phi(v_E)\Bigr)
 \ \ge\ L^{\NA}_B(\phi),
 \end{split}
\tag{50.69}\label{eq:ch50-ding-convex-combination}
\]
the first inequality because a convex combination is at least the minimum,
the second because every \(v_E\) lies in \(X^{\divis}\).  On the other
hand, \(K_{X,B}=-\omega\), so \Cref{prop:ch50-M-proportional} with
\(\gamma=-1\) and \eqref{eq:ch50-min-norm-identities} give
\[
 M^{\NA}_B(\phi)
 =H^{\NA}_B(\phi)-E^\vee_\omega\bigl(\MA_\omega(\phi)\bigr)
 =H^{\NA}_B(\phi)+\int\phi\,\MA_\omega(\phi)-E_\omega(\phi).
\tag{50.70}\label{eq:ch50-M-fano-form}
\]
Combining, \(M^{\NA}_B(\phi)\ge L^{\NA}_B(\phi)-E_\omega(\phi)
=D^{\NA}_B(\phi)\), which is \eqref{eq:stability-ding-less-M}.

If \(\phi\) is anticanonical with constant \(c\), then by
\Cref{lem:ch50-anticanonical-constant} every bracket in
\eqref{eq:ch50-ding-convex-combination} equals \(c\) and
\(L^{\NA}_B(\phi)=c\); so both inequalities in
\eqref{eq:ch50-ding-convex-combination} are equalities, and
\(M^{\NA}_B(\phi)=c-E_\omega(\phi)=D^{\NA}_B(\phi)\).
\end{proof}

\begin{proof}[Imported deep step: the uniform equivalence]
The implication
\[
 M^{\NA}_B\ge\delta J\quad\Longrightarrow\quad D^{\NA}_B\ge\delta J
\]
is the deep one, and it is what we import.  The
opposite implication, from \(D^{\NA}_B\ge\delta J\) to
\(M^{\NA}_B\ge\delta J\), is immediate from
\eqref{eq:stability-ding-less-M}, which has just been proved: \(M\ge D\ge
\delta J\) pointwise.  Nothing needs to be imported for it.  For the deep
direction the mechanism is a
reduction to \emph{special} test configurations by the minimal model
program: given \((\cX,\cL)\), one produces from it an anticanonical
representative with no larger \(\DF\) and comparable \(J\), on which
\(D=M\) by the equality case above.  This is
\cite[Proposition~7.32 and Theorem~8.6]{BHJ17} (preprint 2015, published
2017); the minimal model program input is \cite{BCHM10}.  The assertion is
not a formal consequence of pluripotential theory alone, and no proof of it
is given in this book.
\end{proof}

\begin{remark}[The two halves of \Cref{thm:stability-ding-mabuchi}]
\label{rem:ch50-ding-halves}
The inequality \eqref{eq:stability-ding-less-M} and its equality case are
proved above from the model formula for the Monge--Amp\`ere measure and
from the log-canonical bookkeeping of
\eqref{eq:ch50-log-canonical-atoms}; nothing beyond
\Cref{chap:bj-energy} is used, and the argument is characteristic-free
apart from the standing hypothesis that log discrepancies make sense.  The
equivalence \eqref{eq:stability-ding-M-uniform} is of a different nature:
it is a statement about the \emph{infima} of two functionals over a class
of degenerations, and it is proved by replacing an arbitrary degeneration
with a special one.  The two halves should not be quoted as a single
theorem.  \Cref{ex:ch50-p1-dnc} exhibits a one-parameter family of
positive model metrics on \((\mathbf P^1,-K)\) along which
\(M^{\NA}-D^{\NA}=\tfrac12 c(1-c)\) for \(c\in[0,1]\) and
\(-\tfrac12(c-1)(c-2)\) for \(c\in[1,2]\), vanishing exactly at the
anticanonical member \(c=1\) and at the boundary \(c=2\); it is a complete
illustration of both halves.
\end{remark}

\begin{remark}[Why \(D^{\NA}\) is not defined for all filtrations here]
\label{rem:ch50-ding-scope}
The Mabuchi functional was extended to all of \(\mathcal N_{\mathbf R}\) by
\eqref{eq:stability-M-all-norms}, because
\eqref{eq:stability-beta-measure} is a functional of the measure.  The
Ding functional is not a functional of \(\MA(\phi)\): it is an infimum of
\(\phi\) itself, shifted by \(A_{X,B}\).  Extending it to a general
filtration therefore requires extending the \emph{potential}, and by
\Cref{thm:stability-norm-completion} the potential \(\FS(\chi)\) of a
general norm need not lie in \(\cE^1(\omega)\).  For norms whose
Fubini--Study function is bounded --- which includes every linearly bounded
filtration, by \eqref{eq:ch50-fs-operator} --- the definition
\eqref{eq:stability-Ding} does make sense verbatim, and
\eqref{eq:stability-ding-less-M} continues to hold whenever
\(\FS(\chi)\in\cE^1(\omega)\) and \(\MA(\chi)=\MA_\omega(\FS(\chi))\); we
do not pursue this, because the results of \Cref{sec:ch50-uniform} are
stated for \(M\), and the Ding functional enters only through
\eqref{eq:stability-ding-M-uniform} on model metrics.
\end{remark}

\section{Uniform K-stability, divisorial stability, filtration stability}
\label{sec:ch50-uniform}

\subsection{The classical definition and its model form}
\label{sec:ch50-classical}

For a normal ample test configuration let
\(\|\cX,\cL\|_{\min}\) denote
\eqref{eq:stability-minimum-norm}.  The polarized pair \((X,B;L)\) is
uniformly K-stable if there exists \(\varepsilon>0\) such that
\[
       \DF_B(\cX,\cL)
            \geq\varepsilon\|\cX,\cL\|_{\min}
\tag{50.71}\label{eq:stability-uniform-K}
\]
for every normal ample test configuration.  By
\eqref{eq:stability-M-DF-correction}, normalized base change, and
\eqref{eq:stability-J-L1-comparison}, this is equivalent to
\[
                     M_B^{\NA}(\phi)\geq\varepsilon'J(\phi)
\tag{50.72}\label{eq:stability-M-coercive-model}
\]
on positive model metrics, for some \(\varepsilon'>0\)
\cite[Proposition~8.2]{BHJ17}.

\begin{definition}[The three classical conditions]
\label{def:ch50-k-stability}
\index{K-stability!semistable}\index{K-stability!uniform}
Let \((X,B;L)\) be as in \Cref{sec:ch50-hypotheses}, with \((X,B)\) sublc.
\begin{enumerate}
\item\label{itm:ch50-ks:semi}
      \((X,B;L)\) is \emph{K-semistable} if \(\DF_B(\cX,\cL)\ge0\) for
      every normal ample test configuration \((\cX,\cL)\);
\item\label{itm:ch50-ks:stable}
      it is \emph{K-stable} if moreover \(\DF_B(\cX,\cL)=0\) only for
      \(\|\cX,\cL\|_{\min}=0\);
\item\label{itm:ch50-ks:uniform}
      it is \emph{uniformly K-stable} if \eqref{eq:stability-uniform-K}
      holds for some \(\varepsilon>0\).
\end{enumerate}
\end{definition}

\begin{proposition}
\label{prop:ch50-model-form}
With the notation of \Cref{def:ch50-k-stability}, and with
\(\phi\) running over positive model metrics on \(L\):
\begin{enumerate}
\item\label{itm:ch50-mf:semi}
      \((X,B;L)\) is K-semistable if and only if \(M^{\NA}_B(\phi)\ge0\)
      for every positive model metric \(\phi\);
\item\label{itm:ch50-mf:uniform}
      \((X,B;L)\) is uniformly K-stable if and only if
      \eqref{eq:stability-M-coercive-model} holds for some
      \(\varepsilon'>0\), if and only if
      \(M^{\NA}_B(\phi)\ge\varepsilon''\|\chi\|_{\min}\) for some
      \(\varepsilon''>0\), if and only if
      \(M^{\NA}_B(\phi)\ge\varepsilon'''\|\chi\|_1\) for some
      \(\varepsilon'''>0\).
\end{enumerate}
The constants are related by factors depending only on \(n\).
\end{proposition}

\begin{proof}
\emph{From \(\DF\) to \(M^{\NA}\).}  By
\eqref{eq:stability-M-DF-correction}, \(M^{\NA}_B\le\DF_B\), so
K-semistability does not immediately give \(M^{\NA}_B\ge0\).  Use base
change: for \(d\) sufficiently divisible the normalized base change
\((\cX_d,\cL_d)\) has reduced central fibre and
\(M^{\NA}_B(d\cdot\phi)=\DF_B(\cX_d,\cL_d)\), by Step 5 of the proof of
\Cref{thm:stability-chen-tian}.  Since \(M^{\NA}_B\) is homogeneous of
degree one for the scaling action --- by
\Cref{prop:ch50-M-proportional} in the proportional case, and in general
because every term of \eqref{eq:stability-chen-tian} is, by
\Cref{thm:ch46-energy-homogeneous},
\Cref{prop:ch50-A-properties}\,\cref{itm:ch50-A:homog} and
\eqref{eq:ch33-intersection-base-change} --- we get
\(dM^{\NA}_B(\phi)=\DF_B(\cX_d,\cL_d)\ge0\), whence
\(M^{\NA}_B(\phi)\ge0\).  Conversely \(\DF_B\ge M^{\NA}_B\ge0\).

\emph{The uniform statement.}  The same base change multiplies
\(\|\chi\|_{\min}\) by \(d\), by
\Cref{prop:ch50-minimum-norm}\,\cref{itm:ch50-mn:scaling}; so
\eqref{eq:stability-uniform-K} for all normal ample test configurations is
equivalent to \(M^{\NA}_B(\phi)\ge\varepsilon\|\chi\|_{\min}\) for all
positive model metrics.  The three remaining reformulations follow from the
comparisons of \Cref{rem:ch50-which-functional}: \(\|\chi\|_{\min}\),
\(J\) and \(\|\chi\|_1\) are pairwise comparable with constants depending
only on \(n\), by \eqref{eq:ch50-min-vs-J} and
\eqref{eq:stability-J-L1-comparison}.  This is
\cite[Proposition~8.2]{BHJ17} in the form in which we use it; the proof
just given is complete given \Cref{thm:stability-chen-tian}.
\end{proof}

\begin{remark}[K-stability sits strictly between]
\label{rem:ch50-k-stable-between}
Uniform K-stability implies K-stability implies K-semistability, and both
implications are strict in general.  The first is strict because
\(\DF>0\) on a family with \(\DF/\|\cdot\|_{\min}\to0\) is possible;
the second is strict as soon as \(\Aut(X,L)\) is positive dimensional,
since a product configuration then has \(\|\cdot\|_{\min}>0\) and
\(\DF=0\) whenever the Futaki character vanishes ---
\Cref{prop:ch50-p1-product} is the smallest instance, and
\Cref{ex:ch50-semistable-not-uniform} works it out.  In this book we use
only K-semistability and uniform K-stability; K-stability is recorded for
completeness and is not used, because it is not a metric condition and
does not interact with the completions of
\Cref{thm:stability-norm-completion}.
\end{remark}

\subsection{Divisorial measures and divisorial norms}
\label{sec:ch50-divisorial}

The completed theory has a stronger testing class.  A divisorial measure is
a probability measure
\[
             \mu=\sum_{i=1}^r a_i\delta_{v_i},
\qquad
 v_i\in X^{\divis},\quad a_i>0,\quad\sum_i a_i=1.
\tag{50.73}\label{eq:stability-divisorial-measure}
\]
Write \(\mathcal M^{\divis}\) for their space.  For
\(v\in X^{\divis}\), let \(\chi_v(s):=v(s)\).  A \emph{divisorial norm}
is one of the form
\[
       \chi=\min_{1\leq i\leq r}\{\chi_{v_i}+c_i\},
       \qquad
       \chi_m(s)=\min_i\{v_i(s)+mc_i\};
\]
their space is denoted \(\mathcal N_{\mathbf R}^{\divis}\).  Define
\[
 \sigma_{\div}(X,B;\omega)
 :=\inf_{\substack{\mu\in\mathcal M^{\divis}\\
                    \mu\ne\delta_{v_{\triv}}}}
       \frac{\beta_{X,B;\omega}(\mu)}{\|\mu\|_\omega}.
\tag{50.74}\label{eq:stability-divisorial-threshold}
\]
Divisorial semistability means \(\sigma_{\div}\geq0\), and divisorial
stability means \(\sigma_{\div}>0\).

\begin{lemma}
\label{lem:ch50-divisorial-norms}
Let \(v_1,\dots,v_r\in X^{\divis}\) and \(c_1,\dots,c_r\in\mathbf R\), and
let \(\chi=\min_i\{\chi_{v_i}+c_i\}\).  Then:
\begin{enumerate}
\item\label{itm:ch50-dn:norm}
      \(\chi\in\mathcal N_{\mathbf R}\), it is homogeneous, and
      \(\chi=\chi^{\hom}\);
\item\label{itm:ch50-dn:fs}
      \(\FS(\chi)=\max_i\{-\phi_{v_i}\ \text{shifted}\}\) is a bounded
      \(\omega\)-psh function; more precisely
      \(\FS(\chi)(w)=\sup_{m,s}m^{-1}(\chi_m(s)-w(s))\) is finite for every
      \(w\in X^{\An}\) and \(\lambda_{\min}\le\FS(\chi)\le\lambda_{\max}\);
\item\label{itm:ch50-dn:trivial}
      for \(r=1\), \(v_1=v\) and \(c_1=0\), the measure
      \(\MA(\chi_v)\) is \(\delta_v\), and
      \(\|\chi_v\|_{\min}=E^\vee_\omega(\delta_v)\).
\end{enumerate}
\end{lemma}

\begin{proof}
\cref{itm:ch50-dn:norm}  Each \(\chi_{v_i}+c_i\) is a norm: the ultrametric
inequality is \(v_i(s+s')\ge\min\{v_i(s),v_i(s')\}\), multiplicativity is
\(v_i(ss')=v_i(s)+v_i(s')\) --- an equality, so
\(\chi_{v_i}\) is homogeneous --- and linear boundedness is
\(v_i(s)\le mT_L(v_i)\) with \(T_L(v_i)<\infty\) by
\Cref{prop:ch49-psef-threshold}.  A minimum of finitely many norms is a
norm, by \Cref{lem:ch50-lattice}, and a minimum of homogeneous norms is
homogeneous because \(\chi_{rm}(s^r)=\min_i\{rv_i(s)+rmc_i\}
=r\chi_m(s)\).

\cref{itm:ch50-dn:fs}  Finiteness and the bounds follow from linear
boundedness as in \Cref{rem:ch50-energy-beyond-finite-type}; pshness is the
increasing-supremum description \eqref{eq:ch33-phi-sup-of-flags}.

\cref{itm:ch50-dn:trivial}  For \(r=1\) the function \(\FS(\chi_v)\) is,
by \eqref{eq:ch50-fs-operator}, the supremum over \(m,s\) of
\(m^{-1}(v(s)-w(s))\), which is the Green function of \(v\) normalized by
\(\FS(\chi_v)(v)=0\); this is
\Cref{ex:ch47-dirac-green} and, on the toric slice,
\Cref{prop:ch49-toric-green}\,\cref{itm:ch49-tg:ma}.  Its Monge--Amp\`ere
measure is \(\delta_v\), and \eqref{eq:ch50-min-norm-identities} gives the
last assertion.
\end{proof}

\begin{theorem}\label{thm:stability-divisorial-filtrations}
Assume that \((X,B)\) is subklt.  Then
\[
 \sigma_{\div}(X,B;L)
 =\inf_{\substack{\chi\in\mathcal N_{\mathbf R}\\
                   \|\chi\|_{\min}>0}}
       \frac{M_{X,B;L}(\chi)}{\|\chi\|_{\min}}.
\tag{50.75}\label{eq:stability-filtration-threshold}
\]
Consequently, divisorial semistability is equivalent to
\(M(\chi)\geq0\) for every linearly bounded multiplicative filtration, and
divisorial stability is equivalent to
\[
                M(\chi)\geq\varepsilon\|\chi\|_{\min}
                \quad\text{for every }\chi\in\mathcal N_{\mathbf R}.
\tag{50.76}\label{eq:stability-all-filtrations}
\]
\end{theorem}

\begin{proof}[Proof architecture]
The Monge--Amp\`ere map gives an isometric bijection
\[
       \MA:\mathcal N_{\mathbf R}^{\divis}/\mathbf R
                  \xrightarrow{\sim}\mathcal M^{\divis}
\tag{50.77}\label{eq:stability-divnorm-divmeasure}
\]
between divisorial norms and divisorial measures
\cite[Theorem~C]{BJKS1}.  Strong lower semicontinuity of entropy
approximates any finite-energy measure by divisorial measures without
losing the value of \(\beta\).  Apply
\eqref{eq:stability-norm-ma-isometry} and homogeneity.  This is
\cite[Theorem~5.8]{BJKS2}.
\end{proof}

\begin{remark}[The three inputs of \Cref{thm:stability-divisorial-filtrations},
separated]
\label{rem:ch50-divisorial-inputs}
The proof architecture above compresses three genuinely different
ingredients, and it is worth naming them.
\begin{enumerate}
\item\label{itm:ch50-di:bijection}
      \emph{The divisorial bijection} \eqref{eq:stability-divnorm-divmeasure}.
      That \(\MA\) sends \(\mathcal N^{\divis}_{\mathbf R}/\mathbf R\)
      \emph{into} \(\mathcal M^{\divis}\) is elementary: a divisorial norm
      is a finite minimum, its Fubini--Study function is a finite maximum
      of Green functions, and the Monge--Amp\`ere measure of such a
      function is supported on the finitely many \(v_i\), by
      \Cref{lem:ch50-divisorial-norms} and locality
      \Cref{thm:bj-locality-comparison}.  That the map is
      \emph{bijective and isometric} is
      \cite[Theorem~C]{BJKS1} (preprint 2021, published 2025), and it is
      imported.
\item\label{itm:ch50-di:density}
      \emph{Strong density of divisorial measures} in
      \(\cM^1(\omega)\), which in characteristic zero is
      \Cref{thm:bj-strong-divisorial-density} combined with
      \Cref{thm:bj-linear-topology-complete}\,\cref{itm:ch49-top:embedding}
      and convex combination; this part is proved in this book.
\item\label{itm:ch50-di:lsc}
      \emph{Strong lower semicontinuity of \(\beta\), and the fact that
      it is the maximal such extension from divisorial measures}, which
      is \Cref{imp:ch50-strong-lsc}\,\cref{itm:ch50-slsc:extension},
      imported from \cite[\S\S3.4 and 4.1]{BJKS2} (preprint June 2022,
      published 2023).  This is the step that fails for the \emph{weak}
      topology, by \Cref{thm:ch50-entropy-lsc} and
      \Cref{rem:ch50-entropy-lsc-direction}, and it is why
      \eqref{eq:stability-filtration-threshold} is not a formal
      consequence of density.
\end{enumerate}
Only \cref{itm:ch50-di:density} is book-internal.  We do not reprove
\cref{itm:ch50-di:bijection} or \cref{itm:ch50-di:lsc}.
\end{remark}

\subsection{The three thresholds}
\label{sec:ch50-thresholds}

Write
\[
 \sigma_K(X,B;\omega)
 :=\inf_{\substack{\chi\in\mathcal T_{\mathbf R}\\
        \|\chi\|_{\min}>0}}
      \frac{M^{\NA}_B(\FS(\chi))}{\|\chi\|_{\min}},
 \qquad
 \sigma_{\mathcal N}(X,B;\omega)
 :=\inf_{\substack{\chi\in\mathcal N_{\mathbf R}\\
        \|\chi\|_{\min}>0}}
      \frac{M_{X,B;L}(\chi)}{\|\chi\|_{\min}} ,
\tag{50.78}\label{eq:ch50-three-thresholds}
\]
the \emph{K-stability threshold} and the \emph{filtration threshold}, and
recall \(\sigma_{\div}\) from
\eqref{eq:stability-divisorial-threshold}.  Then
\((X,B;L)\) is K-semistable if and only if \(\sigma_K\ge0\), and uniformly
K-stable if and only if \(\sigma_K>0\), by
\Cref{prop:ch50-model-form}.

\begin{proposition}
\label{prop:ch50-threshold-comparisons}
Assume \((X,B)\) subklt.  Then
\[
 \sigma_K\ \ge\ \sigma_{\mathcal N}\ =\ \sigma_{\div},
\tag{50.79}\label{eq:ch50-threshold-chain}
\]
the inequality because \(\mathcal T_{\mathbf R}\subset\mathcal
N_{\mathbf R}\) and the two Mabuchi functionals agree on
\(\mathcal T_{\mathbf R}\), and the equality being
\eqref{eq:stability-filtration-threshold}.  Moreover all three are
homogeneous of degree zero in \(\omega\) and unchanged by replacing
\(L\) by a positive multiple.
\end{proposition}

\begin{proof}
For \(\chi\in\mathcal T_{\mathbf R}\), the measure
\(\MA(\chi)\) of \eqref{eq:stability-norm-ma-isometry} is
\(\MA_\omega(\FS(\chi))\) by \Cref{prop:ch46-pl-measure}, and
\(M_{X,B;L}(\chi)=\beta_{X,B;\omega}(\MA_\omega(\FS(\chi)))
=M^{\NA}_B(\FS(\chi))\) by \eqref{eq:stability-M-beta}.  Hence the
infimum defining \(\sigma_{\mathcal N}\) is over a larger set of the same
quotients, giving \(\sigma_{\mathcal N}\le\sigma_K\).  The equality
\(\sigma_{\mathcal N}=\sigma_{\div}\) is
\Cref{thm:stability-divisorial-filtrations}.  Homogeneity: replacing
\(L\) by \(rL\) multiplies both \(M\) and \(\|\cdot\|_{\min}\) by \(r\), by
\Cref{thm:ch46-energy-homogeneous} and the corresponding scaling of the
entropy, so the ratio is unchanged.
\end{proof}

\begin{remark}[Ordinary versus filtration stability]
\label{rem:stability-gap-general-polarization}
For a general polarization, \Cref{thm:stability-divisorial-filtrations}
does not identify stability for all filtrations with ordinary uniform
K-stability, which tests only finite-type norms.  If
\(\sigma_K\) is the usual threshold, then
\[
                         \sigma_K\geq\sigma_{\div}.
\tag{50.80}\label{eq:stability-threshold-inequality}
\]
Equality was Conjecture~5.15 of the 2023 source \cite{BJKS2}; it amounts to
an entropy-regularization statement for canonical finite-type
approximants.  Thus no equality in general polarization is silently used
here.  Later developments, when invoked in
\Crefrange{chap:master-dictionary}{chap:major-open-problems}, will be dated
and stated with their own hypotheses.
\end{remark}

\begin{remark}[Where the gap comes from]
\label{rem:ch50-gap-mechanism}
The inequality \eqref{eq:stability-threshold-inequality} can only be strict
if there is a filtration \(\chi\in\mathcal N_{\mathbf R}\) with
\(M_{X,B;L}(\chi)<\sigma_K\|\chi\|_{\min}\) that is not approximable by
finite-type norms without losing entropy.  By
\Cref{imp:ch50-strong-lsc}\,\cref{itm:ch50-slsc:extension} the value of
\(\beta\) \emph{is} recovered by strong approximation with
\emph{divisorial} measures; the difficulty is that a divisorial measure
need not be the Monge--Amp\`ere measure of a \emph{finite-type} norm ---
by \eqref{eq:stability-divnorm-divmeasure} it is the measure of a
divisorial norm, and a divisorial norm is finitely generated only under
extra hypotheses.  \Cref{ex:ch50-harmonic} is a filtration that is not of
finite type, with everything computable; it does not settle the question,
because it is not destabilizing.  In the log Fano case the gap is closed by
the results imported in \Cref{sec:ch50-valuative}, and in particular by the
finite-generation theorem of \cite{LXZ22} (preprint 2021, published 2022);
that closure is Fano-specific and must not be transplanted.
\end{remark}

\subsection{Openness and continuity: an import, dated}
\label{sec:ch50-openness}

\begin{important}{Imported: local H\"older continuity and openness}
\label{imp:ch50-openness}
Let \((X,B)\) be subklt.  Then
\(\omega\mapsto\sigma_{\div}(X,B;\omega)\) is locally H\"older continuous
on the ample cone \(\Amp(X)\subset N^1(X)_{\mathbf R}\)
\cite[Theorem~4.12]{BJKS2} (preprint June 2022, published 2023).  In
particular the divisorial-stability locus
\(\{\omega:\sigma_{\div}(X,B;\omega)>0\}\) is open in the ample cone.

Two cautions.  First, the statement is about \(\sigma_{\div}\) and hence,
by \Cref{prop:ch50-threshold-comparisons}, about \(\sigma_{\mathcal N}\);
it is \emph{not} a statement about \(\sigma_K\), which is only known to
dominate it.  Second, the proof rests on the \(C^1\)-regularity of
\(\omega\mapsto\|\mu\|_\omega\) recorded in
\Cref{imp:ch50-strong-lsc}\,\cref{itm:ch50-slsc:diff}, so it is not
available from the material of this book alone.  Openness of uniform
K-stability for a \emph{fixed} polarization under deformation of \((X,B)\)
is a different and much harder statement, is not used here, and is not
asserted.
\end{important}

\section{The valuative criterion in the log Fano case}
\label{sec:ch50-valuative}

\subsection{The valuative invariants}
\label{sec:ch50-valuative-invariants}

For a divisorial valuation \(v\), define
\[
 \begin{split}
 T_L(v)&:=\sup\{t\geq0\mid\vol(L-tv)>0\},\\
 S_L(v)&:=\frac1V\int_0^\infty\vol(L-tv)\,\rd t,\\
 \beta_L(v)&:=A_{X,B}(v)-S_L(v),\\
 j_L(v)&:=T_L(v)-S_L(v).
 \end{split}
\tag{50.81}\label{eq:stability-valuative-invariants}
\]
Here \(\vol(L-tv)\) means the volume of the filtered linear series
\(\{s\in H^0(mL)\mid v(s)\geq mt\}\), so the definition does not require
the centre of \(v\) to be a Cartier divisor on \(X\).  Both \(A\) and \(S\)
are homogeneous in \(v\), and
\[
                       S_L(v)=E^\vee(\delta_v)
\]
by \cite[Theorem~7.22(iii)]{BJKS1}.  This equality sharpens the
two-sided comparison with \(T_L(v)\) in
\Cref{thm:bj-dirac-finite-energy}.

\begin{proposition}
\label{prop:ch50-valuative-dh}
Let \(v\in X^{\lin}\) and let \(\chi_v\) be the divisorial norm
\(\chi_{v,m}(s)=v(s)\) of \Cref{lem:ch50-divisorial-norms}.  Then:
\begin{enumerate}
\item\label{itm:ch50-vdh:T}
      \(T_L(v)=\lambda_{\max}(\chi_v)=T_\omega(v)\), where the middle term
      is \eqref{eq:ch50-lambda-max}; in particular
      \(T_L(v)<\infty\) exactly when \(v\in X^{\lin}\), by
      \Cref{prop:ch49-T-equals-T};
\item\label{itm:ch50-vdh:S}
      \(S_L(v)=\int\lambda\,\rd\DHm_{\chi_v}(\lambda)
      =\operatorname{vol}(\chi_v)\);
\item\label{itm:ch50-vdh:j}
      \(j_L(v)=J(\FS(\chi_v))\), and
      \(\|\chi_v\|_{\min}=E^\vee_\omega(\delta_v)\);
\item\label{itm:ch50-vdh:comparison}
      \(\displaystyle\frac{T_L(v)}{n+1}\le S_L(v)\le\frac{n\,T_L(v)}{n+1}\)
      and hence
      \(\displaystyle\frac{T_L(v)}{n+1}\le j_L(v)\le\frac{n\,T_L(v)}{n+1}\),
      so that
      \[
       \frac1n\,S_L(v)\ \le\ j_L(v)\ \le\ n\,S_L(v).
      \tag{50.82}\label{eq:ch50-j-vs-S}
      \]
\end{enumerate}
\end{proposition}

\begin{proof}
\cref{itm:ch50-vdh:T}  By definition \(\lambda_{\max}(\chi_v)
=\lim_m m^{-1}\max_{0\ne s\in R_m}v(s)=T_L(v)\) in the sense of
\eqref{eq:bj-maximal-vanishing}, and
\(\vol(L-tv)>0\) exactly for \(t<\lim_mm^{-1}\max_sv(s)\), the volume of a
filtered piece being positive exactly when the piece is nonzero for
\(m\gg0\).  The identification with \(T_\omega\) is
\Cref{prop:ch49-T-equals-T}.

\cref{itm:ch50-vdh:S}  Both sides are \(\int_0^\infty\) of the same
decreasing function.  Indeed, with \(\mathcal F^{mt}R_m
=\{s:v(s)\ge mt\}\), the Duistermaat--Heckman measure of \(\chi_v\)
satisfies \(\DHm_{\chi_v}([t,\infty))=V^{-1}\vol(L-tv)\) by
\Cref{cor:ch33-dh-structure}\,(1) --- which gives the open tail, the two
agreeing for all but the countably many \(t\) carrying an atom, which
suffices for the integration below --- and \eqref{eq:ch33-okounkov-volume}; and
for a probability measure supported in \([0,\infty)\) one has
\(\int\lambda\,\rd\DHm=\int_0^\infty\DHm([t,\infty))\rd t\).  This gives
\(\int\lambda\,\rd\DHm_{\chi_v}=V^{-1}\int_0^\infty\vol(L-tv)\rd t
=S_L(v)\).  The identification with \(\operatorname{vol}(\chi_v)\) is
\eqref{eq:stability-volume-norm}.

\cref{itm:ch50-vdh:j}  By \eqref{eq:stability-dh-energy-J} and
\Cref{prop:ch50-sup-lambdamax},
\(J(\FS(\chi_v))=\lambda_{\max}(\chi_v)-\operatorname{vol}(\chi_v)
=T_L(v)-S_L(v)=j_L(v)\); the last assertion is
\Cref{lem:ch50-divisorial-norms}\,\cref{itm:ch50-dn:trivial}.

\cref{itm:ch50-vdh:comparison}  By
\Cref{lem:ch50-divisorial-norms}\,\cref{itm:ch50-dn:trivial} and
\eqref{eq:stability-minimum-norm}, \(\|\chi_v\|_{\min}
=E^\vee_\omega(\delta_v)\), and by \eqref{eq:bj-dirac-energy-T},
\[
 \frac{T_\omega(v)}{n+1}\ \le\ E^\vee_\omega(\delta_v)
 \ \le\ \frac{n\,T_\omega(v)}{n+1}.
\]
The imported identity \(S_L(v)=E^\vee_\omega(\delta_v)\) of
\Cref{imp:ch50-imported}\,\cref{itm:ch50-import:S} turns this into the
first chain; subtracting from \(T_L(v)\) gives the second.  Dividing the
two chains gives \eqref{eq:ch50-j-vs-S}.
\end{proof}

\begin{remark}[The one imported identity, and how it is checked]
\label{rem:ch50-S-import}
Only \(S_L(v)=E^\vee_\omega(\delta_v)\) is imported in
\Cref{prop:ch50-valuative-dh}; everything else is proved from
\Cref{chap:bj-test-configurations,chap:bj-energy,chap:bj-linear-growth}.
The identity is verified independently at three places in this book: on
\(\mathbf P^1\) in \Cref{prop:ch50-p1-product} and
\Cref{ex:stability-p1-delta}; on the toric slice, where
\Cref{prop:ch49-toric-green}\,\cref{itm:ch49-tg:energy} gives
\(E^\vee_\omega(\delta_{v_u})=f_P(u)+\langle u,b_P\rangle\) and
\Cref{ex:ch50-toric-F1} matches it against a direct volume integral; and
on \(\mathbf P^2\) in \Cref{ex:ch50-infinite-entropy}, where both sides are
computed from a pseudoeffective threshold.  We also note the weaker,
book-internal substitute: by
\Cref{prop:ch50-valuative-dh}\,\cref{itm:ch50-vdh:S} the quantity
\(S_L(v)\) is \(\operatorname{vol}(\chi_v)\), and by
\eqref{eq:stability-d1-volume-energy} that is \(E_\omega(\FS(\chi_v))\) when
\(\chi_v\) is of finite type; since \(\FS(\chi_v)\) is the Green function
of \(v\), normalized to vanish at \(v\), the identity
\(E_\omega(\phi_v)=E^\vee_\omega(\delta_v)\) is then
\Cref{cor:ch49-green-invariants}.  So for divisorial \(v\) whose Green
function is \(\mathbf Q\)-PL the identity is book-internal; by
\Cref{warn:ch49-green-not-pl} that is not every divisorial \(v\), and the
import is needed in general.
\end{remark}

\subsection{The stability threshold}
\label{sec:ch50-delta}

\begin{theorem}\label{thm:stability-valuative-fano}
Let \((X,B)\) be log Fano and \(L=-K_{X,B}\).  Then:
\begin{enumerate}
\item \((X,B)\) is K-semistable if and only if
      \(A_{X,B}(v)\geq S_L(v)\) for every divisorial \(v\);
\item it is uniformly K-stable if and only if there exists
      \(\varepsilon>0\) such that
      \(\beta_L(v)\geq\varepsilon j_L(v)\) for every divisorial \(v\);
\item with
      \[
       \delta(X,B):=
       \inf_{\substack{v\in X^{\divis}\\v\ne v_{\triv}}}
              \frac{A_{X,B}(v)}{S_L(v)},
      \tag{50.83}\label{eq:stability-delta}
      \]
      K-semistability is equivalent to \(\delta(X,B)\geq1\), and uniform
      K-stability is equivalent to \(\delta(X,B)>1\).
\end{enumerate}
\par
\end{theorem}

\begin{proof}[Provenance]
Dreamy divisorial valuations yield special test configurations, on which
the intersection-theoretic Donaldson--Futaki invariant becomes the
valuative \(\beta\)-invariant.  The minimal model program reduces arbitrary
destabilizing configurations to special ones.  Approximation removes
dreaminess.  The \(\beta\)- and \(j\)-form is
\cite{Fuj19,Li17}; the threshold formulation and the equality with the
expected vanishing order are organized in \cite{BJ20}.  This reduction is
special to the anticanonical polarization and must not be transplanted to
an arbitrary \(L\).
\end{proof}

The internal equivalences among (1), (2) and (3) --- that is, everything
except the link with K-stability itself --- are elementary given
\Cref{prop:ch50-valuative-dh}, and we prove them, so that the imported
content is isolated.

\begin{proposition}
\label{prop:ch50-valuative-internal}
Let \((X,B)\) be log Fano, \(L=-K_{X,B}\), and let \(\delta(X,B)\) be as in
\eqref{eq:stability-delta}.  Then the following are equivalent, with
explicit constants:
\begin{enumerate}
\item\label{itm:ch50-vi:beta}
      \(\beta_L(v)\ge0\) for every \(v\in X^{\divis}\);
\item\label{itm:ch50-vi:delta1}
      \(\delta(X,B)\ge1\);
\item\label{itm:ch50-vi:beta-dirac}
      \(\beta_{X,B;\omega}(\delta_v)\ge0\) for every \(v\in X^{\divis}\).
\end{enumerate}
Likewise the following are equivalent:
\begin{enumerate}[resume]
\item\label{itm:ch50-vi:betaj}
      \(\beta_L\ge\varepsilon j_L\) on \(X^{\divis}\) for some
      \(\varepsilon>0\);
\item\label{itm:ch50-vi:deltabig}
      \(\delta(X,B)>1\);
\item\label{itm:ch50-vi:sigmadiv}
      \(\sigma_{\div}(X,B;\omega)>0\).
\end{enumerate}
Explicitly, \(\delta=1+\eta\) with \(\eta>0\) gives
\(\beta_L\ge\frac\eta nj_L\), and \(\beta_L\ge\varepsilon j_L\) gives
\(\delta\ge1+\frac\varepsilon n\); and
\(\sigma_{\div}\ge\frac{\eta}{1+\eta}\) while
\(\delta\ge1+\sigma_{\div}\).
\end{proposition}

\begin{proof}
\cref{itm:ch50-vi:beta}\(\Leftrightarrow\)\cref{itm:ch50-vi:delta1}
is the definition, \(S_L(v)>0\) for \(v\ne v_{\triv}\) by
\Cref{prop:ch50-valuative-dh}\,\cref{itm:ch50-vdh:comparison} and
\(T_\omega(v)>0\).

\cref{itm:ch50-vi:beta}\(\Leftrightarrow\)\cref{itm:ch50-vi:beta-dirac} is
\Cref{cor:ch50-M-dirac} with \(\gamma=-1\):
\(\beta_{X,B;\omega}(\delta_v)=A_{X,B}(v)-E^\vee_\omega(\delta_v)
=A_{X,B}(v)-S_L(v)=\beta_L(v)\).

\cref{itm:ch50-vi:deltabig}\(\Rightarrow\)\cref{itm:ch50-vi:betaj}.  If
\(A\ge(1+\eta)S\) then \(\beta_L=A-S\ge\eta S\ge\frac\eta nj_L\) by
\eqref{eq:ch50-j-vs-S}.

\cref{itm:ch50-vi:betaj}\(\Rightarrow\)\cref{itm:ch50-vi:deltabig}.  If
\(\beta_L\ge\varepsilon j_L\) then, again by \eqref{eq:ch50-j-vs-S},
\(A-S\ge\varepsilon j_L\ge\frac\varepsilon nS\), so
\(A/S\ge1+\varepsilon/n\).

\cref{itm:ch50-vi:beta-dirac}, uniform version,
\(\Leftrightarrow\)\cref{itm:ch50-vi:sigmadiv}.  By
\Cref{thm:ch50-entropy-lsc} and
\Cref{prop:ch50-M-proportional}, \(\beta_{X,B;\omega}\) is affine along
convex combinations of measures in the first summand and homogeneous of
degree one in the second; hence for a divisorial measure
\(\mu=\sum_ia_i\delta_{v_i}\),
\[
 \beta_{X,B;\omega}(\mu)=\sum_ia_iA_{X,B}(v_i)-\|\mu\|_\omega .
\]
Thus \(\beta(\mu)\ge\varepsilon\|\mu\|_\omega\) for all divisorial \(\mu\)
if and only if \(\sum_ia_iA(v_i)\ge(1+\varepsilon)\|\mu\|_\omega\).  Taking
\(\mu=\delta_v\) gives \(A(v)\ge(1+\varepsilon)S_L(v)\), that is
\(\delta\ge1+\varepsilon\).  Conversely, if \(\delta\ge1+\eta\) then
\(\sum_ia_iA(v_i)\ge(1+\eta)\sum_ia_iS_L(v_i)\ge(1+\eta)\|\mu\|_\omega\),
the last step by convexity of \(E^\vee_\omega\)
(\Cref{prop:ch47-dual-elementary}\,\cref{itm:ch47-dual:convex}) applied to
\(\|\mu\|_\omega=E^\vee_\omega(\sum a_i\delta_{v_i})
\le\sum a_iE^\vee_\omega(\delta_{v_i})=\sum a_iS_L(v_i)\).  Hence
\(\beta(\mu)\ge\frac{\eta}{1+\eta}\|\mu\|_\omega\), which gives
\(\sigma_{\div}\ge\eta/(1+\eta)\).  The reverse constant is the
specialization to Dirac masses.
\end{proof}

\begin{important}{Imported, with dates: the link with K-stability}
\label{imp:ch50-fano-imports}
What \Cref{prop:ch50-valuative-internal} does \emph{not} prove is that any
of the valuative conditions is equivalent to K-semistability or to uniform
K-stability in the sense of \Cref{def:ch50-k-stability}.  That is a
separate body of results, each with its own date, source and hypotheses.
Throughout, \((X,B)\) is log Fano over an algebraically closed field of
characteristic zero and \(L=-K_{X,B}\).
\begin{enumerate}
\item\label{itm:ch50-fano:fujita-li}
      \emph{Uniform K-stability \(\Leftrightarrow\)
      \(\beta_L\ge\varepsilon j_L\).}  Fujita, ``A valuative criterion for
      uniform K-stability of \(\mathbf Q\)-Fano varieties'',
      \cite{Fuj19} (preprint 2016, published 2019), together with
      Li, ``K-semistability is equivariant volume minimization'',
      \cite{Li17} (preprint 2015, published 2017).  The hypotheses are:
      \(X\) a \(\mathbf Q\)-Fano variety, that is normal, klt and with
      \(-K_X\) ample \(\mathbf Q\)-Cartier, over an algebraically closed
      field of characteristic zero.  The proof uses the minimal model
      program \cite{BCHM10} through the reduction to special test
      configurations.
\item\label{itm:ch50-fano:bj}
      \emph{The threshold formulation \eqref{eq:stability-delta}, and the
      identification of \(\delta\) with an asymptotic log canonical
      threshold.}  Blum and Jonsson, ``Thresholds, valuations and
      K-stability'', \cite{BJ20} (preprint 2017, published 2020).  The
      hypotheses there are stated for a log pair \((X,B)\) with \(L\) an
      arbitrary big line bundle for the definition of \(\delta_L\), and
      specialize to the log Fano case for the K-stability statement.  The
      same paper proves that \(\delta\) may be computed by \emph{quasi-monomial}
      valuations and, in the presence of a torus action, by
      invariant valuations; this equivariance is the input used in
      \Cref{ex:ch50-toric-F1}.
\item\label{itm:ch50-fano:semistable}
      \emph{K-semistability \(\Leftrightarrow\) \(\delta\ge1\)
      \(\Leftrightarrow\) \(\beta_L\ge0\).}  This is the combination of
      \cref{itm:ch50-fano:fujita-li} and \cref{itm:ch50-fano:bj}; the easy
      direction --- K-semistable implies \(\beta_L\ge0\) --- follows from
      \Cref{prop:ch50-model-form}\,\cref{itm:ch50-mf:semi} together with
      \Cref{cor:ch50-M-dirac} whenever the Green function of \(v\) is
      \(\mathbf Q\)-PL, and in general from \cite{Fuj19}.
\item\label{itm:ch50-fano:lxz}
      \emph{Finite generation, and the equality of the K-stability and
      filtration thresholds in the log Fano case.}  Liu, Xu and Zhuang,
      ``Finite generation for valuations computing stability thresholds
      and applications to K-stability'', \cite{LXZ22} (preprint 2021,
      published 2022).  Hypotheses: \((X,B)\) log Fano over an
      algebraically closed field of characteristic zero.  The theorem
      states that when \(\delta(X,B)\le1\) the infimum in
      \eqref{eq:stability-delta} is attained by a quasi-monomial valuation
      with finitely generated associated graded ring.  This is what closes,
      \emph{for the anticanonical polarization only}, the gap recorded in
      \Cref{rem:stability-gap-general-polarization}.
\item\label{itm:ch50-fano:openness}
      \emph{Related later developments} --- properness of K-moduli and
      optimal destabilizations \cite{BHLLX21} (published 2021),
      \(G\)-uniform stability \cite{Li22} (published 2022), continuity of
      \(\delta\) \cite{Zhang21} (published 2021) --- are not used in this
      chapter.  They are listed here so that the reader can see the
      boundary of what has been imported.
\end{enumerate}
None of these is proved in this book, and none of them is a consequence of
the pluripotential theory developed in \Cref{part:bj}.  Each is used only
where it is cited by name.
\end{important}

\begin{warning}{The anticanonical hypothesis is essential}
\label{warn:ch50-anticanonical-essential}
Every statement in \Cref{imp:ch50-fano-imports} assumes \(L=-K_{X,B}\).
For a general polarization there is no \(\delta\)-invariant criterion for
uniform K-stability, and the reductions to special test configurations are
unavailable, the minimal model program having nothing to say about a pair
whose polarization is unrelated to the canonical class.  What survives for
a general \(L\) is the pluripotential picture of
\Cref{sec:ch50-uniform}: the three thresholds
\eqref{eq:ch50-three-thresholds}, the equality
\(\sigma_{\mathcal N}=\sigma_{\div}\) of
\Cref{thm:stability-divisorial-filtrations}, and the inequality
\eqref{eq:stability-threshold-inequality}.  A reader who transplants
\eqref{eq:stability-delta} to a general \(L\) obtains a well-defined
number, the \(\delta\)-invariant \(\delta_L(X,B)\) of \cite{BJ20}, but its
relation to K-stability is not the one stated in
\Cref{thm:stability-valuative-fano}.
\end{warning}

\subsection{The projective line at the threshold}
\label{sec:ch50-p1-delta}

\begin{example}[\(\mathbf P^1\) at the threshold]
\label{ex:stability-p1-delta}
Take \(X=\mathbf P^1\), \(B=0\), and \(L=-K_X=\cO(2)\).  For
\(v=\ord_p\),
\[
 A_X(v)=1,\qquad
 \vol(L-tv)=\max\{2-t,0\},\qquad
 S_L(v)=\frac12\int_0^2(2-t)\,\rd t=1.
\]
Thus \(T_L(v)=2\), \(\beta_L(v)=0\), and \(j_L(v)=1\).  Hence
\(\delta(\mathbf P^1)=1\): the curve is K-polystable but not uniformly
K-stable in the absolute sense.  The obstruction is exactly its
nontrivial automorphism group, reflected by product configurations of
positive norm and zero Futaki invariant.
\end{example}

\begin{remark}[Three consistency checks on
\Cref{ex:stability-p1-delta}]
\label{rem:ch50-p1-delta-checks}
First, \(S_L(\ord_p)=1\) agrees with \(E^\vee_\omega(\delta_{\ord_p})\):
by \Cref{prop:ch49-psef-threshold}, \(T_\omega(\ord_p)=\tau_{\cO(2)}(p)=2\),
and by \eqref{eq:bj-dirac-energy-T} with \(n=1\) both bounds are
\(T_\omega/2=1\), so \(E^\vee_\omega(\delta_{\ord_p})=1\) exactly.  This
verifies \Cref{imp:ch50-imported}\,\cref{itm:ch50-import:S} in this case.

Second, the toric formula agrees: with \(P=[-1,1]\) the reflexive polytope
of \(-K_{\mathbf P^1}\), one has \(f_P(u)=|u|\), \(b_P=0\), and
\Cref{prop:ch49-toric-green}\,\cref{itm:ch49-tg:energy} gives
\(E^\vee_\omega(\delta_{v_u})=|u|+0=|u|\), which at \(u=1\) is \(1\).  Also
\(A_{\mathbf P^1}(v_u)=f_P(u)=|u|\), so \(\beta_L(v_u)=0\) for every toric
\(u\), and \(\delta=1\).

Third, the Mabuchi functional vanishes on the corresponding metrics:
\Cref{cor:ch50-M-dirac} gives
\(\beta_{X;\omega}(\delta_{\ord_p})=A-S=0\), and this is
\Cref{prop:ch50-p1-product}\,\cref{itm:ch50-p1p:mabuchi} computed with
\(L=\cO(1)\) instead of \(\cO(2)\), the two differing by the scaling
\eqref{eq:ch50-three-thresholds} which leaves the threshold unchanged.
\end{remark}

\section{Four worked examples}
\label{sec:ch50-examples}

\subsection{All the functionals of an explicit test configuration on
\texorpdfstring{\(\mathbf P^1\)}{P1}}
\label{sec:ch50-ex-p1}

\begin{example}[Deformation to the normal cone on \((\mathbf P^1,-K)\)]
\label{ex:ch50-p1-dnc}
Let \(X=\mathbf P^1\), \(B=0\), \(p=[0:1]\), and
\(L=-K_X=\cO(2)\), so that \(n=1\), \(V=2\), \(\gamma=-1\) and
\(\overline S=1\).  Let \(\cX\to\mathbf A^1\) be the blowup of
\(\mathbf P^1\times\mathbf A^1\) at \((p,0)\), as in
\Cref{ex:ch33-dnc-p1}, with exceptional curve \(E\), strict transform
\(\widetilde E\) of the central fibre, and \(b_E=b_{\widetilde E}=1\).
For \(0<c<2\) put
\[
 \cL_c:=\mu^*\operatorname{pr}_1^*\cO(2)-cE,
 \qquad
 \phi_c:=\phi_{\cL_c}=c\,\phi_{\mathfrak a},
 \qquad
 \phi_{\mathfrak a}(t\ord_p)=-\min\{t,1\},
\tag{50.84}\label{eq:ch50-p1-dnc-metric}
\]
with \(\phi_{\mathfrak a}=0\) off the ray through \(p\).  Since
\((\cL_c\cdot E)=c\) and \((\cL_c\cdot\widetilde E)=2-c\), the polarization
is relatively ample for \(0<c<2\) and semiample at \(c=2\).  Then, for
\(0<c<2\):
\[
 \begin{array}{ll}
 \DHm_{\chi_c}=\tfrac{2-c}{2}\,\delta_0
   +\tfrac12\mathbf 1_{[-c,0]}\,\rd\lambda,
 &\quad \lambda_{\max}=0,\\[2mm]
 E(\phi_c)=-\dfrac{c^2}{4},
 &\quad J(\phi_c)=\dfrac{c^2}4,\\[2mm]
 \MA_\omega(\phi_c)=\dfrac c2\delta_{\ord_p}
        +\dfrac{2-c}2\delta_{v_{\triv}},
 &\quad \displaystyle\int\phi_c\,\MA_\omega(\phi_c)=-\dfrac{c^2}2,\\[2mm]
 \|\chi_c\|_{\min}=\dfrac{c^2}4,
 &\quad I(\phi_c)=\dfrac{c^2}2,\\[2mm]
 \|\chi_c\|_1=\dfrac{(2-c)c^2}4+\dfrac{c^4}{32},
 &\quad H^{\NA}(\phi_c)=\dfrac c2,\\[2mm]
 R^{\NA}(\phi_c)=0,
 &\quad M^{\NA}(\phi_c)=\DF(\cX,\cL_c)=\dfrac{c(2-c)}4,\\[2mm]
 L^{\NA}(\phi_c)=\min\{0,1-c\},
 &\quad D^{\NA}(\phi_c)=\min\{0,1-c\}+\dfrac{c^2}4.
 \end{array}
\tag{50.85}\label{eq:ch50-p1-dnc-table}
\]
\end{example}

\begin{proof}[Verification]
\emph{The filtration and its spectral measure.}  A section
\(x^ay^{2m-a}\) of \(mL=\cO(2m)\) has \(\ord_E=a\) on \(\cX\) because
\(\ord_E(z)=1\) for \(z=x/y\); viewing it as a section of
\(m\cL_c\) subtracts \(mc\) from the coefficient along \(E\), and
\(\operatorname{div}(\varpi)=\widetilde E+E\) subtracts \(\lambda\) from
both central components.  Hence, exactly as in
\eqref{eq:ch33-dnc-weights},
\(\chi_{c,m}(x^ay^{2m-a})=\min\{0,\ a-mc\}\) for \(0\le a\le 2m\).  The
normalized jumping values are \(\min\{0,u-c\}\) with \(u=a/m\) ranging over
\([0,2]\), so the concave transform on the Okounkov body \(\Delta=[0,2]\)
is \(G(u)=\min\{0,u-c\}\) and \eqref{eq:ch33-dh-is-pushforward} gives the
displayed \(\DHm_{\chi_c}\), of total mass
\(\frac{2-c}2+\frac c2=1\).  Its top is \(\lambda_{\max}=0\), which is
\(\sup\phi_c\) by \Cref{prop:ch50-sup-lambdamax}, and its first moment is
\(\frac12\int_0^2\min\{0,u-c\}\rd u=-\frac{c^2}4\), which is
\(E(\phi_c)\) by \Cref{thm:ch50-energy-dh}.  As a check,
\((\overline\cL_c^{\,2})=-c^2\) and
\((\overline\cL_c^{\,2})/((n+1)V)=-c^2/4\).

\emph{Monge--Amp\`ere measure.}  By \Cref{prop:ch46-pl-measure},
\(\MA^{\mathrm u}_\omega(\phi_c)=(\cL_c\cdot E)\delta_{v_E}
+(\cL_c\cdot\widetilde E)\delta_{v_{\widetilde E}}
=c\,\delta_{\ord_p}+(2-c)\delta_{v_{\triv}}\), since
\(v_E=\ord_p\) and \(v_{\widetilde E}=v_{\triv}\).  Dividing by \(V=2\)
gives the probability measure.  Then
\(\int\phi_c\,\MA_\omega(\phi_c)=\frac c2\phi_c(\ord_p)
=\frac c2(-c)=-\frac{c^2}2\).

\emph{The three norms.}  \(J=\sup\phi_c-E=0+\frac{c^2}4\);
\(\|\chi_c\|_{\min}=E-\int\phi_c\MA_\omega(\phi_c)=-\frac{c^2}4+\frac{c^2}2
=\frac{c^2}4\) by \eqref{eq:ch50-min-norm-identities}; and
\(I=J+\|\chi_c\|_{\min}=\frac{c^2}2\), consistent with the sharp bounds
\(\frac{I}{n+1}\le J\le\frac{nI}{n+1}\), which for \(n=1\) force
\(J=\|\chi_c\|_{\min}=I/2\).  For the centred \(L^1\)-norm, with
\(\bar\lambda=-c^2/4\),
\[
 \|\chi_c\|_1=2\int(\lambda-\bar\lambda)_+\rd\DHm_{\chi_c}
 =2\Bigl[\frac{2-c}2\cdot\frac{c^2}4
     +\frac12\int_{-c^2/4}^{0}\Bigl(\lambda+\frac{c^2}4\Bigr)\rd\lambda
     \Bigr]
 =\frac{(2-c)c^2}{4}+\frac{c^4}{32},
\]
the second term using \(c^2/4\le c\).  At \(c=1\) this is
\(\frac14+\frac1{32}=\frac9{32}\), and \(J=\frac14\), so the ratio is
\(\frac9{8}\), comfortably inside the window
\([c_1,2]=[\frac12,2]\) of \eqref{eq:stability-J-L1-comparison}.

\emph{Entropy, Ricci energy, Mabuchi.}  \(A_X(\ord_p)=1\) and
\(A_X(v_{\triv})=0\), so \(H^{\NA}(\phi_c)=\frac c2\).  By
\Cref{lem:ch50-ricci-energy} with \(\gamma=-1\),
\(R^{\NA}=-\bigl(2E-\int\phi_c\MA\phi_c\bigr)
=-\bigl(-\frac{c^2}2+\frac{c^2}2\bigr)=0\).  Then
\eqref{eq:stability-chen-tian} gives
\(M^{\NA}=\frac c2+0+1\cdot(-\frac{c^2}4)=\frac{c(2-c)}4\), and
\Cref{prop:ch50-M-proportional} gives the same:
\(M^{\NA}=H+\gamma\|\chi_c\|_{\min}=\frac c2-\frac{c^2}4\).  The central
fibre \(\widetilde E+E\) is reduced, so \(\DF=M^{\NA}\) by
\eqref{eq:stability-M-DF-correction}.

\emph{Ding.}  On the ray through \(p\), \(A_X(t\ord_p)+\phi_c(t\ord_p)
=t-c\min\{t,1\}\), which equals \((1-c)t\) for \(t\le1\) and \(t-c\) for
\(t\ge1\); its infimum over \(t\ge0\) is \(0\) if \(c\le1\) and \(1-c\) if
\(c\ge1\).  On every other ray \(\phi_c=0\) and \(A_X\ge0\), so the
infimum there is \(0\), attained at \(v_{\triv}\).  Hence
\(L^{\NA}(\phi_c)=\min\{0,1-c\}\) and
\(D^{\NA}=L^{\NA}-E=\min\{0,1-c\}+\frac{c^2}4\).
\end{proof}

\begin{remark}[What \Cref{ex:ch50-p1-dnc} verifies]
\label{rem:ch50-p1-dnc-checks}
Four separate statements of this chapter are confirmed on the nose.

\emph{(a) \(D^{\NA}\le M^{\NA}\), with equality exactly at the
anticanonical member.}  From \eqref{eq:ch50-p1-dnc-table},
\[
 M^{\NA}(\phi_c)-D^{\NA}(\phi_c)
 =\begin{cases}
   \tfrac12c(1-c), & 0<c\le1,\\[1mm]
   -\tfrac12(c-1)(c-2), & 1\le c<2,
  \end{cases}
\]
which is \(\ge0\) on \((0,2)\) and vanishes only at \(c=1\).  And
\(c=1\) is precisely the anticanonical member: on \(\cX\) one has
\(K_{\cX/\mathbf A^1}=\mu^*\operatorname{pr}_1^*K_{\mathbf P^1}+E\) and
\(\cX_{0,\mathrm{red}}=\cX_0=\mu^*(\mathbf P^1\times\{0\})\), so
\[
 -K^{\log}_{\cX/\mathbf A^1}+\cX_0
 =\mu^*\operatorname{pr}_1^*\cO(2)-E
 =\cL_1 ,
\]
which is the anticanonical condition of \Cref{sec:ch50-anticanonical} with
constant \(c=0\), the divisor \(\cX_0\) being principal and so contributing
nothing to the line bundle --- but note that \(c\) is not read off the bundle:
it is pinned by \Cref{lem:ch50-anticanonical-constant}, as the computation
just below does.  The pair
\((\cX,\cX_{0,\mathrm{red}})\) is log smooth, hence sublc.  This is
\Cref{thm:stability-ding-mabuchi} and
\Cref{lem:ch50-anticanonical-constant}: indeed
\(A_X(v_E)+\phi_1(v_E)=1+(-1)=0\) and
\(A_X(v_{\widetilde E})+\phi_1(v_{\widetilde E})=0+0=0\), both equal to
the constant \(c=0\) of \eqref{eq:ch50-anticanonical-components}, and
\(L^{\NA}(\phi_1)=0\).

\emph{(b) K-semistability but not uniform K-stability.}  \(\DF>0\) for
every \(c\in(0,2)\), but
\[
 \frac{M^{\NA}(\phi_c)}{J(\phi_c)}
 =\frac{c(2-c)/4}{c^2/4}=\frac{2-c}{c}\ \xrightarrow[c\to2^-]{}\ 0 ,
\]
so no uniform inequality \eqref{eq:stability-M-coercive-model} holds.  At
\(c=2\) the polarization contracts \(\widetilde E\) and the configuration
becomes the product configuration of \Cref{ex:stability-p1-product},
twisted by a character; there \(M^{\NA}=0\) and \(J=1\).  Thus a single
one-parameter family exhibits the failure recorded in
\Cref{rem:ch50-p1-obstruction}.

\emph{(c) The sharp comparison \eqref{eq:stability-J-L1-comparison}
degenerates correctly.}  As \(c\to2\),
\(\|\chi_c\|_1\to\frac{16}{32}=\frac12\) and \(J\to1\), ratio
\(\frac12=c_1\): the extremal constant of \Cref{thm:ch50-J-L1} is attained
in the limit, and \Cref{rem:ch50-extremal-cone} identifies the limiting
measure \(\frac12\mathbf 1_{[-2,0]}\rd\lambda\) as the cone measure for
\(n=1\).

\emph{(d) The valuative invariants.}  \(T_L(\ord_p)=\tau_{\cO(2)}(p)=2\)
and \(S_L(\ord_p)=1=E^\vee_\omega(\delta_{\ord_p})\) by
\Cref{ex:stability-p1-delta}, so
\(\beta_L(\ord_p)=0\) and \(\delta(\mathbf P^1)=1\), consistent with (b).
\end{remark}

\subsection{A toric example decided from the polytope}
\label{sec:ch50-ex-toric}

We first record the general toric computation, which is a two-line
consequence of the formulae of \Cref{chap:bj-linear-growth}.

\begin{proposition}
\label{prop:ch50-toric-beta}
\index{toric geometry!stability threshold}
Let \(X=X_\Sigma\) be a smooth projective toric variety over \(k\) with
\(-K_X\) ample, let \(u_1,\dots,u_r\in N\) be the primitive generators of
the rays of \(\Sigma\), let \(L=-K_X\), and let
\[
 P:=\{m\in M_{\mathbf R}\ :\ \langle u_i,m\rangle\ge-1
     \text{ for }i=1,\dots,r\}
\]
be the associated reflexive polytope, with barycenter
\(b_P=\vol(P)^{-1}\int_Pm\,\rd m\).  Write
\(f_P(u)=-\min_{m\in P}\langle u,m\rangle\) for the support function of
\(-P\), and \(v_u\in X^{\divis}\) for the toric valuation attached to
\(u\in N_{\mathbf R}\).  Then, for every \(u\in N_{\mathbf R}\):
\[
 A_X(v_u)=f_P(u),
 \qquad
 S_L(v_u)=f_P(u)+\langle u,b_P\rangle,
 \qquad
 \beta_L(v_u)=-\langle u,b_P\rangle .
\tag{50.86}\label{eq:ch50-toric-beta}
\]
Consequently, with \(Q:=\operatorname{conv}\{u_1,\dots,u_r\}
=\{u:A_X(v_u)\le1\}\),
\[
 \inf_{u\ne0}\frac{A_X(v_u)}{S_L(v_u)}
 =\frac{1}{1+\max_{u\in Q}\langle u,b_P\rangle}\ \le\ 1,
\tag{50.87}\label{eq:ch50-toric-delta}
\]
with equality to \(1\) if and only if \(b_P=0\).
\end{proposition}

\begin{proof}
\emph{Log discrepancy.}  Since \(X\) is smooth toric,
\((X,\partial X)\) is a log smooth pair whose strata are the torus orbits,
and the quasimonomial valuation attached to a cone
\(\sigma=\operatorname{Cone}(u_{i_1},\dots,u_{i_k})\) with weights
\(\alpha=(\alpha_1,\dots,\alpha_k)\) is \(v_u\) with
\(u=\sum_j\alpha_ju_{i_j}\).  Hence
\Cref{prop:ch32-log-discrepancy-affine} gives
\(A_X(v_u)=\sum_j\alpha_jA_X(\ord_{D_{i_j}})=\sum_j\alpha_j\), because
\(D_{i}\) is a prime divisor on \(X\) itself, so
\(A_X(\ord_{D_i})=1\).  Thus \(A_X\circ v\) is the piecewise linear
function on \(N_{\mathbf R}\) that is linear on each cone of \(\Sigma\) and
takes the value \(1\) at each \(u_i\).  On the other hand
\(f_P\) is the support function of \(-P\), hence is linear on each cone of
the normal fan of \(P\), which is \(\Sigma\); and
\(f_P(u_i)=-\min_{m\in P}\langle u_i,m\rangle=1\), because the inequality
\(\langle u_i,m\rangle\ge-1\) defines a facet of \(P\), \(-K_X\) being
ample.  Two functions linear on the same fan and agreeing on the
primitive generators agree, so \(A_X(v_u)=f_P(u)\).

\emph{The \(S\)-invariant.}  By
\Cref{imp:ch50-imported}\,\cref{itm:ch50-import:S} and
\Cref{prop:ch49-toric-green}\,\cref{itm:ch49-tg:energy},
\(S_L(v_u)=E^\vee_\omega(\delta_{v_u})=f_P(u)+\langle u,b_P\rangle\).
Subtracting gives \(\beta_L(v_u)=-\langle u,b_P\rangle\).  Both sides are
independent of the translation chosen for \(P\), since a translation by
\(m_0\) changes \(f_P\) by \(-\langle\cdot,m_0\rangle\) and \(b_P\) by
\(+m_0\); the reflexive normalization is the one for which the first
identity of \eqref{eq:ch50-toric-beta} holds.

\emph{The threshold.}  The ratio \(A/S\) is homogeneous of degree zero in
\(u\), so it may be computed on \(\{A_X(v_u)=1\}=\partial Q\); there
\(A/S=(1+\langle u,b_P\rangle)^{-1}\), which is minimized when
\(\langle u,b_P\rangle\) is maximal, giving
\eqref{eq:ch50-toric-delta}.  Finally
\(\max_{u\in Q}\langle u,b_P\rangle\ge0\) always, with equality if and only
if \(\langle u,b_P\rangle\le0\) for all \(u\in Q\); since \(Q\) is a
neighbourhood of the origin, this forces \(b_P=0\).
\end{proof}

\begin{example}[\(\mathbb F_1=\mathrm{Bl}_{p}\mathbf P^2\) is not
K-semistable, and \(\delta=6/7\)]
\label{ex:ch50-toric-F1}
Let \(X=\mathbb F_1\) be the blowup of \(\mathbf P^2\) at a torus-fixed
point, a smooth toric del Pezzo surface with
\[
 u_1=(1,0),\quad u_2=(0,1),\quad u_3=(-1,-1),\quad u_4=(1,1),
\]
and \(L=-K_X\).  The reflexive polytope is
\[
 P=\{m_1\ge-1,\ m_2\ge-1,\ -1\le m_1+m_2\le1\},
\]
a quadrilateral with vertices \((-1,0),(-1,2),(2,-1),(0,-1)\), of area
\(4\); accordingly \(V=(-K_X)^2=2!\cdot4=8\).  Its barycenter is
\[
 b_P=\Bigl(\tfrac1{12},\tfrac1{12}\Bigr).
\tag{50.88}\label{eq:ch50-F1-barycenter}
\]
Consequently, by \Cref{prop:ch50-toric-beta}:
\begin{enumerate}
\item\label{itm:ch50-F1:beta}
      \(\beta_L(v_u)=-\tfrac1{12}(u^{(1)}+u^{(2)})\) for
      \(u=(u^{(1)},u^{(2)})\); in particular, for the exceptional divisor
      \(E\), which corresponds to \(u_4=(1,1)\),
      \[
       A_X(\ord_E)=1,\qquad S_L(\ord_E)=\tfrac76,\qquad
       \beta_L(\ord_E)=-\tfrac16<0 ;
      \]
\item\label{itm:ch50-F1:delta}
      \(Q=\operatorname{conv}\{u_1,u_2,u_3,u_4\}\) and
      \(\max_{u\in Q}\langle u,b_P\rangle=\langle u_4,b_P\rangle=\tfrac16\),
      so the toric infimum in \eqref{eq:ch50-toric-delta} equals
      \(\tfrac{1}{1+1/6}=\tfrac67\);
\item\label{itm:ch50-F1:destab}
      the Green function \(\phi_{u_4}\) of \(v_{u_4}=\ord_E\) is
      \(\mathbf Q\)-PL, and
      \(M^{\NA}(\phi_{u_4})=\beta_L(\ord_E)=-\tfrac16<0\); hence
      \((\mathbb F_1,-K)\) is not K-semistable.
\end{enumerate}
\end{example}

\begin{proof}[Verification]
\emph{The polytope.}  The four inequalities \(\langle u_i,m\rangle\ge-1\)
read \(m_1\ge-1\), \(m_2\ge-1\), \(m_1+m_2\le1\) and \(m_1+m_2\ge-1\).  The
listed vertices are the pairwise intersections that satisfy all four.  The
shoelace formula on \((-1,0),(-1,2),(2,-1),(0,-1)\) gives twice the area
\(=8\), so \(\vol(P)=4\) and \(V=2!\cdot4=8\), matching
\((-K_{\mathbb F_1})^2=8\).  For the barycenter, decompose \(P\) as the
triangle \(T_1\) with vertices \((-1,-1),(-1,2),(2,-1)\), of area
\(\tfrac92\) and centroid \((0,0)\), minus the triangle \(T_0\) with
vertices \((-1,-1),(-1,0),(0,-1)\), of area \(\tfrac12\) and centroid
\((-\tfrac23,-\tfrac23)\).  Then
\(\int_Pm\,\rd m=\tfrac92(0,0)-\tfrac12(-\tfrac23,-\tfrac23)
=(\tfrac13,\tfrac13)\), and dividing by \(\vol(P)=4\) gives
\eqref{eq:ch50-F1-barycenter}.

\emph{\cref{itm:ch50-F1:beta}.}  \(\beta_L(v_u)=-\langle u,b_P\rangle\) is
\eqref{eq:ch50-toric-beta}.  At \(u_4=(1,1)\):
\(f_P(u_4)=-\min_{m\in P}(m_1+m_2)=1\), which is \(A_X(\ord_E)=1\) as it
must be, \(E\) being a prime divisor on \(X\); and
\(\langle u_4,b_P\rangle=\tfrac16\), so \(S_L=1+\tfrac16=\tfrac76\).

As an independent check of \(S_L(\ord_E)=\tfrac76\), compute the volume
integral directly.  Write \(H\) for the pullback of a line and
\(E\) for the exceptional curve, so \(-K_X=3H-E\), \(H^2=1\),
\(H\cdot E=0\), \(E^2=-1\).  For \(s\ge0\) the class \(3H-sE\) is nef
exactly for \(0\le s\le3\), since \((3H-sE)\cdot E=s\ge0\) and
\((3H-sE)\cdot(H-E)=3-s\ge0\), and then
\(\vol(3H-sE)=(3H-sE)^2=9-s^2\).  Hence
\[
 \begin{split}
 S_L(\ord_E)
 &=\frac1V\int_0^\infty\vol(-K_X-t\ord_E)\,\rd t
  =\frac18\int_0^2\bigl(9-(1+t)^2\bigr)\rd t\\
 &=\frac18\int_1^3(9-s^2)\,\rd s
  =\frac18\Bigl(18-\frac{26}{3}\Bigr)=\frac76 .
 \end{split}
\]
This confirms simultaneously the toric formula
\eqref{eq:ch50-toric-beta} and the imported identity
\(S_L=E^\vee_\omega(\delta_\cdot)\) of
\Cref{imp:ch50-imported}\,\cref{itm:ch50-import:S} in this case.

\emph{\cref{itm:ch50-F1:delta}.}  The linear functional
\(u\mapsto\langle u,b_P\rangle=\tfrac1{12}(u^{(1)}+u^{(2)})\) takes the
values \(\tfrac1{12},\tfrac1{12},-\tfrac16,\tfrac16\) at
\(u_1,u_2,u_3,u_4\); its maximum over the polytope \(Q\) is attained at a
vertex, hence equals \(\tfrac16\).  Now apply
\eqref{eq:ch50-toric-delta}.

\emph{\cref{itm:ch50-F1:destab}.}  By
\Cref{prop:ch49-toric-green}\,\cref{itm:ch49-tg:ma} the Green function
\(\phi_{u_4}\) is \(\mathbf Q\)-PL, \(u_4\) being a lattice point, and
\(\MA_\omega(\phi_{u_4})=\delta_{v_{u_4}}\).  By
\Cref{thm:ch34-pl-density} and \Cref{thm:tc-pl-correspondence} a
\(\mathbf Q\)-PL \(\omega\)-psh function is the Fubini--Study function of a
semiample normal test configuration; so \Cref{cor:ch50-M-dirac} applies and
gives \(M^{\NA}(\phi_{u_4})=A_X(\ord_E)-S_L(\ord_E)=-\tfrac16\).  By
\Cref{prop:ch50-model-form}\,\cref{itm:ch50-mf:semi}, a K-semistable pair
has \(M^{\NA}\ge0\) on every positive model metric.  Hence
\((\mathbb F_1,-K)\) is not K-semistable.
\end{proof}

\begin{remark}[Toric log Fanos: the complete picture, and its dating]
\label{rem:ch50-toric-picture}
\Cref{prop:ch50-toric-beta} shows that on the toric slice the whole theory
collapses to one vector, the barycenter.  Three consequences, with their
statuses.

\emph{(a) Book-internal.}  If \(b_P\ne0\), then \(X\) is not K-semistable:
this needs only \Cref{prop:ch50-toric-beta}, the rationality of some
\(u\in N\) with \(\langle u,b_P\rangle>0\), and
\Cref{prop:ch50-model-form}, as in
\Cref{ex:ch50-toric-F1}\,\cref{itm:ch50-F1:destab}.

\emph{(b) Book-internal.}  No smooth toric log Fano is uniformly
K-stable.  Indeed for any \(u\ne0\), \(\beta_L(v_u)+\beta_L(v_{-u})
=-\langle u,b_P\rangle-\langle-u,b_P\rangle=0\), while
\(j_L(v_u)+j_L(v_{-u})>0\) by
\Cref{prop:ch50-valuative-dh}\,\cref{itm:ch50-vdh:comparison}; so
\(\beta_L\ge\varepsilon j_L\) fails for every \(\varepsilon>0\) on the pair
\(\{v_u,v_{-u}\}\).  This is the valuative shadow of the fact that the
torus sits inside \(\Aut(X)\).

\emph{(c) Imported.}  If \(b_P=0\), then \(X\) \emph{is} K-semistable.
This direction needs \(\delta=\delta_{\mathrm{toric}}\), that is, that the
infimum \eqref{eq:stability-delta} over all divisorial valuations is
computed by torus-invariant ones, which is the equivariance statement of
\cite{BJ20} (preprint 2017, published 2020) recorded in
\Cref{imp:ch50-fano-imports}\,\cref{itm:ch50-fano:bj}; and then
K-semistability follows from
\Cref{thm:stability-valuative-fano}\,(3), imported from
\cite{Fuj19,Li17}.  We do not prove either.  With those inputs,
\(\delta(X)=(1+\max_{u\in Q}\langle u,b_P\rangle)^{-1}\) for every smooth
toric log Fano, and \(\delta(\mathbb F_1)=\tfrac67\).

For \(\mathbf P^2\) and \(\mathbf P^1\times\mathbf P^1\) one has
\(b_P=0\) by symmetry, so \(\delta=1\): both are K-semistable --- indeed
K-polystable --- and neither is uniformly K-stable.  This is the promised
family of examples separating the two notions in every dimension.
\end{remark}

\subsection{A filtration that is not finitely generated}
\label{sec:ch50-ex-harmonic}

\begin{example}[The harmonic filtration on \(\mathbf P^1\), with all its
invariants]
\label{ex:ch50-harmonic}
Let \(X=\mathbf P^1\), \(B=0\), \(L=\cO(1)\), so \(n=1\), \(V=1\),
\(K_X=-2\omega\) and \(\gamma=-2\), \(\overline S=2\).  Let \(\chi\) be the
multiplicative, linearly bounded filtration of
\Cref{ex:ch33-harmonic-not-fg}, defined on monomials by
\[
 \chi_m(x^ay^{b})=\Psi(a,b)=\frac{ab}{a+b},\qquad a+b=m,
\]
and by the minimum over the monomial support in general.  It is not of
finite type, and no test configuration induces it.  Nevertheless every
functional of this chapter is defined and computable:
\[
 \begin{array}{ll}
 \DHm_\chi=\dfrac{2\,\rd\tau}{\sqrt{1-4\tau}}
     \mathbf 1_{[0,1/4)}(\tau),
 &\quad \lambda_{\max}=\dfrac14,\\[3mm]
 \phi_\chi(t\ord_p)=\dfrac{(1-t)^2}{4}\ (t\le1),\ 0\ (t\ge1),
 &\quad E(\phi_\chi)=\dfrac16,\\[3mm]
 \MA_\omega(\phi_\chi)=\dfrac12\mathbf 1_{[0,1]}\,\rd t
       \ \text{on each ray},
 &\quad J(\phi_\chi)=\dfrac1{12},\\[3mm]
 \displaystyle\int\phi_\chi\,\MA_\omega(\phi_\chi)=\dfrac1{12},
 &\quad \|\chi\|_{\min}=\dfrac1{12},\\[3mm]
 I(\phi_\chi)=\dfrac16,
 &\quad \|\chi\|_1=\dfrac{\sqrt3}{27},\\[3mm]
 H^{\NA}(\phi_\chi)=\dfrac12,
 &\quad M_{X;L}(\chi)=\dfrac13 .
 \end{array}
\tag{50.89}\label{eq:ch50-harmonic-table}
\]
\end{example}

\begin{proof}[Verification]
\emph{The spectral data.}  The concave transform on the Okounkov body
\(\Delta=[0,1]\) is \(G(u)=u(1-u)\), so
\(\DHm_\chi([\tau,\infty))=\lvert\{u\in[0,1]:u(1-u)\ge\tau\}\rvert
=\sqrt{1-4\tau}\) for \(0\le\tau\le\tfrac14\); differentiating gives the
displayed density, which is \eqref{eq:ch33-harmonic-dh}.  Its top is
\(\lambda_{\max}=\tfrac14\), and its first moment is
\(\int_0^1u(1-u)\rd u=\tfrac16\).  The Fubini--Study function is
\eqref{eq:ch33-harmonic-fs}, and \(\sup\phi_\chi=\phi_\chi(v_{\triv})
=\tfrac14=\lambda_{\max}\), as \Cref{prop:ch50-sup-lambdamax} requires.
Hence \(J=\tfrac14-\tfrac16=\tfrac1{12}\).

\emph{The Monge--Amp\`ere measure is diffuse.}  On the ray through
\(p=[0:1]\), \(g(t):=\phi_\chi(t\ord_p)=\tfrac14(1-t)^2\) for
\(t\in[0,1]\) and \(0\) for \(t\ge1\).  Its distributional second
derivative on \((0,\infty)\) is \(\tfrac12\mathbf 1_{[0,1]}\rd t\): the
function is \(C^1\) at \(t=1\), so there is no atom there.  The same holds
on the ray through \([1:0]\) by symmetry, and \(g\equiv\tfrac14\) on every
other ray.  The two rays carry total mass
\(2\cdot\tfrac12=1=V\), so no mass remains at \(v_{\triv}\); this is
\eqref{eq:bj-curve-laplacian}.  \emph{In particular
\(\MA_\omega(\phi_\chi)\) has no atoms at all}, in sharp contrast with the
atomic formula \eqref{eq:bj-model-ma-atomic} valid for model metrics.

\emph{The energy, twice.}  From the first moment, \(E=\tfrac16\).
Independently, for \(n=1\) the Bernstein expansion of
\Cref{thm:ch47-segment-expansion} along the segment from \(0\) to
\(\phi_\chi\) gives
\(E(\phi_\chi)=\tfrac12\bigl(\int\phi_\chi\,\MA_\omega(0)
+\int\phi_\chi\,\MA_\omega(\phi_\chi)\bigr)
=\tfrac12(\sup\phi_\chi+\tfrac1{12})
=\tfrac12(\tfrac14+\tfrac1{12})=\tfrac16\), where
\(\int\phi_\chi\MA_\omega(\phi_\chi)
=2\int_0^1\tfrac14(1-t)^2\cdot\tfrac12\rd t=\tfrac1{12}\).  The two agree,
which also verifies \eqref{eq:stability-d1-volume-energy} for this
non-finitely-generated \(\chi\): it holds because \(\chi\) is homogeneous,
\(\Psi\) being positively homogeneous of degree one.

\emph{The norms.}  \(\|\chi\|_{\min}=E-\int\phi_\chi\MA_\omega(\phi_\chi)
=\tfrac16-\tfrac1{12}=\tfrac1{12}\), and
\(I=J+\|\chi\|_{\min}=\tfrac16\); as in \Cref{ex:ch50-p1-dnc} the case
\(n=1\) forces \(J=\|\chi\|_{\min}=I/2\), and indeed both are
\(\tfrac1{12}\).  For the centred \(L^1\)-norm, with
\(\bar\lambda=\tfrac16\) and the substitution \(u=\tfrac12+w\),
\[
 \tfrac12\|\chi\|_1=\int_0^1\Bigl(u(1-u)-\tfrac16\Bigr)_+\rd u
 =\int_{-1/\sqrt{12}}^{1/\sqrt{12}}\Bigl(\tfrac1{12}-w^2\Bigr)\rd w
 =\frac{1}{18\sqrt3},
\]
so \(\|\chi\|_1=\tfrac1{9\sqrt3}=\tfrac{\sqrt3}{27}\approx0.0642\).  The
window of \eqref{eq:stability-J-L1-comparison} for \(n=1\) is
\([c_1J,2J]=[\tfrac1{24},\tfrac16]\approx[0.0417,0.1667]\), and
\(0.0642\) lies inside it.

\emph{Entropy and Mabuchi.}  \(A_{\mathbf P^1}(t\ord_p)=t\), so
\(H^{\NA}(\phi_\chi)=\int A\,\rd\MA_\omega(\phi_\chi)
=2\int_0^1t\cdot\tfrac12\rd t=\tfrac12\).  Since
\(K_X\equiv-2\omega\), \Cref{prop:ch50-M-proportional} gives
\[
 M_{X;L}(\chi)=H^{\NA}(\phi_\chi)-2\,\|\chi\|_{\min}
 =\tfrac12-2\cdot\tfrac1{12}=\tfrac13 .
\]
As a check by the Chen--Tian route,
\(R^{\NA}=\gamma\bigl(2E-\int\phi_\chi\MA_\omega\phi_\chi\bigr)
=-2\bigl(\tfrac13-\tfrac1{12}\bigr)=-\tfrac12\) and
\(M=H+R+\overline SE=\tfrac12-\tfrac12+2\cdot\tfrac16=\tfrac13\).
\end{proof}

\begin{remark}[What changes when finite generation fails]
\label{rem:ch50-harmonic-what-changes}
It is worth being precise about which statements of this chapter survive
for \Cref{ex:ch50-harmonic} and which do not.

\emph{Lost.}  There is no test configuration, hence no total space, no
central fibre, no flag ideal in a single degree, and no
Donaldson--Futaki invariant \(\DF\): \Cref{def:ch50-df} does not apply, and
\eqref{eq:stability-M-DF-correction} is vacuous.  The Fubini--Study
function is not PL --- it is strictly convex along the ray, by
\eqref{eq:ch33-harmonic-fs} --- so \Cref{prop:ch46-pl-measure} does not
apply and the Monge--Amp\`ere measure is not atomic; the model entropy
formula \eqref{eq:stability-model-entropy} is likewise unavailable, and
\(H^{\NA}\) must be computed as an integral against a diffuse measure.
Neither \(\lambda_{\max}\) nor any jumping value is attained in a single
degree: the denominators grow with \(m\), by
\Cref{ex:ch33-harmonic-not-fg}\,(2).

\emph{Kept.}  Everything else.  The Duistermaat--Heckman measure exists
(\Cref{thm:dh-limit-filtration}) and satisfies the Brunn--Minkowski
concavity used in \Cref{thm:ch50-J-L1}; the Fubini--Study function is
\(\omega\)-psh and bounded; \(\MA_\omega(\phi_\chi)\) is a probability
measure of finite dual energy; \(E\), \(J\), \(I\),
\(\|\cdot\|_{\min}\), \(\|\cdot\|_1\), \(H^{\NA}\) and \(M_{X;L}(\chi)\)
are all defined and finite, the last by
\eqref{eq:stability-M-all-norms}; and the sharp comparison
\eqref{eq:stability-J-L1-comparison} holds, as verified above.  This is the
precise sense in which \(\mathcal N_{\mathbf R}\) is the right category:
the functionals do not notice that the filtration is not finitely
generated.

\emph{Why it is not destabilizing.}  \(M_{X;L}(\chi)=\tfrac13>0\) and
\(\|\chi\|_{\min}=\tfrac1{12}\), so \(M/\|\cdot\|_{\min}=4\).  This is far
from the infimum \(\sigma_K(\mathbf P^1;\cO(1))=0\) computed in
\Cref{ex:ch50-semistable-not-uniform}, which is approached along a family
of \emph{finitely generated} filtrations.  The example therefore does not
bear on \Cref{rem:stability-gap-general-polarization}; it shows only that
the larger category is nonempty and computable.
\end{remark}

\subsection{K-semistable but not uniformly K-stable}
\label{sec:ch50-ex-semistable}

\begin{example}[The projective line, and the general mechanism]
\label{ex:ch50-semistable-not-uniform}
Let \(X=\mathbf P^1\), \(B=0\), \(L=\cO(1)\).  Then \((X;L)\) is
K-semistable but not uniformly K-stable; more precisely
\[
 \begin{gathered}
 \sigma_K(\mathbf P^1;\cO(1))=\sigma_{\div}(\mathbf P^1;\cO(1))=0,\\
 M^{\NA}(\phi)\ge0\ \text{for every positive model metric},
 \end{gathered}
\tag{50.90}\label{eq:ch50-p1-thresholds}
\]
and the infimum \(0\) is attained, on the product configuration of
\Cref{ex:stability-p1-product}.
\end{example}

\begin{proof}
\emph{K-semistability.}  By \Cref{prop:ch50-threshold-comparisons} the
thresholds are unchanged when \(L\) is replaced by a positive multiple, so
we may test with \(L=\cO(2)=-K_{\mathbf P^1}\).  By
\Cref{ex:stability-p1-delta}, \(\delta(\mathbf P^1)=1\), so
\Cref{thm:stability-valuative-fano}\,(3) --- imported, see
\Cref{imp:ch50-fano-imports}\,\cref{itm:ch50-fano:semistable} --- gives
K-semistability.  A book-internal argument is also available for
\(n=1\): by \Cref{prop:ch50-M-proportional} with \(\gamma=-2\) and
\(L=\cO(1)\),
\[
 M^{\NA}(\phi)=\Ent_X\bigl(\MA_\omega(\phi)\bigr)
   -2\,E^\vee_\omega\bigl(\MA_\omega(\phi)\bigr),
\]
and for a probability measure \(\mu\) on \((\mathbf P^1)^{\An}\) one has
\(E^\vee_\omega(\mu)\le\sum_ia_iE^\vee_\omega(\delta_{v_i})\) by convexity
when \(\mu=\sum a_i\delta_{v_i}\) is atomic, while
\(E^\vee_\omega(\delta_{t\ord_q})=\tfrac12T_\omega(t\ord_q)=\tfrac t2\) by
\eqref{eq:bj-dirac-energy-T} with \(n=1\) and
\(T_{\cO(1)}(\ord_q)=1\); since
\(\Ent_X(\delta_{t\ord_q})=A_X(t\ord_q)=t\), each atom contributes
\(t-2\cdot\tfrac t2=0\), whence \(M^{\NA}(\phi)\ge0\).  Taking a general
\(\mu\in\cM^1\) and using \Cref{thm:ch50-entropy-lsc} together with weak
approximation by atomic measures gives the same conclusion.

\emph{Failure of uniformity.}  By
\Cref{prop:ch50-p1-product}\,\cref{itm:ch50-p1p:mabuchi} the product
configuration has \(M^{\NA}=0\) and
\(\|\chi\|_{\min}=\tfrac12>0\), so the quotient
\(M^{\NA}/\|\cdot\|_{\min}\) vanishes there and
\(\sigma_K=0\).  By \eqref{eq:ch50-threshold-chain},
\(0=\sigma_K\ge\sigma_{\div}\); and \(\sigma_{\div}\ge0\) because
K-semistability gives \(M\ge0\) on divisorial measures, by the computation
above.  Hence both thresholds vanish.  Alternatively, the family
\(\phi_c\) of \Cref{ex:ch50-p1-dnc} has
\(M^{\NA}/J=(2-c)/c\to0\) as \(c\to2\), which exhibits the failure through
a family of \emph{relatively ample} configurations, without appealing to
the semiample boundary case.
\end{proof}

\begin{remark}[The separation is not an accident of dimension one]
\label{rem:ch50-separation-general}
Three mechanisms are known to force K-semistability without uniform
K-stability, and it is worth distinguishing them.

\emph{(a) Positive-dimensional automorphisms.}  If \(\Aut(X,L)\) has
positive dimension, a one-parameter subgroup \(\lambda\) gives a product
configuration whose Fubini--Study function is, by
\Cref{prop:ch49-toric-green}\,\cref{itm:ch49-tg:ma}, the Green function of
the associated valuation \(v_\lambda\); its Monge--Amp\`ere measure is
\(\delta_{v_\lambda}\) and its minimum norm is
\(E^\vee_\omega(\delta_{v_\lambda})>0\).  If in addition the Futaki
character vanishes on \(\lambda\), then \(M^{\NA}=0\) there and uniformity
fails.  \Cref{ex:ch50-semistable-not-uniform} is the case
\(X=\mathbf P^1\), and \Cref{rem:ch50-toric-picture}\,(b) is the toric
case in every dimension: \emph{no} smooth toric log Fano is uniformly
K-stable, while those with \(b_P=0\) --- for instance \(\mathbf P^n\) and
products of projective spaces --- are K-semistable.

\emph{(b) The Ding and Mabuchi functionals may vanish without the metric
being trivial.}  On \(\mathbf P^1\) with \(L=-K\), the anticanonical
member \(c=1\) of \Cref{ex:ch50-p1-dnc} has
\(D^{\NA}=M^{\NA}=\tfrac14>0\); but as \(c\to2\) both tend to \(0\) while
\(J\to1\).  The vanishing in the limit is again explained by (a), the
limit being a product configuration.

\emph{(c) Degenerate polarizations.}  For \(L\) on the boundary of the
ample cone the thresholds can degenerate for reasons unrelated to
automorphisms; \Cref{imp:ch50-openness} shows that this cannot happen in
the interior, \(\sigma_{\div}\) being locally H\"older continuous there.
We do not pursue the boundary case.
\end{remark}

\section{Scope and hypotheses}
\label{sec:ch50-scope}

\subsection{Where each hypothesis is used}
\label{sec:ch50-hypothesis-use}

\begingroup
\begin{longtable}{p{.30\textwidth}p{.62\textwidth}}
\toprule
hypothesis & where it is used\\
\midrule
\endfirsthead
\toprule
hypothesis & where it is used\\
\midrule
\endhead
\(k\) algebraically closed, trivially valued &
Everywhere.  It is what makes \(X^{\An}\) a cone with vertex
\(v_{\triv}\), makes \(\sup\phi=\phi(v_{\triv})\) a maximum principle
(\Cref{prop:ch50-sup-lambdamax}), makes a non-Archimedean norm on a
finite-dimensional space a filtration by subspaces
(\Cref{lem:ch50-norm-is-filtration}), and makes a test configuration a
\(\mathbf G_m\)-degeneration.\\
\(X\) integral projective, \(L\) ample &
Throughout.  Ampleness is needed for the existence of the
Duistermaat--Heckman measure (\Cref{thm:dh-limit-filtration}), for
\Cref{prop:ch49-psef-threshold}, and for every energy identity of
\Cref{sec:ch50-dh}.  Integrality is needed for the Okounkov-body
description \Cref{thm:ch33-dh-concave-transform} and hence for
\Cref{thm:ch50-J-L1}.\\
\(X\) unibranch &
Only through the envelope theorem, hence only where
\Cref{imp:ch50-envelope} is invoked: the Calabi--Yau theorem behind
\eqref{eq:stability-M-all-norms}, the homeomorphism property behind
\Cref{thm:stability-norm-completion}\,(3), the uniqueness statement used to
identify metrics from their measures in
\Cref{prop:ch50-p1-product} and \Cref{ex:ch50-toric-F1}, and the
completeness of \(\cE^1(\omega)\).  \Cref{sec:ch50-norms} through
\Cref{prop:ch50-d1-volume}, all of \Cref{sec:ch50-dh} except
\eqref{eq:stability-minimum-norm}, and
\Cref{thm:stability-chen-tian} do \emph{not} use it.\\
\(X\) normal, \(K_{X,B}\) \(\mathbf Q\)-Cartier &
From \Cref{sec:ch50-entropy} onwards: the definition
\eqref{eq:stability-log-discrepancy} of log discrepancy, hence entropy,
the Ricci energy \eqref{eq:stability-ricci-energy} and everything
downstream.  \Cref{sec:ch50-norms,sec:ch50-dh} are free of it.\\
\(\operatorname{char}k=0\) &
Three places.  (i) \Cref{def:ch50-A-extension} and
\Cref{prop:ch50-A-properties}, through resolution of singularities and
the SNC pairs of \Cref{prop:ch32-log-discrepancy-affine}.  (ii)
\Cref{prop:ch50-entropy-finite-cases}\,\cref{itm:ch50-ef:dense}, through
\Cref{thm:bj-strong-divisorial-density}.  (iii) Every statement of
\Cref{imp:ch50-fano-imports}, whose proofs use the minimal model program
\cite{BCHM10}.\\
\((X,B)\) sublc &
Nonnegativity of \(A_{X,B}\) and hence of entropy
(\Cref{prop:ch50-A-properties}\,\cref{itm:ch50-A:sign},
\Cref{thm:ch50-entropy-lsc}); without it,
\eqref{eq:stability-entropy} needs the truncation conventions of
\Cref{app:normalizations}.\\
\((X,B)\) subklt &
\Cref{thm:stability-divisorial-filtrations} and
\Cref{imp:ch50-strong-lsc}\,\cref{itm:ch50-slsc:extension};
\Cref{prop:ch50-threshold-comparisons}.\\
\((X,B)\) log Fano, \(L=-K_{X,B}\) &
All of \Cref{sec:ch50-ding,sec:ch50-valuative}, and
\Cref{ex:ch50-p1-dnc,ex:ch50-toric-F1}.  See
\Cref{warn:ch50-anticanonical-essential}.\\
\(K_{X,B}\equiv\gamma\omega\) &
\Cref{prop:ch50-M-proportional} and \Cref{cor:ch50-M-dirac}, which are the
only places where \eqref{eq:stability-M-beta} is proved rather than
imported; and hence every example of \Cref{sec:ch50-examples}.\\
\(X\) toric, \(L\) toric &
\Cref{prop:ch50-toric-beta} and \Cref{ex:ch50-toric-F1}, and the
consistency checks in
\Cref{prop:ch50-p1-product},
\Cref{rem:ch50-p1-delta-checks}.\\
\bottomrule
\end{longtable}
\endgroup

\subsection{What is not claimed}
\label{sec:ch50-not-claimed}

\begin{enumerate}
\item\label{itm:ch50-nc:erratum}
      \emph{No asymptotic slope statement.}  Nothing in this chapter
      identifies a limit of a complex energy functional along a ray with a
      non-Archimedean functional.  \Cref{warn:ch50-erratum-firewall} is
      binding, and the quarantine
      \eqref{eq:rf-quarantined-results} is respected without exception.
      Such identifications appear only in \Cref{chap:rays-filtrations},
      each with its own post-erratum source.
\item\label{itm:ch50-nc:gap}
      \emph{No equality of \(\sigma_K\) and \(\sigma_{\div}\) for a general
      polarization.}  Only \eqref{eq:stability-threshold-inequality} is
      asserted.  \Cref{rem:stability-gap-general-polarization} and
      \Cref{rem:ch50-gap-mechanism} record what would be needed.  In the
      log Fano case the gap is closed by \cite{LXZ22}, an import.
\item\label{itm:ch50-nc:dp}
      \emph{The triangle inequality for \(d_p\), \(1<p<\infty\).}  Proved
      here only for \(p=1\) and \(p=\infty\); see
      \Cref{rem:ch50-dp-import}.  Nothing below uses the intermediate
      cases.
\item\label{itm:ch50-nc:completion}
      \emph{The isometry \eqref{eq:stability-norm-ma-isometry} for
      arbitrary norms.}  Imported.  What is proved here is its restriction
      to finite-type norms, where it is
      \Cref{prop:ch46-pl-measure}, and the completeness statements of
      clause (3); see \Cref{rem:ch50-completion-status}.
\item\label{itm:ch50-nc:entropy-bound}
      \emph{No bound of entropy by energy.}  There is no constant \(C\)
      with \(\Ent_{X,B}\le C\,\|\cdot\|_\omega\) on \(\cM^1(\omega)\), by
      \Cref{ex:ch50-infinite-entropy}, and no constant \(C\) with
      \(A_{X,B}\le C\,T_\omega\) on \(X^{\divis}\), by
      \Cref{warn:ch50-A-vs-T}.  Consequently \(M_{X,B;L}\) genuinely takes
      the value \(+\infty\).
\item\label{itm:ch50-nc:coercivity}
      \emph{No coercivity statement in K\"ahler geometry.}  The
      equivalence \eqref{eq:stability-ding-M-uniform} is between two
      non-Archimedean functionals on positive model metrics, imported from
      \cite[Theorem~8.6]{BHJ17}.  It is not, and must not be read as, an
      implication from uniform K-stability to properness of a functional on
      the space of K\"ahler metrics.
\item\label{itm:ch50-nc:ding-filtrations}
      \emph{The Ding functional is not extended to all filtrations.}  See
      \Cref{rem:ch50-ding-scope}.  Only \(M\) is, and only through the
      measure formulation.
\item\label{itm:ch50-nc:openness}
      \emph{No openness of uniform K-stability under deformation of the
      variety.}  \Cref{imp:ch50-openness} concerns variation of the
      polarization within the ample cone of a fixed \((X,B)\), and it is
      imported.
\end{enumerate}

\section{Chapter dictionary}
\label{sec:ch50-dictionary}

\begingroup
\begin{longtable}{p{.23\textwidth}p{.41\textwidth}p{.26\textwidth}}
\toprule
object & content & where\\
\midrule
\endfirsthead
\toprule
object & content & where\\
\midrule
\endhead
graded norm &
\(\chi_m:R_m\to\mathbf R\cup\{+\infty\}\), ultrametric, superadditive,
linearly bounded; equivalently a decreasing multiplicative filtration &
\eqref{eq:stability-norm-axioms},
\eqref{eq:stability-norm-filtration}\\
diagonalization &
a norm on a finite-dimensional space is a filtration by subspaces; adapted
bases are orthogonal &
\Cref{lem:ch50-norm-is-filtration},
\eqref{eq:ch50-orthogonal}\\
joint orthogonality &
two norms admit a common orthogonal basis; the relative spectrum is
intrinsic, with multiplicities \(g(\lambda,\mu)\) &
\Cref{thm:ch50-joint-orthogonal},
\eqref{eq:ch50-bigraded}\\
\(d_\infty\) &
\(=\sup_{s\ne0}|\chi(s)-\chi'(s)|\); on finite-type norms
\(=\sup_{X^{\An}}|\phi-\phi'|\) &
\Cref{prop:ch50-dinfty-sup},
\Cref{prop:ch50-dinfty-fs}\\
\(d_1\) &
\(=\int(\dim F^\lambda_\chi+\dim F^\lambda_{\chi'}-2\dim(F^\lambda_\chi\cap
F^\lambda_{\chi'}))\rd\lambda\); triangle inequality from a dimension
inequality for three subspaces &
\Cref{thm:ch50-d1-integral},
\eqref{eq:ch50-three-subspaces}\\
volume of a norm &
\(\operatorname{vol}(\chi)=\int\lambda\,\rd\DHm_\chi\);
\(d_1=\operatorname{vol}(\chi)+\operatorname{vol}(\chi')
-2\operatorname{vol}(\chi\wedge\chi')\) &
\eqref{eq:stability-volume-norm},
\Cref{prop:ch50-d1-volume}\\
Galois connection &
\(\FS(\chi)\le\phi\iff\chi\le\chi^\phi\);
\(\chi\mapsto\chi^{\FS(\chi)}\) is a closure operator &
\Cref{thm:ch50-galois},
\eqref{eq:ch50-galois}\\
homogenization &
\(\chi^{\hom}_m(s)=\lim_rr^{-1}\chi_{rm}(s^r)\);
\(\FS(\chi^{\hom})=\FS(\chi)\), \(d_1(\chi,\chi^{\hom})=0\) &
\Cref{prop:ch50-homogenization},
\Cref{rem:stability-three-equivalences}\\
Hausdorff quotient &
\((\mathcal T_\Lambda,d_p)\to\mathcal H_\Lambda(L)\); completion is
\(\cM^1(L)\), and \(\cE^1(L)\) under the envelope property &
\Cref{thm:stability-fs-quotient},
\Cref{thm:stability-norm-completion}\\
maximum principle &
\(\sup\phi_\chi=\phi_\chi(v_{\triv})=\lambda_{\max}(\chi)
=\max\operatorname{supp}\DHm_\chi\) &
\Cref{prop:ch50-sup-lambdamax}\\
energy as a moment &
\(E(\FS\chi)=\int\lambda\,\rd\DHm_\chi\);
\(J=\lambda_{\max}-\int\lambda\,\rd\DHm_\chi\) &
\Cref{thm:ch50-energy-dh},
\eqref{eq:stability-dh-energy-J}\\
sharp \(J\)--\(L^1\) comparison &
\(c_nJ\le\|\chi\|_1\le2J\), \(c_n=2n^n/(n+1)^{n+1}\), extremal on cones &
\Cref{thm:ch50-J-L1},
\Cref{rem:ch50-extremal-cone}\\
minimum norm &
\(\|\chi\|_{\min}=E^\vee(\MA\chi)=E-\int\phi\MA\phi=I-J\);
\(\frac1nJ\le\|\chi\|_{\min}\le nJ\) &
\eqref{eq:stability-minimum-norm},
\Cref{prop:ch50-minimum-norm}\\
log discrepancy &
lsc \([0,+\infty]\)-valued, homogeneous, affine on quasimonomial cones;
\(+\infty\) at proper support &
\Cref{def:ch50-A-extension},
\Cref{prop:ch50-A-properties}\\
entropy &
\(\Ent(\mu)=\int A\,\rd\mu\); affine, weakly lsc, \(+\infty\) on part of
\(\cM^1\) &
\Cref{thm:ch50-entropy-lsc},
\Cref{ex:ch50-infinite-entropy}\\
model entropy &
\(H^{\NA}=V^{-1}\sum_EA(v_E)b_E(E\cdot\overline\cL^n)\) &
\eqref{eq:stability-model-entropy},
\Cref{prop:ch50-model-entropy}\\
Ricci energy &
\(R^{\NA}=V^{-1}(K_{X,B},0)(\omega,\phi)^n\); \(=\gamma((n+1)E-\int\phi\MA
\phi)\) when \(K_{X,B}\equiv\gamma\omega\) &
\Cref{lem:ch50-ricci-energy}\\
Chen--Tian &
\(M^{\NA}=H^{\NA}+R^{\NA}+\overline SE\);
\(K^{\log}=p^*K_{X,B}+\sum b_EA(v_E)E\) &
\Cref{thm:stability-chen-tian},
\eqref{eq:ch50-log-canonical-atoms}\\
\(M\) versus \(\DF\) &
\(M^{\NA}=\DF+V^{-1}((\cX_{0,\mathrm{red}}-\cX_0)\cdot\overline\cL^n)
\le\DF\), equality iff reduced fibre &
\eqref{eq:stability-M-DF-correction}\\
proportional case &
\(M^{\NA}=H^{\NA}+\gamma\|\MA\phi\|_\omega\); for a Dirac,
\(=A(v)+\gamma S_L(v)\) &
\Cref{prop:ch50-M-proportional},
\Cref{cor:ch50-M-dirac}\\
Ding &
\(L^{\NA}=\inf_v(A+\phi)\), concave, \(1\)-Lipschitz;
\(D^{\NA}=L^{\NA}-E\le M^{\NA}\), equality on anticanonical metrics &
\Cref{prop:ch50-L-elementary},
\Cref{thm:stability-ding-mabuchi},
\Cref{lem:ch50-anticanonical-constant}\\
three thresholds &
\(\sigma_K\ge\sigma_{\mathcal N}=\sigma_{\div}\); K-semistable
\(\iff\sigma_K\ge0\), uniformly K-stable \(\iff\sigma_K>0\) &
\eqref{eq:ch50-three-thresholds},
\Cref{prop:ch50-threshold-comparisons},
\Cref{prop:ch50-model-form}\\
valuative invariants &
\(T=\lambda_{\max}(\chi_v)\), \(S=\operatorname{vol}(\chi_v)
=E^\vee(\delta_v)\), \(j=J(\FS\chi_v)\);
\(\frac1nS\le j\le nS\) &
\Cref{prop:ch50-valuative-dh},
\eqref{eq:ch50-j-vs-S}\\
\(\delta\)-invariant &
\(\delta=\inf A/S\); the valuative conditions are equivalent to each other
with explicit constants &
\eqref{eq:stability-delta},
\Cref{prop:ch50-valuative-internal}\\
toric dictionary &
\(A(v_u)=f_P(u)\), \(S(v_u)=f_P(u)+\langle u,b_P\rangle\),
\(\beta(v_u)=-\langle u,b_P\rangle\); \(\delta_{\mathrm{tor}}
=(1+\max_Q\langle\cdot,b_P\rangle)^{-1}\) &
\Cref{prop:ch50-toric-beta},
\eqref{eq:ch50-toric-beta},
\eqref{eq:ch50-toric-delta}\\
\bottomrule
\end{longtable}
\endgroup

\section{The dated scope table}
\label{sec:ch50-dated}

The chapter draws on four layers of literature.  The table records, for
each imported statement, the source, the year of first public availability
and of publication, and the exact use made of it here.  Statements proved
in this book are listed separately below.

\subsection{Layer one: the foundational sources}
\label{sec:ch50-dated-foundational}

\begin{center}
\small
\begin{tabular}{p{.27\textwidth}p{.12\textwidth}p{.13\textwidth}p{.34\textwidth}}
\toprule
statement & source & dates & used here for\\
\midrule
Duistermaat--Heckman measure of a filtration; concave transform &
\cite{BC11} &
2011 pre., 2011 pub. &
\Cref{thm:ch33-dh-concave-transform}, hence
\Cref{thm:ch50-J-L1} and \Cref{prop:ch50-valuative-dh}\\
\(\DF\), \(M^{\NA}\), Chen--Tian, \(D^{\NA}\le M^{\NA}\), the
\(J\)--\(\|\cdot\|_1\) comparison (\S\S6--8) &
\cite{BHJ17} &
2015 pre., 2017 pub. &
\Cref{sec:ch50-mabuchi,sec:ch50-ding}; proved here except
\eqref{eq:stability-ding-M-uniform}\\
special test configurations (Prop.~7.32, Thm.~8.6);
\(M\ge\delta J\Rightarrow D\ge\delta J\) &
\cite{BHJ17} &
2015 pre., 2017 pub. &
the coercive half of \eqref{eq:stability-ding-M-uniform}; uses
\cite{BCHM10}\\
trivially valued pluripotential theory; \(E\), \(\cE^1\), \(\cM^1\),
\(T_\omega\), toric formulae &
\cite{BJ22,BJ22b} &
2018 pre., 2022 and 2024 pub. &
the whole of
\Crefrange{chap:bj-energy}{chap:bj-linear-growth}, quoted here as proved
material\\
log discrepancy on valuation spaces; monotonicity under refinement &
\cite{JM12} &
2010 pre., 2012 pub. &
\Cref{prop:ch50-A-properties}\,\cref{itm:ch50-A:consistent}\\
\bottomrule
\end{tabular}
\end{center}

\subsection{Layer two: the later sources, each dated}
\label{sec:ch50-dated-later}

\begin{center}
\small
\begin{tabular}{p{.27\textwidth}p{.14\textwidth}p{.12\textwidth}p{.33\textwidth}}
\toprule
statement & source & dates & used here for\\
\midrule
metric geometry of norms; \(d_p\) for \(1<p<\infty\); Theorems A, B, C;
\(S_L=E^\vee(\delta_\cdot)\) &
\cite{BJKS1} &
2021 pre., 2025 pub. &
\Cref{imp:ch50-imported}\,\cref{itm:ch50-import:dp}--\cref{itm:ch50-import:S};
\Cref{thm:stability-fs-quotient},
\Cref{thm:stability-norm-completion}\\
\(C^1\)-regularity of \(\omega\mapsto\|\mu\|_\omega\); strong lsc extension
of \(\beta\); divisorial stability; H\"older continuity &
\cite{BJKS2} &
June 2022 pre., 2023 pub. &
\Cref{imp:ch50-strong-lsc},
\Cref{thm:stability-divisorial-filtrations},
\Cref{imp:ch50-openness}\\
Conjecture~5.15 on \(\sigma_K=\sigma_{\div}\) for general \(L\) &
\cite{BJKS2} &
June 2022 pre., 2023 pub. &
recorded as \emph{open} in
\Cref{rem:stability-gap-general-polarization}; not used\\
K-semistability is equivariant volume minimization &
\cite{Li17} &
2015 pre., 2017 pub. &
\Cref{thm:stability-valuative-fano}\,(1)\\
valuative criterion for uniform K-stability of \(\mathbf Q\)-Fano
varieties &
\cite{Fuj19} &
2016 pre., 2019 pub. &
\Cref{thm:stability-valuative-fano}\,(2)\\
thresholds, valuations and K-stability; equivariance of \(\delta\) &
\cite{BJ20} &
2017 pre., 2020 pub. &
\Cref{thm:stability-valuative-fano}\,(3);
\Cref{rem:ch50-toric-picture}\,(c)\\
finite generation for valuations computing \(\delta\) &
\cite{LXZ22} &
2021 pre., 2022 pub. &
cited in \Cref{rem:ch50-gap-mechanism} for the log Fano closure of the
gap; not used elsewhere\\
K-moduli, optimal degenerations; \(G\)-uniform stability; continuity of
\(\delta\) &
\cite{BHLLX21}, \cite{Li22}, and \cite{Zhang21} &
2021--2022 pub. &
\emph{not used}; listed in
\Cref{imp:ch50-fano-imports}\,\cref{itm:ch50-fano:openness} to mark the
boundary\\
\bottomrule
\end{tabular}
\end{center}

\begin{warning}{Foundational versus later}
\label{warn:ch50-foundational-boundary}
The foundational layer of this part is \cite{BJ22} together with its
addendum \cite{BJ22b} and, for the intersection-theoretic functionals,
\cite{BHJ17}.  Everything in
\Crefrange{chap:bj-energy}{chap:bj-linear-growth} belongs to that layer.
The metric geometry of \emph{all} norms
(\Cref{thm:stability-fs-quotient,thm:stability-norm-completion}), the
divisorial Monge--Amp\`ere bijection
\eqref{eq:stability-divnorm-divmeasure}, the identity
\(S_L=E^\vee(\delta_\cdot)\), the differentiability of
\(\omega\mapsto\|\mu\|_\omega\) and divisorial stability
\emph{do not} belong to it: they are \cite{BJKS1} and \cite{BJKS2}, dated
2021 and 2022 as preprints.  The Fano valuative criteria form a third,
independent line, \cite{Li17,Fuj19,BJ20,LXZ22}, which is not part of the
trivially valued pluripotential theory at all and which uses the minimal
model program.  We have not folded any of the later material backwards
into \cite{BJ22}, and no proof presented here as foundational uses a later
result.
\end{warning}

\section{Source and provenance}
\label{sec:ch50-source}

\subsection{Source note}
\label{sec:ch50-source-note}

The model identities begin with \eqref{eq:stability-dh-energy-J} and continue
through \eqref{eq:stability-M-DF-correction}; they are proved in
\cite[\S\S6--8]{BHJ17}.  The norm/filtration
Hausdorff quotients, completions, minimum norm, and divisorial
Monge--Amp\`ere correspondence are Theorems~A--C and
Proposition~7.20 of \cite{BJKS1}.  The differentiability in the ample class,
the extended entropy, and \Cref{thm:stability-divisorial-filtrations} are
Theorem~A, \S\S3--4, and Theorem~5.8 of \cite{BJKS2}; the finite-energy
differentiability statement used above is \cite[Theorem~2.15]{BJKS2}.
The latter paper also proves
local H\"older continuity of \(\sigma_{\div}\) on the ample cone
\cite[Theorem~4.12]{BJKS2}.

The chronology matters.  BJ22 supplies the pluripotential foundation but
does not contain the later completion theorem for arbitrary norms or the
later equivalence with divisorial stability.  Conversely, the general
polarization equality discussed in
\Cref{rem:stability-gap-general-polarization} was conjectural in that
source.  No later assertion has been folded backward into BJ22.

Statement by statement, the remaining attributions are as follows.  The
Duistermaat--Heckman material of \Cref{sec:ch50-dh} rests on
\Cref{chap:bj-test-configurations}, hence on \cite{BC11} and
\cite[\S3]{BHJ17}.  The energy identities rest on
\Cref{chap:bj-energy,chap:bj-finite-energy}, hence on
\cite[\S\S7--10]{BJ22}.  The Dirac comparison
\eqref{eq:bj-dirac-energy-T} used in
\Cref{prop:ch50-valuative-dh} and \Cref{ex:ch50-infinite-entropy} is
\cite[Proposition~11.1]{BJ22}, proved in
\Cref{thm:bj-dirac-finite-energy}.  The toric formulae are
\cite[Appendix~B]{BJ22}, proved in
\Cref{thm:bj-toric-ma-energy,prop:ch49-toric-green}.  Log discrepancy on
valuation spaces is \cite{JM12}.  The Fano line is
\cite{Li17,Fuj19,BJ20,LXZ22}, as tabulated in
\Cref{sec:ch50-dated}.

\begin{important}{Erratum confirmation}
\label{imp:ch50-erratum-confirmation}
No statement of this chapter uses \cite[Theorem~5.6]{BHJ19},
\cite[Theorem~5.4]{BHJ19}, \cite[Theorem~C]{BHJ19} or
\cite[Corollaries~D--E]{BHJ19}, directly or transitively, nor any
reformulation of one of them.  The only paper of that group of authors used
here is \cite{BHJ17}, a different and earlier paper, unaffected by the
erratum \cite{BHJ19err}; see
\Cref{warn:ch50-erratum-firewall}\,\cref{itm:ch50-fw:bhj17ok}.  The
citations \cite[Theorem~C]{BJKS1} and \cite[Theorem~A]{BJKS1} refer to
\cite{BJKS1}, not to \cite{BHJ19}.  \Cref{thm:rf-erratum-audit} is the
book's dependency ledger and is respected here without exception.
\end{important}

\subsection{What is supplied here and is not in the sources in this form}
\label{sec:ch50-additions}

\Cref{lem:ch50-norm-is-filtration} and
\Cref{thm:ch50-joint-orthogonal}: the elementary theory of
non-Archimedean norms over a trivially valued field, and in particular the
complete proof, by a bigraded dimension count, that two such norms admit a
joint orthogonal basis.  The sources take diagonalization for granted; we
give the argument because it is exactly what makes the relative spectrum,
and hence every \(d_p\), well defined, and because
\Cref{rem:ch50-three-flags} explains why it cannot be extended to three
norms and hence why the triangle inequality is a theorem.

\Cref{prop:ch50-dinfty-sup} and \Cref{thm:ch50-d1-integral}: the closed
forms \eqref{eq:ch50-dinfty-sup} and \eqref{eq:ch50-d1-integral} for
\(d_\infty\) and \(d_1\), and the reduction of the triangle inequality for
\(d_1\) to the elementary inequality \eqref{eq:ch50-three-subspaces} for
three subspaces of a vector space.  The lattice identity of
\Cref{lem:ch50-lattice}, and \Cref{prop:ch50-d1-volume}, are stated in the
sources; the proofs given here are self-contained.

\Cref{thm:ch50-galois}, \Cref{prop:ch50-supnorm} and
\Cref{prop:ch50-homogenization}: the adjunction between the Fubini--Study
operator and the sup-norm operator, the resulting closure operator, and the
proof that homogenization does not change the Fubini--Study function.  The
one-line proof of \eqref{eq:ch50-fs-homog} --- multiply by \(r\), use
\(v(s^r)=rv(s)\), take the infimum --- is the reason the whole theory is
insensitive to homogenization, and \Cref{ex:ch50-nonhomogeneous} shows that
the insensitivity is not a pointwise identity.

\Cref{prop:ch50-sup-lambdamax}: the identification of
\(\sup\FS(\chi)\) with \(\lambda_{\max}(\chi)\), proved directly from
\(v(s)\ge0\) without using pshness, so that it is available for arbitrary
linearly bounded filtrations.

\Cref{thm:ch50-J-L1}: a complete and self-contained proof of the sharp
two-sided comparison \eqref{eq:stability-J-L1-comparison}, isolated as a
statement about probability measures whose tail function has a concave
\(n\)-th root.  The lower bound is obtained in two steps --- the mass above
the mean is at least \((n/(n+1))^n\), by
\eqref{eq:ch50-mass-above-mean}, and the tail integral above the mean is at
least \(F(\bar\lambda)A/(n+1)\) --- both of which are elementary
consequences of concavity; the extremal measure is exhibited.  The
statement is \cite[Theorem~7.9]{BHJ17}; this proof, and the identification
of the extremal cone in \Cref{rem:ch50-extremal-cone} with the extremal
configuration of \Cref{cor:ch49-minkowski}, are supplied here.

\Cref{prop:ch50-minimum-norm}: the identity
\eqref{eq:ch50-min-norm-identities} in the form
\(\|\chi\|_{\min}=E^\vee(\MA\phi)=J_\phi(0)=I-J\), derived from the Fenchel
equality of \Cref{prop:bj-fenchel-equalities}, together with the explicit
comparison \eqref{eq:ch50-min-vs-J} with \(J\).

\Cref{def:ch50-A-extension}, \Cref{prop:ch50-A-properties} and
\Cref{thm:ch50-entropy-lsc}: the construction of \(A_{X,B}\) on
\(X^{\An}\) as a supremum of continuous functions along model retractions,
and the resulting one-line proof of weak lower semicontinuity of the
entropy via \eqref{eq:ch50-entropy-sup-continuous}.  The
failure of upper semicontinuity for entropy,
\Cref{ex:ch50-entropy-jumps}, and the
warning \Cref{warn:ch50-A-vs-T} that \(A/T_\omega\) is unbounded are
supplied here.

\Cref{ex:ch50-infinite-entropy}: an explicit measure on
\((\mathbf P^2)^{\An}\) with dual energy at most \(2/3\) and infinite
entropy, built from smooth plane curves of increasing degree.  That entropy
is not finite on \(\cM^1\) is stated in the sources; a concrete witness,
with all constants, is supplied here, and it is what makes
\Cref{rem:ch50-M-all-norms-scope} precise.

\Cref{lem:ch50-ricci-energy} and \Cref{prop:ch50-M-proportional}: the
identification of the Ricci energy with a mixed energy pairing, and the
closed formula \(M^{\NA}=H^{\NA}+\gamma\|\MA\phi\|_\omega\) in the
numerically proportional case.  The latter proves
\eqref{eq:stability-M-beta} in that case \emph{without} the
differentiability import
\Cref{imp:ch50-strong-lsc}\,\cref{itm:ch50-slsc:diff}, and it is what makes
every example of \Cref{sec:ch50-examples} computable by hand.
\Cref{cor:ch50-M-dirac} then identifies \(M^{\NA}\) of a Green metric with
the Fujita--Li \(\beta\)-invariant, book-internally.

The complete proof of \Cref{thm:stability-chen-tian}, including the
bookkeeping \eqref{eq:ch50-log-canonical-atoms} of the coefficients
\(b_EA_{X,B}(v_E)\) via \Cref{lem:ch40-gauss-log-discrepancy}, and the
warning \Cref{rem:ch50-chen-tian-provenance} about the two places where a
multiplicity or a sign is easily lost.

\Cref{lem:ch50-anticanonical-constant} and the complete proof of
\eqref{eq:stability-ding-less-M} together with its equality case.  The
identity \eqref{eq:ch50-anticanonical-defect}, which expresses the defect
\(A_{X,B}(v)+\phi(v)-c\) of an anticanonical metric as a log discrepancy of
the pair \((\cX,B_\cX+\cX_{0,\mathrm{red}})\) divided by \(b_{E'}\), is the
precise mechanism by which sublc-ness enters, and it is written out here.

\Cref{prop:ch50-model-form},
\Cref{prop:ch50-threshold-comparisons} and
\Cref{prop:ch50-valuative-internal}: the reduction of
\eqref{eq:stability-uniform-K} to the model form
\eqref{eq:stability-M-coercive-model} with explicit constants; the chain
\(\sigma_K\ge\sigma_{\mathcal N}=\sigma_{\div}\); and the proof that all
the valuative conditions of \Cref{thm:stability-valuative-fano} are
equivalent to \emph{each other}, with explicit constants, so that the
imported content is exactly the link with K-stability.

\Cref{prop:ch50-toric-beta}: the closed formulae
\eqref{eq:ch50-toric-beta} and \eqref{eq:ch50-toric-delta} expressing
\(A\), \(S\), \(\beta\) and the toric stability threshold of a smooth toric
log Fano in terms of the reflexive polytope and its barycenter.  The
statement that \(A_X(v_u)=f_P(u)\) --- the log discrepancy \emph{is} the
support function of \(-P\) --- is the reason the whole computation
collapses to the barycenter, and it is proved here from
\Cref{prop:ch32-log-discrepancy-affine}.

All four worked examples:
\Cref{ex:ch50-p1-dnc}, whose table \eqref{eq:ch50-p1-dnc-table} evaluates
every functional of the chapter on a one-parameter family and locates the
anticanonical member;
\Cref{ex:ch50-toric-F1}, where \(\delta(\mathbb F_1)=6/7\) is computed
twice, from the polytope and from a volume integral;
\Cref{ex:ch50-harmonic}, where every invariant of a
non-finitely-generated filtration is computed and
\Cref{rem:ch50-harmonic-what-changes} separates what is lost from what is
kept; and \Cref{ex:ch50-semistable-not-uniform}.  Also
\Cref{ex:ch50-nonhomogeneous} and \Cref{ex:ch50-entropy-jumps}.

\subsection{Historical note}
\label{sec:ch50-history}

The subject begins with a variational principle and a numerical
obstruction.  Futaki's character \cite{Fut83} obstructs the existence of a
K\"ahler--Einstein metric on a Fano manifold with nontrivial holomorphic
vector fields; Mabuchi's energy \cite{Mab87} makes it the derivative of a
functional on the space of K\"ahler metrics.  Tian \cite{Tia97} introduced
K-stability as a limiting form of the Futaki obstruction over
degenerations, and Donaldson \cite{Don02} recast it algebraically in terms
of test configurations and the invariant that bears his and Futaki's names.
Ross and Thomas \cite{RT07} made the Hilbert--Mumford analogy systematic;
Odaka \cite{Od13,Od13b} discovered that K-semistability forces the pair to
have mild singularities, which is the first appearance of log discrepancy
in the subject and the ancestor of the entropy term of
\eqref{eq:stability-chen-tian}.

The functional-analytic reorganization is due to Boucksom, Hisamoto and
Jonsson \cite{BHJ17}: the Donaldson--Futaki invariant becomes a sum of an
entropy, a Ricci energy and a scalar-curvature term; the norms
\(\|\cdot\|_p\) of \eqref{eq:stability-spectral-lp} become moments of a
Duistermaat--Heckman measure; and uniform K-stability, in the sense of
Dervan \cite{Der16}, becomes a coercivity statement.  That paper also
isolates the minimum norm and proves the sharp comparison
\eqref{eq:stability-J-L1-comparison} reproved here.  Its companion
\cite{BHJ19} treated the K\"ahler-geometric asymptotics; part of that
companion was retracted in \cite{BHJ19err}, and
\Cref{warn:ch50-erratum-firewall} records the resulting boundary, which
this chapter does not cross.

In parallel, Fujita \cite{Fuj19} and Li \cite{Li17} discovered that for a
\(\mathbf Q\)-Fano variety the whole of K-stability can be read off
individual divisorial valuations, through the invariant
\(\beta_L(v)=A(v)-S_L(v)\); Blum and Jonsson \cite{BJ20} organized this
into the threshold \(\delta\) of \eqref{eq:stability-delta} and connected
it to log canonical thresholds of complete linear systems.  The reduction
to special test configurations, which is what converts an arbitrary
degeneration into a valuation, rests on the minimal model program
\cite{BCHM10}; and Liu, Xu and Zhuang \cite{LXZ22} completed the picture by
showing that a valuation computing \(\delta\le1\) has a finitely generated
associated graded ring.  Xu's survey \cite{Xu21} is the standard account of
this line.

The thread pursued in this chapter is the third one.  Boucksom and Jonsson
\cite{BJ22} built the pluripotential theory over a trivially valued field;
its sequels \cite{BJKS1,BJKS2} used it to metrize the space of \emph{all}
graded norms and to identify its completion with the space of finite-energy
measures, so that stability becomes a positivity statement for a
functional on a complete metric space rather than a family of inequalities
indexed by degenerations.  The move is the exact analogue of passing from
smooth K\"ahler metrics to \(\cE^1\) in the complex theory, and it has the
same payoff: existence of minimizers, continuity of thresholds, and a
notion --- divisorial stability --- that is stable under the limits one
actually needs.  The synthetic account of finite-energy measures is
\cite{BJ25FiniteEnergy}, and the weighted extension is
\cite{BJ25WeightedYTD}.

Two threads leave this chapter.  The first goes to
\Cref{chap:rays-filtrations}, where the filtrations of
\Cref{sec:ch50-norms} are matched with geodesic rays in the space of
positively curved metrics on a complex manifold, and where the asymptotic
slope statements deliberately avoided here are proved under their own,
post-erratum, hypotheses.  The second goes to
\Cref{chap:ytd-bridge}, where the functionals of
\Cref{sec:ch50-mabuchi,sec:ch50-ding} are matched with their complex
counterparts and the Yau--Tian--Donaldson correspondence is assembled;
\Cref{chap:master-dictionary} records the resulting table, and
\Cref{chap:mpwn-k-stability} treats the transcendental case, where no
section ring is available and the filtration picture must be replaced by
test curves.
